% Traductions : l'ordre, l'interdiction, « chez », « devenir » ; debates : educación de los niños y NTIC — Espagnol Tle A — Cours signé OnBuch+
% Programme officiel MINESEC. Compiler avec : tectonic E06-orden-prohibicion-chez-devenir.tex
\newcommand{\DOCMATIERE}{Espagnol}
\newcommand{\DOCNIVEAU}{Tle A}
\newcommand{\DOCMODULE}{Módulo 1 : Vie familiale et intégration sociale}
\newcommand{\DOCLECON}{E06}
\newcommand{\DOCTITRE}{Traductions : l'ordre, l'interdiction, « chez », « devenir » ; debates : educación de los niños y NTIC}
% OnBuch+ — préambule « cours signés » ENRICHI (compilé avec tectonic / XeLaTeX)
% Système visuel « Pop » d'OnBuch+ : fond crème (jamais blanc), bordures encre
% épaisses, ombres « dures » décalées, titres Archivo Black en capitales.
% Version MATHÉMATIQUES : graphes (pgfplots/tikz), boîtes pédagogiques
% (définition, propriété, méthode, attention, le savais-tu, exercice résolu,
% exercice, TP, bilan), page de garde et signature OnBuch+ ancrée en bas.
%
% Macros attendues dans main.tex AVANT \input{preamble} :
%   \DOCMATIERE  \DOCNIVEAU  \DOCMODULE  \DOCLECON (numéro)  \DOCTITRE
\documentclass[11pt]{article}

\usepackage{fontspec}
\usepackage[spanish,french]{babel} % français = langue principale ; l'espagnol via \esp{..} / espagnol
\usepackage[a4paper,margin=2cm,top=3.4cm,bottom=2.8cm,headheight=1.4cm,footskip=1.3cm]{geometry}
\usepackage{amsmath,amssymb,amsfonts}
\usepackage{mathtools} % \xmapsto, \coloneqq... souvent utilisés par les modèles
\usepackage{cancel} % simplifications barrées (bilans de forces)
% \abs{z} (module, valeur absolue) et \norm{u} (norme), utilisés par les modèles
\providecommand{\abs}{}\let\abs\relax\DeclarePairedDelimiter{\abs}{\lvert}{\rvert}
\providecommand{\norm}{}\let\norm\relax\DeclarePairedDelimiter{\norm}{\lVert}{\rVert}
\usepackage[table]{xcolor} % [table] : fournit aussi \rowcolors (alternance de lignes)
\usepackage{pagecolor}
\usepackage{tikz}
\usetikzlibrary{arrows.meta,positioning,calc,shapes.geometric,shapes.misc,shapes.arrows,decorations.pathmorphing,decorations.pathreplacing,decorations.markings,patterns,shadows,mindmap,trees,fit,backgrounds,matrix,intersections,angles,quotes}
\usepackage{pgfplots}
\usepackage[version=4]{mhchem} % équations nucléaires \ce{^{235}_{92}U}
\usepackage[european,siunitx]{circuitikz} % schémas électriques
\pgfplotsset{compat=1.18}
\usepgfplotslibrary{fillbetween,groupplots} % aires entre courbes (calcul intégral)
\usepackage{enumitem}
\usepackage{fancyhdr}
% [most] charge toutes les bibliothèques tcolorbox (listings, minted, raster,
% documentation...) et gonfle inutilement la mémoire TeX sur un document long
% (~300 boîtes) : on ne charge que ce qui est réellement utilisé.
\usepackage[skins,breakable]{tcolorbox}
\usepackage{graphicx}
\usepackage{colortbl}
\usepackage{booktabs}
\usepackage{tabularx}
\usepackage{diagbox} % case d'en-tête diagonale des tableaux à double entrée
% Colonnes extensibles centrée / gauche / droite (C, L, R), souvent utilisées par les modèles
\newcolumntype{C}{>{\centering\arraybackslash}X}
\newcolumntype{L}{>{\raggedright\arraybackslash}X}
\newcolumntype{R}{>{\raggedleft\arraybackslash}X}
\usepackage{array}
\usepackage{multicol}
\usepackage{siunitx}
% parse-numbers=false : accepte aussi \SI{2,5 \times 10^{-3}}{...}
\sisetup{locale=FR,output-decimal-marker={,},group-minimum-digits=4,parse-numbers=false}
\DeclareSIUnit\ohm{\ensuremath{\symup{\Omega}}} % U+2126 absent de la police
\DeclareSIUnit\division{div} % réglages d'oscilloscope (ms/div, V/div)
\DeclareSIUnit\tr{tr} % tours (vitesse de rotation, ex. \SI{1500}{\tr\per\minute})
\DeclareSIUnit\molar{M} % concentration molaire (ex. \SI{3}{\molar}, \SI{1}{\micro\molar})
\DeclareSIUnit\an{an} % année (le modèle confond parfois avec \year, un compteur TeX) — ex. \SI{5}{\kilo\meter\per\an}
\DeclareSIUnit\FCFA{FCFA} % franc CFA (ex. \SI{500}{\FCFA\per\kilo\gram})
\DeclareSIUnit\degreeBrix{\textdegree Brix} % teneur en sucre d'un sirop/jus (réfractométrie)
\DeclareSIUnit\month{mois} % le modèle utilise parfois \month, un compteur TeX, comme unité de durée
\usepackage{qrcode}
% \degree : commande fréquemment utilisée par les modèles pour un angle ou une
% température en dehors de \SI{}{} (non fournie par les paquets chargés).
\providecommand{\degree}{^{\circ}}
% \qed (fin de démonstration), utilisé par les modèles sans amsthm
\providecommand{\qed}{\hfill$\square$}
% \barre : double barre des valeurs interdites dans un tableau de variations (tkz-tab)
\providecommand{\barre}{\ensuremath{\|}}
% environnement proof (amsthm non chargé) utilisé par les modèles
\ifdefined\proof\else\newenvironment{proof}[1][Démonstration]{\par\noindent\textit{#1.}\ }{\qed\par}\fi
\usepackage{lastpage}
\usepackage{hyperref}
\hypersetup{hidelinks}

% ============================================================
% Palette « Pop » OnBuch+ — mêmes hex que la classe P (plus_ui.dart)
% ============================================================
\definecolor{poporange}{HTML}{F59321}
\definecolor{popdark}{HTML}{E07A0C}
\definecolor{popcream}{HTML}{FFF6E8}
\definecolor{popink}{HTML}{1C1714}
\definecolor{poppurple}{HTML}{7A5AE0}
\definecolor{popgreen}{HTML}{2E9E6A}
\definecolor{poppink}{HTML}{E0457B}
\definecolor{popblue}{HTML}{2F7DE1}
\definecolor{popgold}{HTML}{C98A12}
\definecolor{popmuted}{HTML}{8A7F76}
\definecolor{popline}{HTML}{EADFD2}
\definecolor{popgray}{HTML}{8A7F76}
\colorlet{popgrey}{popgray}
% Alias tolérants (le modèle nomme parfois les couleurs différemment)
\colorlet{popwhite}{white}
\colorlet{popblack}{popink}
\colorlet{popred}{poppink}
\colorlet{popyellow}{popgold}
% Teintes claires (fonds de tableaux/encadrés) — le modèle nomme parfois
% les couleurs avec un suffixe L (ex. popblueL) sans les avoir définies.
\colorlet{poporangeL}{poporange!20}
\colorlet{popdarkL}{popdark!20}
\colorlet{poppurpleL}{poppurple!20}
\colorlet{popgreenL}{popgreen!20}
\colorlet{poppinkL}{poppink!20}
\colorlet{popblueL}{popblue!20}
\colorlet{popgoldL}{popgold!20}
\colorlet{popredL}{popred!20}
\colorlet{popgrayL}{popgray!20}
% Teintes claires pour fonds de tableaux / zones
\colorlet{popblueL}{popblue!10!white}
\colorlet{poporangeL}{poporange!14!white}
\colorlet{popgreenL}{popgreen!12!white}
\colorlet{poppurpleL}{poppurple!12!white}
\colorlet{poppinkL}{poppink!12!white}
\colorlet{popgoldL}{popgold!14!white}

\pagecolor{popcream}

% ============================================================
% Typographie
% ============================================================
\defaultfontfeatures{Ligatures=TeX}
\setmainfont{PlusJakartaSans-Medium.ttf}[Path=../../../fonts/,BoldFont=PlusJakartaSans-Bold.ttf]
% Maths en sans-serif (Fira Math) avec chiffres et lettres droites de Plus
% Jakarta, pour que les formules se fondent dans le texte.
\usepackage[math-style=ISO,bold-style=ISO]{unicode-math}
\setmathfont{FiraMath-Regular.otf}[Path=../../../fonts/]
\setmathfont{PlusJakartaSans-Medium.ttf}[Path=../../../fonts/,range={up/num,up/{Latin,latin},\mathup}]
\usepackage{icomma} % virgule décimale sans espace en mode math
% Caractères mathématiques Unicode absents des polices de texte (Plus Jakarta,
% Archivo Black) : rendus par leur équivalent mathématique, en texte comme en
% formule — sinon ils s'affichent en carré vide (ex. « Arithmétique dans ℤ »).
\usepackage{newunicodechar}
\newunicodechar{ℝ}{\ensuremath{\mathbb{R}}}
\newunicodechar{ℤ}{\ensuremath{\mathbb{Z}}}
\newunicodechar{ℂ}{\ensuremath{\mathbb{C}}}
\newunicodechar{ℕ}{\ensuremath{\mathbb{N}}}
\newunicodechar{ℚ}{\ensuremath{\mathbb{Q}}}
\newunicodechar{𝔻}{\ensuremath{\mathbb{D}}}
\newunicodechar{α}{\ensuremath{\alpha}}
\newunicodechar{β}{\ensuremath{\beta}}
\newunicodechar{γ}{\ensuremath{\gamma}}
\newunicodechar{δ}{\ensuremath{\delta}}
\newunicodechar{ε}{\ensuremath{\varepsilon}}
\newunicodechar{θ}{\ensuremath{\theta}}
\newunicodechar{λ}{\ensuremath{\lambda}}
\newunicodechar{μ}{\ensuremath{\mu}}
\newunicodechar{σ}{\ensuremath{\sigma}}
\newunicodechar{τ}{\ensuremath{\tau}}
\newunicodechar{φ}{\ensuremath{\varphi}}
\newunicodechar{ψ}{\ensuremath{\psi}}
\newunicodechar{ω}{\ensuremath{\omega}}
\newunicodechar{Σ}{\ensuremath{\Sigma}}
% majuscules produites par \MakeUppercase dans les titres (α-aminés -> Α)
\newunicodechar{Α}{\ensuremath{\alpha}}
\newunicodechar{Β}{\ensuremath{\beta}}
\newunicodechar{Γ}{\ensuremath{\Gamma}}
\newunicodechar{Δ}{\ensuremath{\Delta}}
\newunicodechar{Ε}{\ensuremath{\varepsilon}}
\newunicodechar{Θ}{\ensuremath{\Theta}}
\newunicodechar{Λ}{\ensuremath{\Lambda}}
\newunicodechar{Μ}{\ensuremath{\mu}}
\newunicodechar{Τ}{\ensuremath{\tau}}
\newunicodechar{Φ}{\ensuremath{\Phi}}
\newunicodechar{Ψ}{\ensuremath{\Psi}}
\newunicodechar{Ω}{\ensuremath{\Omega}}
\newunicodechar{↦}{\ensuremath{\mapsto}}
\newunicodechar{∈}{\ensuremath{\in}}
\newunicodechar{∉}{\ensuremath{\notin}}
\newunicodechar{∧}{\ensuremath{\wedge}}
\newunicodechar{⇔}{\ensuremath{\Leftrightarrow}}
\newunicodechar{⟺}{\ensuremath{\Longleftrightarrow}}
\newunicodechar{⇒}{\ensuremath{\Rightarrow}}
\newunicodechar{⟹}{\ensuremath{\Longrightarrow}}
\newunicodechar{∘}{\ensuremath{\circ}}
\newunicodechar{∪}{\ensuremath{\cup}}
\newunicodechar{∩}{\ensuremath{\cap}}
\newunicodechar{≡}{\ensuremath{\equiv}}
\newunicodechar{⊂}{\ensuremath{\subset}}
\newunicodechar{⊥}{\ensuremath{\perp}}
\newunicodechar{∃}{\ensuremath{\exists}}
\newunicodechar{∀}{\ensuremath{\forall}}
\newunicodechar{∛}{\ensuremath{\sqrt[3]{\,}}}
\newunicodechar{∜}{\ensuremath{\sqrt[4]{\,}}}
\newunicodechar{⃗}{\ensuremath{{}^{\to}}}
\newunicodechar{ⁿ}{\ifmmode^{n}\else\textsuperscript{\ensuremath{n}}\fi}
\newunicodechar{ˣ}{\ifmmode^{x}\else\textsuperscript{\ensuremath{x}}\fi}
\newunicodechar{ᵇ}{\ifmmode^{b}\else\textsuperscript{\ensuremath{b}}\fi}
\newunicodechar{ʳ}{\ifmmode^{r}\else\textsuperscript{\ensuremath{r}}\fi}
\newunicodechar{ᵘ}{\ifmmode^{u}\else\textsuperscript{\ensuremath{u}}\fi}
\newunicodechar{ʸ}{\ifmmode^{y}\else\textsuperscript{\ensuremath{y}}\fi}
\newunicodechar{ᵃ}{\ifmmode^{a}\else\textsuperscript{\ensuremath{a}}\fi}
\newunicodechar{ᵉ}{\ifmmode^{e}\else\textsuperscript{\ensuremath{e}}\fi}
\newunicodechar{ᵏ}{\ifmmode^{k}\else\textsuperscript{\ensuremath{k}}\fi}
\newunicodechar{ᵖ}{\ifmmode^{p}\else\textsuperscript{\ensuremath{p}}\fi}
\newunicodechar{ᴺ}{\ifmmode^{N}\else\textsuperscript{\ensuremath{N}}\fi}
\newunicodechar{ᶜ}{\ifmmode^{c}\else\textsuperscript{\ensuremath{c}}\fi}
\newunicodechar{ʰ}{\ifmmode^{h}\else\textsuperscript{\ensuremath{h}}\fi}
\newunicodechar{ᵠ}{\ifmmode^{q}\else\textsuperscript{\ensuremath{q}}\fi}
\newunicodechar{ₙ}{\ifmmode_{n}\else\textsubscript{\ensuremath{n}}\fi}
\newunicodechar{ₐ}{\ifmmode_{a}\else\textsubscript{\ensuremath{a}}\fi}
\newunicodechar{ᵢ}{\ifmmode_{i}\else\textsubscript{\ensuremath{i}}\fi}
\newunicodechar{ᵣ}{\ifmmode_{r}\else\textsubscript{\ensuremath{r}}\fi}
\newunicodechar{ₖ}{\ifmmode_{k}\else\textsubscript{\ensuremath{k}}\fi}
\newunicodechar{ᵦ}{\ifmmode_{\beta}\else\textsubscript{\ensuremath{\beta}}\fi}
\newunicodechar{ₓ}{\ifmmode_{x}\else\textsubscript{\ensuremath{x}}\fi}
\newfontfamily\popdisplay{ArchivoBlack-Regular.ttf}[Path=../../../fonts/]
\newfontfamily\poplogo{SpaceGrotesk-Bold.ttf}[Path=../../../fonts/]
\newfontfamily\popsemi{PlusJakartaSans-SemiBold.ttf}[Path=../../../fonts/]
\newfontfamily\popxbold{PlusJakartaSans-ExtraBold.ttf}[Path=../../../fonts/]
\newfontfamily\popxboldls{PlusJakartaSans-ExtraBold.ttf}[Path=../../../fonts/,LetterSpace=3.0]

\setlist[enumerate]{leftmargin=1.7em,itemsep=5pt,topsep=4pt}
\setlist[itemize]{leftmargin=1.7em,itemsep=4pt,topsep=4pt,label={\color{poporange}\textbullet}}
\setlength{\parindent}{0pt}
\setlength{\parskip}{5pt}
\renewcommand{\arraystretch}{1.25}

% pgfplots : style Pop par défaut
\pgfplotsset{
  every axis/.append style={axis line style={popink,line width=1pt},tick style={popink},
    grid style={popline,line width=0.5pt},label style={font=\small},tick label style={font=\footnotesize}},
  popcurve/.style={poporange,line width=1.8pt,no markers,smooth},
}
% Même style disponible dans un \draw TikZ simple (hors pgfplots)
\tikzset{popcurve/.style={poporange,line width=1.8pt}}
\tikzset{popgrid/.style={popline,very thin}}
\colorlet{popcurve}{poporange} % le nom du style est parfois utilisé comme couleur

% ============================================================
% Pill / squiggle
% ============================================================
\newcommand{\poppill}[3]{%
  \tikz[baseline=-0.55ex]\node[draw=popink,line width=1.2pt,rounded corners=2.6pt,%
    fill=#2,inner xsep=6.5pt,inner ysep=3pt]{%
    \popxboldls\fontsize{7}{7}\selectfont\color{#3}\MakeUppercase{#1}};%
}
\newcommand{\popsquiggle}[1]{%
  \tikz[baseline=-0.4ex]{
    \draw[#1,line width=2.6pt,line cap=round]
      plot [smooth] coordinates {(0,0.09) (0.34,0.01) (0.68,0.21) (1.02,0.01) (1.36,0.10)};
  }%
}
% Mot-clé surligné (marqueur orange sous le texte)
\newcommand{\cle}[1]{{\popxbold\color{popdark}#1}}

% ============================================================
% En-tête / pied
% ============================================================
\pagestyle{fancy}
\fancyhf{}
\renewcommand{\headrulewidth}{0pt}
\renewcommand{\footrulewidth}{0pt}
\lhead{\raisebox{-0.15ex}{\poplogo\fontsize{13}{13}\selectfont\color{popink}OnBuch\color{poporange}+}}
\rhead{\poppill{\DOCMATIERE\ \textbullet\ \DOCNIVEAU\ \textbullet\ Leçon \DOCLECON}{white}{popink}}
\lfoot{\popxbold\fontsize{7.6}{7.6}\selectfont\color{popmuted}\MakeUppercase{Cours signé OnBuch+}\ \textbullet\ {\popsemi\DOCTITRE}}
\rfoot{\tikz[baseline=-0.6ex]\node[rounded corners=8.5pt,fill=popink,minimum height=17pt,minimum width=17pt,inner xsep=5pt,inner sep=0]{\popxbold\fontsize{8}{8}\selectfont\color{popcream}\thepage\,/\,\pageref*{LastPage}};}

% ============================================================
% Titres
% ============================================================
\newcommand{\chaptitre}[2]{%
  \begin{tcolorbox}[enhanced,colback=poporange,colframe=popink,boxrule=2.6pt,arc=11pt,
      drop shadow=popink,left=15pt,right=15pt,top=12pt,bottom=11pt,before skip=3pt,after skip=18pt]
    \poppill{#2}{popink}{popcream}\par\vspace{9pt}
    {\raggedright\hyphenpenalty=10000\exhyphenpenalty=10000\popdisplay\fontsize{22}{24}\selectfont\color{popcream}\MakeUppercase{#1}\par}
  \end{tcolorbox}
}
% Section numérotée : pastille orange avec le numéro + titre Archivo
\newcounter{popsec}
\newcommand{\coursec}[1]{%
  \stepcounter{popsec}\setcounter{popsub}{0}%
  \par\vspace{16pt}\noindent
  \begin{minipage}{\linewidth}
  \tikz[baseline=-0.7ex]\node[draw=popink,line width=1.6pt,fill=poporange,rounded corners=4pt,minimum size=22pt,inner sep=0]{\popdisplay\fontsize{12}{12}\selectfont\color{popcream}\thepopsec};%
  \hspace{8pt}{\popdisplay\fontsize{15}{17}\selectfont\color{popink}\MakeUppercase{#1}}\par
  \vspace{4pt}\noindent{\color{poporange}\rule{36pt}{3.6pt}}
  \end{minipage}\par\nopagebreak\vspace{8pt}
}
\newcounter{popsub}
\newcommand{\courssub}[1]{%
  \stepcounter{popsub}%
  \par\vspace{9pt}\noindent{\popxbold\fontsize{11.5}{13}\selectfont\color{popink}{\color{poporange}\thepopsec.\thepopsub}\hspace{6pt}#1}\par\nopagebreak\vspace{4pt}
}

% ============================================================
% Boîtes pédagogiques — même recette : fond blanc, bordure épaisse colorée,
% ombre dure encre, badge plein. Titre optionnel [#1].
% ============================================================
\tcbset{popbase/.style={breakable,enhanced,colback=white,boxrule=2.3pt,arc=9pt,
  shadow={3pt}{-3pt}{0pt}{popink},left=14pt,right=14pt,top=11pt,bottom=11pt,before skip=11pt,after skip=15pt,
  fontupper=\popsemi\fontsize{10.6}{14.5}\selectfont\color{popink},
  fontlower=\popsemi\fontsize{10.6}{14.5}\selectfont\color{popink}}}
% Coche dessinée (le glyphe ✓ n'existe pas dans les polices du document)
\newcommand{\popcheck}[1]{\tikz[baseline=-0.5ex]\draw[#1,line width=1.6pt,line cap=round,line join=round](-0.13,0.0)--(-0.04,-0.09)--(0.13,0.1);}
\providecommand{\coche}{\popcheck{popgreen}}
\providecommand{\resultat}[1]{\par\smallskip\noindent\fcolorbox{popgreen}{popgreenL}{\parbox{\dimexpr\linewidth-2\fboxsep-2\fboxrule\relax}{\textbf{Résultat.} #1}}\par\smallskip}
\newcommand{\popboxhead}[3]{\poppill{#1}{#2}{white}\ifx\relax#3\relax\else\hspace{6pt}{\popxbold\fontsize{10.6}{12}\selectfont\color{popink}#3}\fi\par\vspace{7pt}}

\newtcolorbox{aretenir}[1][]{popbase,colframe=popgreen,before upper={\popboxhead{À retenir}{popgreen}{#1}}}
% Texte en espagnol (document de lecture, dialogue, poème, extrait) : cadre
% d'étude. Un titre optionnel (source, auteur) s'affiche dans l'en-tête.
\newtcolorbox{textebox}[1][]{popbase,colframe=popblue,before upper={\popboxhead{Texto}{popblue}{#1}\begin{otherlanguage}{spanish}},after upper={\end{otherlanguage}}}
% Marques ✓ / ✗ (forme correcte / fautive) souvent utilisées par les modèles
\providecommand{\cross}{\textcolor{poppink}{\ensuremath{\times}}}
\providecommand{\xmark}{\cross}\providecommand{\crossmark}{\cross}\providecommand{\cmark}{\ensuremath{\checkmark}}
% Ligne de réponse / justification à compléter (inventée par les modèles)
\providecommand{\justification}[1]{\par\noindent\textit{Justification~:} \dotfill\par}
% \textipa (package tipa non chargé) : la police ne couvre pas l'API -> simple passage
\providecommand{\textipa}[1]{#1}
% Mot ou phrase en espagnol à l'intérieur d'un paragraphe français
\newcommand{\esp}[1]{\begingroup\selectlanguage{spanish}\itshape #1\endgroup}
\newtcolorbox{exemplebox}[1][]{popbase,colframe=poporange,before upper={\popboxhead{Exemple}{poporange}{#1}}}
\newtcolorbox{definition}[1][]{popbase,colframe=popblue,before upper={\popboxhead{Définition}{popblue}{#1}}}
\newtcolorbox{propriete}[1][]{popbase,colframe=poppurple,before upper={\popboxhead{Propriété}{poppurple}{#1}}}
\newtcolorbox{methode}[1][]{popbase,colframe=popgold,before upper={\popboxhead{Méthode}{popgold}{#1}}}
\newtcolorbox{attention}[1][]{popbase,colframe=poppink,before upper={\popboxhead{Attention}{poppink}{#1}}}
\newtcolorbox{experience}[1][]{popbase,colframe=popblue,colback=popblueL,before upper={\popboxhead{Expérience}{popblue}{#1}}}
% « Le savais-tu ? » : carte violette pleine (contexte, culture, Cameroun)
\newtcolorbox{savaistu}[1][]{popbase,colback=poppurple,colframe=popink,
  fontupper=\popsemi\fontsize{10.4}{14.2}\selectfont\color{white},
  before upper={\poppill{Le savais-tu\,?}{white}{poppurple}\ifx\relax#1\relax\else\hspace{6pt}{\popxbold\color{white}#1}\fi\par\vspace{7pt}}}
% Exercice résolu : énoncé en haut, \tcblower puis la solution sur fond crème
\newcounter{popexo}\newcounter{popexr}
\newtcolorbox{exoresolu}[1][]{popbase,colframe=poporange,colbacklower=popcream,
  segmentation style={popink,line width=1.2pt,dashed},
  before upper={\stepcounter{popexr}\popboxhead{Exercice résolu \thepopexr}{poporange}{#1}},
  before lower={\poppill{Solution}{popink}{popcream}\par\vspace{6pt}}}
% Exercice (non résolu en place) — la correction est dans la partie Corrigés
\newtcolorbox{exercice}[1][]{popbase,colframe=popink,
  before upper={\stepcounter{popexo}\popboxhead{Exercice \thepopexo}{popink}{#1}}}
% Corrigé
\newtcolorbox{corrige}[1][]{popbase,colframe=popgreen,colback=popgreenL,
  before upper={\popboxhead{Corrigé}{popgreen}{#1}}}
% Objectifs / prérequis / bilan
\newtcolorbox{objectifs}{popbase,colframe=popink,colback=poporangeL,
  before upper={\popboxhead{Objectifs de la leçon}{popink}{}}}
\newtcolorbox{prerequis}{popbase,colframe=popink,colback=popblueL,
  before upper={\popboxhead{Prérequis}{popblue}{}}}
\newtcolorbox{bilan}{popbase,colframe=popink,colback=white,boxrule=2.8pt,
  before upper={\popboxhead{Fiche bilan}{poporange}{L'essentiel à connaître pour l'examen}}}
% Figure encadrée avec légende
\newtcolorbox{popfigure}[1][]{enhanced,breakable=false,colback=white,colframe=popline,boxrule=1.4pt,arc=8pt,
  left=10pt,right=10pt,top=10pt,bottom=8pt,before skip=10pt,after skip=12pt,halign=center,
  lower separated=false,halign lower=center,
  fontlower=\popsemi\fontsize{9}{11.5}\selectfont\color{popmuted},
  #1}
\newcommand{\legende}[1]{\par\vspace{6pt}{\popsemi\fontsize{9}{11.5}\selectfont\color{popmuted}#1}}

% ============================================================
% Signature OnBuch+
% ============================================================
\newcommand{\poplogoinline}[1]{{\poplogo\fontsize{#1}{#1}\selectfont\color{popink}OnBuch\color{poporange}+}}
\newcommand{\signatureonbuch}{%
  \begin{tcolorbox}[enhanced,colback=popink,colframe=popink,boxrule=2.2pt,arc=10pt,
      drop shadow=poporange,left=14pt,right=14pt,top=10pt,bottom=10pt,before skip=0pt,after skip=0pt]
    \tikz[baseline=-0.7ex]\node[circle,fill=popgreen,draw=popcream,line width=1.3pt,minimum size=22pt,inner sep=0]{\popcheck{white}};%
    \hspace{9pt}\begin{minipage}[c]{0.62\linewidth}
      {\popxbold\fontsize{12}{13}\selectfont\color{popcream}Signé\ \textbullet\ {\poplogo OnBuch\color{poporange}+}}\par\vspace{2pt}
      {\raggedright\popsemi\fontsize{8}{10.5}\selectfont\color{popline}\MakeUppercase{Équipe pédagogique OnBuch+}\\\MakeUppercase{Conforme au programme officiel MINESEC}\par}
    \end{minipage}\hfill
    {\poplogo\fontsize{22}{22}\selectfont\color{popcream}OnBuch\color{poporange}+}
  \end{tcolorbox}%
}

% ============================================================
% Page de garde
% #1 titre · #2 matière · #3 niveau · #4 module · #5 n° leçon · #6 durée ·
% #7 description / objectifs (texte ou itemize)
% ============================================================
\newcommand{\pagegarde}[7]{%
  \thispagestyle{empty}
  \newgeometry{a4paper,margin=1.7cm,top=1.6cm,bottom=1.4cm}%
  \begin{tikzpicture}[remember picture,overlay]
    \fill[poporange!18] ([xshift=-3.2cm,yshift=-3.6cm]current page.north east) circle (4.6cm);
    \fill[poppurple!14] ([xshift=2.4cm,yshift=7.6cm]current page.south west) circle (3.8cm);
  \end{tikzpicture}%
  {\poplogo\fontsize{30}{30}\selectfont\color{popink}OnBuch\color{poporange}+}\hfill
  \raisebox{0.6ex}{\poppill{Cours signé \textbullet\ Premium}{popink}{popcream}}\par
  \vspace{1pt}\popsquiggle{poppurple}\par
  \vspace{8pt}
  \begin{tcolorbox}[enhanced,colback=poporange,colframe=popink,boxrule=2.8pt,arc=13pt,
      drop shadow=popink,left=18pt,right=18pt,top=12pt,bottom=13pt,after skip=10pt]
    \poppill{#2 \textbullet\ #3}{popink}{popcream}\hspace{5pt}\poppill{#4}{white}{popink}\par\vspace{8pt}
    {\popdisplay\fontsize{14}{14}\selectfont\color{popink}LEÇON #5}\par\vspace{4pt}
    {\raggedright\hyphenpenalty=10000\exhyphenpenalty=10000\popdisplay\fontsize{25}{27}\selectfont\color{popcream}\MakeUppercase{#1}\par}
    \vspace{8pt}
    {\popxbold\fontsize{9.5}{9.5}\selectfont\color{popink}\MakeUppercase{Durée conseillée : #6}}
  \end{tcolorbox}
  \begin{tcolorbox}[enhanced,colback=white,colframe=popink,boxrule=2.2pt,arc=10pt,
      drop shadow=popink,left=16pt,right=16pt,top=10pt,bottom=10pt,after skip=8pt]
    \poppill{Dans cette leçon}{poppurple}{white}\par\vspace{5pt}
    \popsemi\fontsize{9.3}{12.6}\selectfont\color{popink}#7
  \end{tcolorbox}
  \noindent\begin{minipage}[t]{0.487\linewidth}
    \begin{tcolorbox}[enhanced,colback=popgreen,colframe=popink,boxrule=2.2pt,arc=10pt,
        drop shadow=popink,left=11pt,right=11pt,top=10pt,bottom=10pt,height=3.1cm,valign=center]
      \tcbox[colback=white,colframe=popink,boxrule=1.3pt,arc=5pt,boxsep=0pt,nobeforeafter,
        left=3pt,right=3pt,top=3pt,bottom=3pt,valign=center]{\qrcode[height=1.9cm]{https://play.google.com/store/apps/details?id=com.luvvix.onbuch&hl=fr}}%
      \hspace{8pt}\begin{minipage}[c]{0.5\linewidth}
        {\popdisplay\fontsize{11}{12.5}\selectfont\color{white}\MakeUppercase{Télécharge OnBuch}}\par\vspace{4pt}
        {\raggedright\popsemi\fontsize{8.4}{11}\selectfont\color{white}Gratuit \textbullet\ Google Play\\Tuteur IA, annales, concours, résultats\par}
      \end{minipage}
    \end{tcolorbox}
  \end{minipage}\hfill
  \begin{minipage}[t]{0.487\linewidth}
    \begin{tcolorbox}[enhanced,colback=poppurple,colframe=popink,boxrule=2.2pt,arc=10pt,
        drop shadow=popink,left=11pt,right=11pt,top=10pt,bottom=10pt,height=3.1cm,valign=center]
      \tikz[baseline=-0.7ex]\node[circle,fill=white,draw=popink,line width=1.3pt,minimum size=1.6cm,inner sep=0]{\popdisplay\fontsize{24}{24}\selectfont\color{poppurple}+};%
      \hspace{8pt}\begin{minipage}[c]{0.6\linewidth}
        {\popdisplay\fontsize{11}{12.5}\selectfont\color{white}\MakeUppercase{Passe à OnBuch+}}\par\vspace{4pt}
        {\raggedright\popsemi\fontsize{8.4}{11}\selectfont\color{white}Tous les cours signés, Tuteur IA illimité, fiches PDF — onglet OnBuch+ de l'app\par}
      \end{minipage}
    \end{tcolorbox}
  \end{minipage}
  \vspace{8pt}
  \popcontenu
  \vfill
  \signatureonbuch
  \restoregeometry
  \newpage
}

% Rangée « Ce cours contient » (page de garde)
\newcommand{\poptile}[3]{% #1 couleur #2 gros texte #3 légende
  \begin{tcolorbox}[enhanced,colback=#1,colframe=popink,boxrule=2pt,arc=9pt,shadow={2.5pt}{-2.5pt}{0pt}{popink},
      left=6pt,right=6pt,top=9pt,bottom=9pt,halign=center,valign=center,height=1.75cm,width=\linewidth,nobeforeafter]
    {\popdisplay\fontsize{12.5}{13}\selectfont\color{white}\MakeUppercase{#2}}\par\vspace{3pt}
    {\centering\popsemi\fontsize{7.8}{9.5}\selectfont\color{white}#3\par}
  \end{tcolorbox}}
\newcommand{\popcontenu}{%
  \par\noindent{\popxboldls\fontsize{7.5}{7.5}\selectfont\color{popmuted}CE COURS SIGNÉ CONTIENT}\par\vspace{7pt}
  \noindent\begin{minipage}[t]{0.235\linewidth}\poptile{popblue}{Cours}{complet, illustré, pas à pas}\end{minipage}\hfill
  \begin{minipage}[t]{0.235\linewidth}\poptile{poporange}{Méthodes}{et exercices résolus}\end{minipage}\hfill
  \begin{minipage}[t]{0.235\linewidth}\poptile{poppink}{Exercices}{type examen + corrigés}\end{minipage}\hfill
  \begin{minipage}[t]{0.235\linewidth}\poptile{popgreen}{Fiche bilan}{l'essentiel à retenir}\end{minipage}\par}

% Sommaire
\newtcolorbox{sommaire}{enhanced,colback=white,colframe=popink,boxrule=2.2pt,arc=10pt,
  drop shadow=popink,left=18pt,right=18pt,top=13pt,bottom=13pt,before skip=6pt,after skip=20pt,
  before upper={\poppill{Sommaire}{poppurple}{white}\par\vspace{9pt}\popsemi\fontsize{10.6}{14.5}\selectfont\color{popink}}}

% Signature de fin de cours — poussée tout en bas de la dernière page
\newcommand{\findecours}{%
  \par\vfill\nopagebreak
  \begin{center}\popsquiggle{poporange}\end{center}\vspace{4pt}
  \signatureonbuch
}
\AtBeginDocument{\def\llbracket{\mathopen{[\mkern-3.5mu[}}\def\rrbracket{\mathclose{]\mkern-3.5mu]}}} % intervalles d'entiers (glyphe ⟦ absent de la police)
% Glyphes absents de Fira Math : redéfinis à partir de glyphes disponibles
\AtBeginDocument{%
  \def\setminus{\mathbin{\backslash}}\let\smallsetminus\setminus
  \def\vdots{\mathord{\vbox{\baselineskip3.2pt\lineskiplimit0pt\kern4pt\hbox{.}\hbox{.}\hbox{.}}}}%
  \def\ddots{\mathinner{\mkern1mu\raise6.5pt\hbox{.}\mkern2mu\raise3.6pt\hbox{.}\mkern2mu\raise0.7pt\hbox{.}\mkern1mu}}%
  \def\triangle{\mathord{\scalebox{0.9}{$\symup{\Delta}$}}}%
  \def\nabla{\mathord{\raisebox{\depth}{\scalebox{1}[-1]{$\symup{\Delta}$}}}}%
  \def\checkmark{\popcheck{popdark}}%
  \def\star{\mathbin{\ast}}\def\bigstar{\mathord{\ast}}%
  \def\diamond{\mathbin{\scalebox{0.6}{$\Diamond$}}}%
  \def\bigcup{\mathop{\mathchoice{\raisebox{-0.3em}{\scalebox{1.7}{$\cup$}}}{\scalebox{1.25}{$\cup$}}{\cup}{\cup}}\displaylimits}%
  \def\bigcap{\mathop{\mathchoice{\raisebox{-0.3em}{\scalebox{1.7}{$\cap$}}}{\scalebox{1.25}{$\cap$}}{\cap}{\cap}}\displaylimits}%
  \def\lesseqqgtr{\mathrel{\lessgtr}}\def\bigoplus{\mathop{\oplus}\displaylimits}%
}
\providecommand{\ssi}{si et seulement si} % abréviation fréquente
\AtBeginDocument{\providecommand{\wideparen}{\overparen}} % arc de cercle

\begin{document}

\pagegarde{Traductions : l'ordre, l'interdiction, « chez », « devenir » ; debates : educación de los niños y NTIC}{Espagnol}{Tle A}{Módulo 1 : Vie familiale et intégration sociale}{E06}{2 h}{Leçon mixte de Terminale A : maîtriser l'expression de l'ordre et de l'interdiction en espagnol, choisir le bon équivalent de « chez » et de « devenir », puis mobiliser ces outils dans un débat argumenté sur l'éducation des enfants et les nouvelles technologies.
\vspace{4pt}
\begin{multicols}{2}\small
\begin{itemize}
\item À la fin de cette leçon, je sais exprimer un ordre affirmatif et négatif en espagnol (impératif, subjonctif, infinitif).
\item À la fin de cette leçon, je sais exprimer l'interdiction par l'infinitif, « prohibido + infinitif » et « se prohíbe + infinitif ».
\item À la fin de cette leçon, je sais traduire « chez » selon le contexte : en casa de, a casa de, en el/la..., entre los...
\item À la fin de cette leçon, je sais choisir le bon équivalent de « devenir » : ponerse, volverse, hacerse, convertirse en, llegar a ser, quedarse.
\item À la fin de cette leçon, je sais repérer et éviter les pièges de traduction fréquents des francophones.
\item À la fin de cette leçon, je sais comprendre et commenter un texte sur l'éducation des enfants et les NTIC.
\end{itemize}\end{multicols}}

\begin{sommaire}
\begin{multicols}{2}
\begin{enumerate}
\item Texte support : « Los niños y las pantallas »
\item Lexique : éducation des enfants et NTIC
\item Grammaire 1 : l'ordre et l'interdiction
\item Grammaire 2 : traduire « chez »
\item Grammaire 3 : traduire « devenir »
\item Débat : la educación de los niños y las NTIC
\item Activité d'intégration
\item Méthodes et exercices résolus
\item Exercices
\item Corrigés des exercices
\item Fiche bilan
\end{enumerate}
\end{multicols}
\end{sommaire}

\begin{objectifs}
\begin{itemize}
\item À la fin de cette leçon, je sais exprimer un ordre affirmatif et négatif en espagnol (impératif, subjonctif, infinitif).
\item À la fin de cette leçon, je sais exprimer l'interdiction par l'infinitif, « prohibido + infinitif » et « se prohíbe + infinitif ».
\item À la fin de cette leçon, je sais traduire « chez » selon le contexte : en casa de, a casa de, en el/la..., entre los...
\item À la fin de cette leçon, je sais choisir le bon équivalent de « devenir » : ponerse, volverse, hacerse, convertirse en, llegar a ser, quedarse.
\item À la fin de cette leçon, je sais repérer et éviter les pièges de traduction fréquents des francophones.
\item À la fin de cette leçon, je sais comprendre et commenter un texte sur l'éducation des enfants et les NTIC.
\item À la fin de cette leçon, je sais argumenter oralement et par écrit dans un débat (donner une opinion, la justifier, réfuter).
\item À la fin de cette leçon, je sais traduire un court texte français mobilisant l'ordre, l'interdiction, « chez » et « devenir ».
\end{itemize}
\end{objectifs}

\begin{prerequis}
\begin{itemize}
\item L'impératif affirmatif et négatif (classe de Première)
\item Le présent du subjonctif : formation et emplois de base
\item Les pronoms compléments et leur place avec l'impératif (dímelo, no me lo digas)
\item Le lexique de la famille et de l'éducation (Módulo 1)
\item Les structures d'opinion : pienso que, me parece que, estoy de acuerdo
\end{itemize}
\end{prerequis}

\newpage

\coursec{Texte support : « Los niños y las pantallas »}

\courssub{Situation d'accroche et problématique}

À l'entrée du lycée de Nkolbisson, un panneau affiche en lettres rouges : \esp{PROHIBIDO FUMAR}. Dans la cour, un père camerounais dit à son fils : \esp{¡No uses el teléfono en la mesa!} « N'utilise pas le téléphone à table ! ». Le soir, toute la famille se réunit \esp{en casa de tía María} « chez tante María » pour débattre : les enfants \esp{se vuelven} « deviennent »-ils trop dépendants des écrans ?

Ordre, interdiction, « chez », « devenir » : voilà quatre outils grammaticaux indispensables pour traduire et débattre au Bac. Avant de les étudier en détail dans les sections suivantes, nous allons les \emph{observer en situation} dans un texte authentique. C'est la démarche du commentaire de texte : lire, comprendre, repérer, analyser.

\textbf{Problématique :} comment la langue espagnole exprime-t-elle l'autorité (l'ordre, l'interdiction) et le changement (devenir) dans le débat familial sur les nouvelles technologies ?

\courssub{Lecture du texte}

\begin{textebox}[Los niños y las pantallas]
¡Apaga el teléfono ahora mismo! — grita mamá desde la cocina. En nuestra casa de Yaoundé, la cena es un momento sagrado: se prohíbe usar el móvil durante la comida, pero nadie respeta la regla.

Mi hermano pequeño se ha vuelto muy callado desde que tiene una tableta. Antes jugaba al fútbol en el barrio; ahora pasa horas delante de la pantalla. Mi hermana, en cambio, se ha hecho más responsable: usa su ordenador para estudiar y ha llegado a ser la primera de su clase.

Los fines de semana, los niños van a casa de mi prima Sandrine, que vive cerca del mercado. En casa de ella no hay reglas: cada uno hace lo que quiere con su teléfono. Cuando vuelven, se ponen insoportables y no quieren ni hablar con nosotros.

— ¡Ven aquí y deja ese aparato! — le dice papá a mi hermano. Pero él no obedece. Mi madre cree que hay que prohibir las pantallas por completo; mi padre piensa que es mejor educar que castigar. «La brecha digital entre los padres y los hijos es enorme — dice tía María —, pero criar a un niño siempre ha sido difícil, con o sin tecnología.»

Yo creo que las nuevas tecnologías no son ni buenas ni malas: todo depende del uso que hacemos de ellas. Si los niños se convierten en esclavos de las pantallas, la culpa no es de las máquinas, sino de los adultos que no saben poner límites.
\end{textebox}

\textbf{Traduction du texte.} « Éteins le téléphone tout de suite ! — crie maman depuis la cuisine. Chez nous, à Yaoundé, le dîner est un moment sacré : il est interdit d'utiliser le portable pendant le repas, mais personne ne respecte la règle. Mon petit frère est devenu très silencieux depuis qu'il a une tablette. Avant, il jouait au football dans le quartier ; maintenant, il passe des heures devant l'écran. Ma sœur, en revanche, est devenue plus responsable : elle utilise son ordinateur pour étudier et elle est devenue la première de sa classe. Le week-end, les enfants vont chez ma cousine Sandrine, qui habite près du marché. Chez elle, il n'y a pas de règles : chacun fait ce qu'il veut avec son téléphone. Quand ils reviennent, ils deviennent insupportables et ne veulent même pas nous parler. — Viens ici et laisse cet appareil ! — dit papa à mon frère. Mais il n'obéit pas. Ma mère pense qu'il faut interdire complètement les écrans ; mon père pense qu'il vaut mieux éduquer que punir. "La fracture numérique entre les parents et les enfants est énorme — dit tante María —, mais élever un enfant a toujours été difficile, avec ou sans technologie." Moi, je crois que les nouvelles technologies ne sont ni bonnes ni mauvaises : tout dépend de l'usage que nous en faisons. Si les enfants deviennent esclaves des écrans, la faute n'est pas aux machines, mais aux adultes qui ne savent pas fixer de limites. »

\courssub{Glossaire et lexique}

\begin{definition}[Vocabulaire clé du texte]
\begin{itemize}
\item \esp{la pantalla} : l'écran ; \esp{delante de la pantalla} : devant l'écran.
\item \esp{prohibir} : interdire ; \esp{se prohíbe usar el móvil} : il est interdit d'utiliser le portable.
\item \esp{castigar} : punir ; \esp{el castigo} : la punition.
\item \esp{volverse} : devenir (changement psychologique, souvent involontaire) ; \esp{se ha vuelto callado} : il est devenu silencieux.
\item \esp{criar} : élever (un enfant) ; \esp{la crianza} : l'éducation, l'élevage.
\item \esp{la brecha digital} : la fracture numérique.
\item \esp{poner límites} : fixer des limites ; \esp{respetar la regla} : respecter la règle.
\item \esp{hacerse} : devenir (par effort, volonté) ; \esp{llegar a ser} : devenir (aboutissement) ; \esp{convertirse en} : se transformer en ; \esp{ponerse} : devenir (changement passager, physique ou d'humeur).
\end{itemize}
\end{definition}

\begin{center}
\begin{tabularx}{\linewidth}{|X|X|X|}\hline
\rowcolor{popblueL} Español & Français & Remarque \\ \hline
la pantalla & l'écran & nom féminin \\ \hline
prohibir & interdire & \esp{prohibido + infinitivo} sur les panneaux \\ \hline
se prohíbe & il est interdit & tournure impersonnelle \\ \hline
castigar & punir & \esp{el castigo} : la punition \\ \hline
criar & élever (un enfant) & \esp{educar} = éduquer, instruire \\ \hline
la brecha digital & la fracture numérique & \esp{brecha} = brèche, fossé \\ \hline
el móvil / el celular & le portable & \esp{móvil} en Espagne, \esp{celular} en Amérique latine \\ \hline
en casa de & chez (lieu, état) & \esp{a casa de} avec les verbes de mouvement \\ \hline
\end{tabularx}
\end{center}

\begin{savaistu}[Móvil ou celular ?]
Au Cameroun comme en Espagne, les familles débattent de la \esp{brecha digital} entre les générations : les grands-parents nés avant Internet, les parents qui ont découvert le téléphone portable adultes, et les enfants « natifs du numérique ». Petit détail de vocabulaire : le mot \esp{móvil} est le terme utilisé en Espagne, tandis que \esp{celular} domine en Amérique latine (et en Guinée équatoriale, on entend les deux). Au Bac, les deux sont acceptés, mais sois cohérent dans ta copie !
\end{savaistu}

\courssub{Repérage dans le texte}

Avant d'analyser, repérons les quatre phénomènes grammaticaux de la leçon.

\textbf{1. Deux ordres à l'impératif :}
\begin{itemize}
\item \esp{¡Apaga el teléfono ahora mismo!} « Éteins le téléphone tout de suite ! » — impératif affirmatif, 2\up{e} personne du singulier (\esp{apagar}).
\item \esp{¡Ven aquí y deja ese aparato!} « Viens ici et laisse cet appareil ! » — deux impératifs affirmatifs : \esp{ven} (irrégulier, de \esp{venir}) et \esp{deja} (de \esp{dejar}).
\end{itemize}

\textbf{2. Une interdiction impersonnelle :}
\begin{itemize}
\item \esp{Se prohíbe usar el móvil durante la comida.} « Il est interdit d'utiliser le portable pendant le repas. » — tournure impersonnelle avec \esp{se} + 3\up{e} personne du singulier.
\end{itemize}

\textbf{3. Deux emplois de « chez » :}
\begin{itemize}
\item \esp{van a casa de mi prima Sandrine} « ils vont chez ma cousine Sandrine » — \esp{a casa de} avec un verbe de mouvement (\esp{ir}).
\item \esp{En casa de ella no hay reglas} « Chez elle, il n'y a pas de règles » — \esp{en casa de} avec un verbe d'état (\esp{haber}).
\end{itemize}

\textbf{4. Trois verbes de « devenir » :}
\begin{itemize}
\item \esp{se ha vuelto muy callado} « il est devenu très silencieux » — \esp{volverse} + adjectif (changement psychologique durable).
\item \esp{se ha hecho más responsable} « elle est devenue plus responsable » — \esp{hacerse} + adjectif (résultat d'un effort).
\item \esp{se convierten en esclavos} « ils deviennent esclaves » — \esp{convertirse en} + nom (transformation profonde).
\end{itemize}

\courssub{Questions de compréhension}

\textbf{A. Compréhension globale}
\begin{enumerate}
\item ¿Quién habla en el texto? ¿Quiénes son los personajes? (Qui parle ? Quels sont les personnages ?)
\item ¿Cuál es el problema principal de esta familia? (Quel est le problème principal ?)
\item ¿Qué soluciones proponen la madre, el padre y la tía María? (Quelles solutions chacun propose-t-il ?)
\end{enumerate}

\textbf{B. Compréhension fine : arguments pour et contre les NTIC}
\begin{enumerate}
\item Releva en el texto dos argumentos \emph{contra} las pantallas. (Relève deux arguments contre les écrans.)
\item Releva un argumento \emph{a favor} de las nuevas tecnologías. (Relève un argument en faveur des NTIC.)
\item ¿Qué opina el narrador al final? ¿Estás de acuerdo? (Quelle est l'opinion du narrateur ? Es-tu d'accord ?)
\end{enumerate}

\begin{corrige}[Éléments de réponse]
\textbf{A.} 1. Le narrateur est un enfant (ou adolescent) de la famille ; les personnages sont la mère, le père, le frère, la sœur, la cousine Sandrine et tante María. 2. Le problème : la place des écrans dans la vie familiale et l'éducation des enfants. 3. La mère veut interdire totalement les écrans ; le père préfère éduquer plutôt que punir ; tante María rappelle qu'élever un enfant a toujours été difficile.

\textbf{B.} Contre : le frère \esp{se ha vuelto muy callado} et ne joue plus dehors ; les enfants reviennent \esp{insoportables} de chez la cousine. Pour : la sœur utilise l'ordinateur pour étudier et \esp{ha llegado a ser la primera de su clase}. Le narrateur conclut que tout dépend de l'usage : les adultes doivent \esp{poner límites}.
\end{corrige}

\textbf{C. Questions de grammaire}
\begin{enumerate}
\item Identifica el modo y el tiempo de : \esp{apaga}, \esp{ven}, \esp{deja}, \esp{se prohíbe}. (Identifie le mode et le temps.)
\item ¿Por qué se usa \esp{a casa de} con \esp{van} y \esp{en casa de} con \esp{hay}? (Pourquoi \esp{a casa de} avec « vont » et \esp{en casa de} avec « il y a » ?)
\end{enumerate}

\begin{corrige}[Éléments de réponse]
1. \esp{Apaga}, \esp{ven}, \esp{deja} : mode impératif (l'impératif n'a qu'un « présent »), 2\up{e} personne du singulier. \esp{Se prohíbe} : mode indicatif, présent, 3\up{e} personne du singulier, tournure impersonnelle. 2. \esp{A casa de} s'emploie avec les verbes de \emph{mouvement} (\esp{ir}, \esp{venir}) ; \esp{en casa de} avec les verbes d'\emph{état} (\esp{estar}, \esp{haber}, \esp{comer}).
\end{corrige}

\courssub{Méthode : réussir le commentaire de texte au Bac}

\begin{methode}[Les étapes du commentaire de texte]
\begin{enumerate}
\item \textbf{Lire deux fois} le texte : une fois pour le sens global, une fois en soulignant les mots-clés et les faits de grammaire.
\item \textbf{Situer le texte} dans l'introduction : qui parle ? à qui ? de quoi ? dans quelles circonstances ? Présente le thème général (ici : l'éducation des enfants face aux nouvelles technologies).
\item \textbf{Dégager la problématique} : une question que le texte pose (ici : faut-il interdire les écrans ou éduquer à leur usage ?).
\item \textbf{Construire un plan en 2 ou 3 parties} : par exemple I. Le problème (les effets des écrans sur les enfants) ; II. Les positions des adultes ; III. La solution proposée par le narrateur.
\item \textbf{Citer et traduire} chaque passage analysé : une citation espagnole doit toujours être suivie de sa traduction française et d'un commentaire grammatical ou stylistique.
\item \textbf{Conclure} en répondant à la problématique et en ouvrant sur une question plus large (la fracture numérique au Cameroun, le rôle des parents...).
\end{enumerate}
\end{methode}

\begin{popfigure}
\begin{center}
\begin{tikzpicture}[
intro/.style={rectangle, rounded corners, draw=popblue, fill=popblueL, align=center, minimum width=3.2cm, minimum height=1.3cm, font=\small},
des/.style={rectangle, rounded corners, draw=poporange, fill=poporangeL, align=center, minimum width=3.2cm, minimum height=1.3cm, font=\small},
conc/.style={rectangle, rounded corners, draw=popgreen, fill=popgreenL, align=center, minimum width=3.2cm, minimum height=1.3cm, font=\small},
flecha/.style={-{Stealth[length=3mm]}, thick, popdark}]
\node[intro, align=center] (i){INTRODUCCIÓN\\ situación, autor,\\ tema, problemática};
\node[des, right=1.4cm of i, align=center] (d){DESARROLLO\\ plan en 2 o 3 partes\\ citar + traducir};
\node[conc, right=1.4cm of d, align=center] (c){CONCLUSIÓN\\ respuesta a la\\ problemática + apertura};
\draw[flecha] (i) -- (d);
\draw[flecha] (d) -- (c);
\end{tikzpicture}
\end{center}
\legende{La structure du commentaire de texte : introducción, desarrollo, conclusión.}
\end{popfigure}

\courssub{Commentaire modèle rédigé}

\begin{exemplebox}[La educación de los hijos frente a las nuevas tecnologías]
\textbf{Introducción.} Este texto es un relato en primera persona: un joven de Yaoundé describe los conflictos de su familia a causa de las pantallas. El tema es la educación de los hijos frente a las nuevas tecnologías, un problema actual tanto en Camerún como en España. La problemática es la siguiente: ¿hay que prohibir las pantallas o educar a los niños en su uso?

\textbf{Desarrollo.} En primer lugar, el narrador muestra los efectos negativos de las pantallas: su hermano \esp{se ha vuelto muy callado} (« il est devenu très silencieux »), es decir, ha cambiado de carácter. La madre reacciona con órdenes al imperativo: \esp{¡Apaga el teléfono!} (« Éteins le téléphone ! »), y con una prohibición impersonal: \esp{se prohíbe usar el móvil} (« il est interdit d'utiliser le portable »). En segundo lugar, el texto presenta los aspectos positivos: la hermana \esp{se ha hecho más responsable} (« elle est devenue plus responsable ») gracias al ordenador y \esp{ha llegado a ser la primera de su clase} (« elle est devenue la première de sa classe »). Finalmente, el debate familiar opone a la madre, partidaria de prohibir, al padre, que prefiere \esp{educar que castigar} (« éduquer plutôt que punir »). El narrador concluye que todo depende del uso y que los adultos deben \esp{poner límites} (« fixer des limites »).

\textbf{Conclusión.} El texto defiende una posición equilibrada: ni prohibición total ni libertad total. Esta reflexión puede extenderse a la \esp{brecha digital} entre generaciones: ¿cómo pueden los padres guiar a sus hijos en un mundo que ellos mismos descubren?
\end{exemplebox}

\begin{exoresolu}[Analyse grammaticale d'une phrase du texte]
Énoncé. Relève dans la phrase \esp{En casa de ella no hay reglas: cada uno hace lo que quiere} : a) l'équivalent de « chez » ; b) le temps et le mode des verbes ; c) traduis la phrase entière en français.

\tcblower

Solution détaillée. a) « Chez » est traduit par \esp{en casa de} car le verbe \esp{hay} (il y a) exprime un \emph{état}, non un mouvement : avec un verbe de mouvement, on aurait \esp{a casa de}. b) \esp{Hay} : présent de l'indicatif, forme impersonnelle de \esp{haber} ; \esp{hace} et \esp{quiere} : présent de l'indicatif, 3\up{e} personne du singulier (\esp{hacer}, \esp{querer}). c) Traduction : « Chez elle, il n'y a pas de règles : chacun fait ce qu'il veut. »
\end{exoresolu}

\begin{attention}[Pièges fréquents des francophones]
\begin{itemize}
\item Ne traduis jamais « chez » par \esp{en la casa de} quand il s'agit simplement du foyer : on dit \esp{en casa de mi prima} (sans article). \esp{En la casa de} insisterait sur le bâtiment.
\item \esp{Prohibir} se conjugue avec une diphtongaison : \esp{prohíbo, prohíbes, prohíbe, prohibimos, prohibís, prohíben}. Attention à l'accent : \esp{se prohíbe}.
\item « Devenir » ne se traduit pas par un seul verbe espagnol : \esp{volverse} (changement subi), \esp{hacerse} (effort), \esp{ponerse} (passager), \esp{convertirse en} + nom, \esp{llegar a ser} (aboutissement). Le choix dépend du sens !
\item Faux ami : \esp{educar} signifie « éduquer » mais aussi « élever » au sens large ; \esp{criar} = élever, nourrir (un enfant, un animal).
\end{itemize}
\end{attention}

\begin{popfigure}
\begin{center}
\begin{tikzpicture}[mindmap, grow cyclic,
root concept/.append style={concept color=popblue, text=white, font=\bfseries},
level 1 concept/.append style={sibling angle=90, level distance=3.4cm, font=\small},
concept color=popblue]
\node[concept] {Educación y NTIC}
child[concept color=popgreen] { node[concept] {ventajas} child[level distance=2.6cm, sibling angle=60] { node[concept, font=\scriptsize] {estudiar, informarse} } child[level distance=2.6cm, sibling angle=60] { node[concept, font=\scriptsize] {comunicarse} } }
child[concept color=poporange] { node[concept] {inconvenientes} child[level distance=2.6cm, sibling angle=60] { node[concept, font=\scriptsize] {adicción a las pantallas} } child[level distance=2.6cm, sibling angle=60] { node[concept, font=\scriptsize] {aislamiento} } }
child[concept color=poppurple] { node[concept] {soluciones} child[level distance=2.6cm, sibling angle=60] { node[concept, font=\scriptsize] {poner límites} } child[level distance=2.6cm, sibling angle=60] { node[concept, font=\scriptsize] {educar, no castigar} } }
child[concept color=popgold] { node[concept] {vocabulario clave} child[level distance=2.6cm, sibling angle=60] { node[concept, font=\scriptsize] {brecha digital} } child[level distance=2.6cm, sibling angle=60] { node[concept, font=\scriptsize] {prohibir, criar} } };
\end{tikzpicture}
\end{center}
\legende{Carte mentale : les quatre axes du débat « Educación y NTIC ».}
\end{popfigure}

\begin{aretenir}[L'essentiel de cette section]
\begin{itemize}
\item Le texte \esp{Los niños y las pantallas} met en scène un débat familial camerounais : la mère interdit, le père éduque, le narrateur nuance.
\item L'ordre s'exprime à l'\cle{impératif} (\esp{¡Apaga!}, \esp{¡Ven!}) ; l'interdiction impersonnelle avec \esp{se prohíbe + infinitivo}.
\item « Chez » se traduit \esp{en casa de} (état) ou \esp{a casa de} (mouvement).
\item « Devenir » : \esp{volverse}, \esp{hacerse}, \esp{ponerse}, \esp{convertirse en}, \esp{llegar a ser} selon le type de changement.
\item Le commentaire de texte suit toujours : introduction (situation + thème + problématique), développement (plan + citations traduites), conclusion (réponse + ouverture).
\end{itemize}
\end{aretenir}

\coursec{Grammaire 1 : l'ordre et l'interdiction}

\courssub{L'ordre : l'impératif affirmatif}

\begin{propriete}[Formation de l'impératif affirmatif (tú / vosotros)]
\begin{itemize}
\item \textbf{tú} : comme la 3\up{e} personne du singulier du présent de l'indicatif : \esp{hablar → habla}, \esp{comer → come}, \esp{escribir → escribe}.
\item \textbf{vosotros} : infinitif dont le \textbf{-r} final devient \textbf{-d} : \esp{hablar → hablad}, \esp{comer → comed}, \esp{escribir → escribid}.
\item \textbf{Impératifs irréguliers de tú} (à connaître par cœur) : \esp{decir → di}, \esp{hacer → haz}, \esp{ir → ve}, \esp{poner → pon}, \esp{salir → sal}, \esp{ser → sé}, \esp{tener → ten}, \esp{venir → ven}.
\end{itemize}
\end{propriete}

\begin{center}
\begin{tabular}{|l|l|l|}\hline
\rowcolor{popblueL} Infinitivo & tú & vosotros \\ \hline
apagar & apaga & apagad \\ \hline
venir (irrégulier) & ven & venid \\ \hline
poner (irrégulier) & pon & poned \\ \hline
hacer (irrégulier) & haz & haced \\ \hline
\end{tabular}
\end{center}

\esp{¡Haz los deberes antes de jugar!} « Fais tes devoirs avant de jouer ! » ; \esp{¡Venid a cenar, niños!} « Venez dîner, les enfants ! »

\courssub{L'interdiction : l'impératif négatif et les tournures impersonnelles}

\begin{propriete}[L'ordre négatif = no + subjonctif présent]
Pour défendre quelque chose à quelqu'un, l'espagnol n'utilise \emph{pas} l'impératif : il emploie \esp{no} + \textbf{subjonctif présent}.
\begin{itemize}
\item \esp{¡No fumes!} « Ne fume pas ! » (et non un impératif négatif).
\item \esp{¡No uses el teléfono en la mesa!} « N'utilise pas le téléphone à table ! »
\item \esp{¡No hagáis ruido!} « Ne faites pas de bruit ! » (vous, pluriel).
\end{itemize}
Rappel de formation du subjonctif présent : radical de la 1\up{re} personne du présent (\esp{hago} → \esp{hag-}) + terminaisons « croisées » : \textbf{-e} pour les verbes en \textbf{-ar}, \textbf{-a} pour les verbes en \textbf{-er/-ir}.
\end{propriete}

\begin{propriete}[L'interdiction impersonnelle]
Sur les panneaux et dans les règlements, l'interdiction s'exprime de façon impersonnelle :
\begin{itemize}
\item \esp{Prohibido + infinitivo} : \esp{Prohibido fumar} « Défense de fumer » ; \esp{Prohibido usar el móvil} « Portable interdit ».
\item \esp{Se prohíbe + infinitivo} : \esp{Se prohíbe la entrada} « Entrée interdite » ; \esp{Se prohíbe usar el móvil durante la comida} « Il est interdit d'utiliser le portable pendant le repas ».
\item \esp{Está prohibido + infinitivo} : \esp{Está prohibido correr en el pasillo} « Il est interdit de courir dans le couloir ».
\end{itemize}
\end{propriete}

\begin{attention}[Ne calque pas le français]
\begin{itemize}
\item « Ne fume pas » ne se dit jamais avec un impératif négatif : toujours \esp{¡No fumes!} (subjonctif).
\item Sur un panneau, pas de sujet : \esp{Prohibido fumar}, jamais une phrase complète avec « vous ».
\item N'oublie pas les points d'exclamation inversés : \esp{¡Apaga el teléfono!}
\end{itemize}
\end{attention}

\coursec{Grammaire 2 : traduire « chez »}

\begin{propriete}[Les équivalents espagnols de « chez »]
\begin{itemize}
\item \textbf{État (être, rester, manger chez...)} : \esp{en casa de + nom/pronom}. \esp{Ceno en casa de mi tío.} « Je dîne chez mon oncle. » \esp{En casa de ella no hay reglas.} « Chez elle, il n'y a pas de règles. »
\item \textbf{Mouvement (aller, venir chez...)} : \esp{a casa de + nom/pronom}. \esp{Vamos a casa de la tía María.} « Nous allons chez tante María. »
\item \textbf{Chez soi, sans précision} : \esp{en casa} / \esp{a casa}. \esp{Me quedo en casa.} « Je reste chez moi. » \esp{Vuelvo a casa.} « Je rentre à la maison. »
\item \textbf{« Chez » au sens de « parmi, dans la société de »} : \esp{entre}. \esp{Entre los españoles, la cena es tardía.} « Chez les Espagnols, le dîner est tardif. »
\item \textbf{Chez le commerçant, le professionnel} : \esp{en la + lieu} ou \esp{en el + lieu}. \esp{Compro el pan en la panadería.} « J'achète le pain chez le boulanger. »
\end{itemize}
\end{propriete}

\begin{center}
\begin{tabularx}{\linewidth}{|X|X|X|}\hline
\rowcolor{popblueL} Français & Español & Cas \\ \hline
Je suis chez ma sœur. & Estoy en casa de mi hermana. & état \\ \hline
Je vais chez ma sœur. & Voy a casa de mi hermana. & mouvement \\ \hline
Chez les Espagnols, on dîne tard. & Entre los españoles, se cena tarde. & « parmi » \\ \hline
Je reste chez moi. & Me quedo en casa. & sans complément \\ \hline
\end{tabularx}
\end{center}

\begin{attention}[Erreurs à éviter]
\begin{itemize}
\item Pas d'article après \esp{casa de} quand on parle du foyer d'une personne : \esp{en casa de Pedro}, et non \esp{en la casa de Pedro}.
\item « Chez les jeunes » au sens de « parmi les jeunes » = \esp{entre los jóvenes}, jamais \esp{en casa de los jóvenes} !
\end{itemize}
\end{attention}

\coursec{Grammaire 3 : traduire « devenir »}

\begin{propriete}[Cinq verbes pour cinq types de changement]
\begin{itemize}
\item \esp{ponerse + adjectif} : changement \textbf{passager}, physique ou d'humeur, souvent soudain. \esp{Se puso rojo de vergüenza.} « Il est devenu rouge de honte. » \esp{Se ponen insoportables.} « Ils deviennent insupportables. »
\item \esp{volverse + adjectif} : changement \textbf{profond, durable, souvent involontaire} (caractère, état mental). \esp{Se ha vuelto muy callado.} « Il est devenu très silencieux. » \esp{Se volvió loco.} « Il est devenu fou. »
\item \esp{hacerse + adjectif/nom} : changement obtenu \textbf{par l'effort ou la volonté} ; aussi pour les idées, la religion, la profession. \esp{Se ha hecho más responsable.} « Elle est devenue plus responsable. » \esp{Se hizo médico.} « Il est devenu médecin. »
\item \esp{convertirse en + nom} : \textbf{transformation} complète, souvent spectaculaire. \esp{Los niños se convierten en esclavos de las pantallas.} « Les enfants deviennent esclaves des écrans. »
\item \esp{llegar a ser + nom/adjectif} : \textbf{aboutissement} d'un long processus. \esp{Ha llegado a ser la primera de su clase.} « Elle est devenue la première de sa classe. »
\item \esp{quedarse} : « devenir » au sens de \textbf{rester dans un état} (souvent négatif et définitif). \esp{Se quedó ciego.} « Il est devenu aveugle. » \esp{Se quedó soltero.} « Il est resté célibataire. »
\end{itemize}
\end{propriete}

\begin{center}
\begin{tabularx}{\linewidth}{|X|X|X|}\hline
\rowcolor{popblueL} Verbe & Type de changement & Exemple \\ \hline
ponerse + adj. & passager, humeur & Se puso triste. « Il est devenu triste. » \\ \hline
volverse + adj. & durable, subi & Se ha vuelto perezoso. « Il est devenu paresseux. » \\ \hline
hacerse + adj./nom & effort, volonté & Se hizo rico. « Il est devenu riche. » \\ \hline
convertirse en + nom & transformation & Se convirtió en un problema. « C'est devenu un problème. » \\ \hline
llegar a ser & aboutissement & Llegó a ser director. « Il est devenu directeur. » \\ \hline
quedarse & état définitif & Se quedó sordo. « Il est devenu sourd. » \\ \hline
\end{tabularx}
\end{center}

\begin{attention}[Pièges de traduction]
\begin{itemize}
\item « Devenir » + adjectif d'humeur momentanée → \esp{ponerse} ; de caractère durable → \esp{volverse}. Compare : \esp{Se puso nervioso antes del examen} « Il est devenu nerveux avant l'examen » (passager) / \esp{Se ha vuelto muy nervioso desde el accidente} « Il est devenu très nerveux depuis l'accident » (durable).
\item \esp{Hacerse} est irrégulier au passé simple : \esp{hice, hiciste, hizo, hicimos, hicisteis, hicieron}.
\item Ne traduis jamais « devenir » par \esp{devenir} : ce verbe n'existe pas en espagnol !
\end{itemize}
\end{attention}

\begin{exoresolu}[Thème guidé : choisir le bon verbe]
Énoncé. Traduis en espagnol : a) « Elle est devenue rouge de colère. » b) « Mon frère est devenu professeur après de longues études. » c) « Il est devenu très bavard depuis qu'il a un téléphone. »

\tcblower

Solution détaillée. a) Changement passager d'humeur → \esp{ponerse} : \esp{Se puso roja de ira.} b) Aboutissement d'un long processus → \esp{llegar a ser} : \esp{Mi hermano llegó a ser profesor después de largos estudios.} c) Changement durable du caractère, subi → \esp{volverse} : \esp{Se ha vuelto muy hablador desde que tiene un teléfono.}
\end{exoresolu}

\coursec{Débat : la educación de los niños y las NTIC}

\courssub{Outils linguistiques pour débattre}

\begin{definition}[Expressions utiles pour le débat]
\begin{itemize}
\item Donner son avis : \esp{en mi opinión} « à mon avis » ; \esp{yo creo que / pienso que} « je crois que » ; \esp{me parece que} « il me semble que ».
\item Être d'accord : \esp{estoy de acuerdo (contigo)} « je suis d'accord (avec toi) » ; \esp{tienes razón} « tu as raison » ; \esp{es verdad que...} « il est vrai que... ».
\item Ne pas être d'accord : \esp{no estoy de acuerdo} « je ne suis pas d'accord » ; \esp{al contrario} « au contraire » ; \esp{sin embargo} « cependant » ; \esp{no creo que + subjonctif} « je ne crois pas que... ».
\item Nuancer : \esp{por un lado... por otro lado...} « d'un côté... de l'autre... » ; \esp{todo depende de...} « tout dépend de... ».
\item Conclure : \esp{en conclusión} « en conclusion » ; \esp{para terminar} « pour terminer ».
\end{itemize}
\end{definition}

\begin{attention}[Après « je ne crois pas que » : le subjonctif]
\esp{Creo que las pantallas son útiles.} « Je crois que les écrans sont utiles » (indicatif, opinion affirmative). Mais : \esp{No creo que las pantallas sean útiles.} « Je ne crois pas que les écrans soient utiles » (subjonctif après la négation d'un verbe d'opinion). C'est un point chaud du Bac !
\end{attention}

\begin{experience}[Débat en classe : ¿Prohibir o educar?]
\textbf{Organisation (20 min).} La classe se divise en deux équipes :
\begin{itemize}
\item \textbf{Équipe A} : \esp{Hay que prohibir las pantallas a los niños.} (Il faut interdire les écrans aux enfants.)
\item \textbf{Équipe B} : \esp{Es mejor educar que prohibir.} (Il vaut mieux éduquer qu'interdire.)
\end{itemize}
\textbf{Déroulement :} 1. Chaque équipe prépare trois arguments avec des exemples (5 min). 2. Débat dirigé par le professeur : chaque élève parle au moins une fois en utilisant une expression du lexique (10 min). 3. Vote final et conclusion commune (5 min).

\textbf{Contraintes grammaticales :} utilise au moins une fois \esp{se prohíbe}, un impératif (\esp{¡Piensa en los niños!}), un verbe de « devenir » (\esp{los niños se vuelven adictos}) et \esp{no creo que + subjonctif}.
\end{experience}

\begin{exemplebox}[Modèle de prise de parole]
\esp{En mi opinión, no hay que prohibir las pantallas por completo. Es verdad que algunos niños se vuelven adictos al móvil, pero no creo que la prohibición sea la solución. Por un lado, las nuevas tecnologías ayudan a estudiar; por otro lado, los padres deben poner límites claros. En conclusión, es mejor educar que castigar.}

« À mon avis, il ne faut pas interdire complètement les écrans. Il est vrai que certains enfants deviennent accros au portable, mais je ne crois pas que l'interdiction soit la solution. D'un côté, les nouvelles technologies aident à étudier ; de l'autre, les parents doivent fixer des limites claires. En conclusion, il vaut mieux éduquer que punir. »
\end{exemplebox}

\begin{aretenir}[Bilan de la leçon]
\begin{itemize}
\item \textbf{Ordre} : impératif affirmatif (\esp{¡Apaga!}, \esp{¡Ven!}, \esp{¡Haz!}) ; \textbf{ordre négatif} : \esp{no + subjonctif} (\esp{¡No fumes!}).
\item \textbf{Interdiction impersonnelle} : \esp{prohibido + infinitivo}, \esp{se prohíbe + infinitivo}, \esp{está prohibido + infinitivo}.
\item \textbf{Chez} : \esp{en casa de} (état), \esp{a casa de} (mouvement), \esp{entre} (« chez les Espagnols » = \esp{entre los españoles}).
\item \textbf{Devenir} : \esp{ponerse} (passager), \esp{volverse} (durable, subi), \esp{hacerse} (effort), \esp{convertirse en} (transformation), \esp{llegar a ser} (aboutissement), \esp{quedarse} (état définitif).
\item \textbf{Débat} : \esp{no creo que + subjonctif} ; articuler avec \esp{por un lado... por otro lado...}.
\end{itemize}
\end{aretenir}

\begin{exercice}[À toi de jouer]
\begin{enumerate}
\item Traduis en espagnol : « Éteins la tablette et viens dîner ! » ; « Ne touche pas à cet appareil ! » ; « Il est interdit de fumer dans le lycée » ; « Nous allons chez ma tante le week-end » ; « Mon frère est devenu très paresseux » ; « Chez les jeunes camerounais, le téléphone portable est partout. »
\item Remplace « devenir » par le bon verbe espagnol : a) « Elle est devenue médecin par son travail. » b) « Il est devenu pâle en entendant la nouvelle. » c) « L'eau devient glace à zéro degré. »
\item Rédige une introduction de commentaire (5 lignes) pour le texte \esp{Los niños y las pantallas} en suivant la méthode.
\item Prépare trois arguments oraux pour le débat : \esp{¿Deben los padres prohibir el móvil a sus hijos?}
\end{enumerate}
\end{exercice}

\begin{corrige}[Exercice « À toi de jouer » (1 et 2)]
\textbf{1.} \esp{¡Apaga la tableta y ven a cenar!} ; \esp{¡No toques ese aparato!} ; \esp{Se prohíbe fumar en el liceo} (ou \esp{Prohibido fumar en el liceo}) ; \esp{Vamos a casa de mi tía el fin de semana} ; \esp{Mi hermano se ha vuelto muy perezoso} ; \esp{Entre los jóvenes cameruneses, el teléfono móvil está en todas partes.}

\textbf{2.} a) \esp{Se hizo médica gracias a su trabajo} (effort) ; b) \esp{Se puso pálido al oír la noticia} (réaction passagère) ; c) \esp{El agua se convierte en hielo a cero grados} (transformation).
\end{corrige}

\coursec{Lexique : éducation des enfants et NTIC}

Dans la section précédente, nous avons lu et commenté le texte \esp{« Los niños y las pantallas »}, qui décrivait la vie quotidienne d'une famille partagée entre l'école, les devoirs et la fascination des écrans. Ce texte contenait de nombreux mots nouveaux : \esp{criar}, \esp{los deberes}, \esp{estar enganchado}, \esp{el control parental}... Avant d'attaquer la grammaire (l'ordre, l'interdiction, « chez », « devenir »), il faut d'abord \cle{maîtriser ce vocabulaire} : c'est lui qui nous permettra de traduire correctement au thème comme à la version, et de participer au débat final sur l'éducation des enfants et les nouvelles technologies. Organisons donc ce lexique en quatre champs : la famille et l'éducation, les NTIC, l'argumentation, et enfin les effets et les solutions.

\begin{definition}[Qu'est-ce que les NTIC ?]
\cle{NTIC} est l'abréviation de \esp{Nuevas Tecnologías de la Información y la Comunicación}, c'est-à-dire « nouvelles technologies de l'information et de la communication ». En espagnol, on dit aussi simplement \esp{las nuevas tecnologías} ou \esp{las tecnologías}. Ce terme regroupe l'ordinateur, le téléphone portable, la tablette, Internet et les réseaux sociaux. Attention au genre : on dit \esp{las NTIC} (féminin pluriel), car le mot \esp{tecnologías} est féminin.
\end{definition}

\courssub{Champ 1 : L'éducation et la famille}

Observons d'abord quelques phrases tirées de notre texte support et de la vie quotidienne :

\esp{Los padres educan a sus hijos con paciencia.} « Les parents éduquent leurs enfants avec patience. »

\esp{María fue criada por sus abuelos en el pueblo.} « María a été élevée par ses grands-parents au village. »

\esp{El niño hace sus deberes antes de jugar.} « L'enfant fait ses devoirs avant de jouer. »

\esp{No hay que castigar a los niños sin explicarles el motivo.} « Il ne faut pas punir les enfants sans leur expliquer la raison. »

On remarque une distinction essentielle que le français ne fait pas toujours clairement : \esp{educar} (éduquer, former l'esprit et le comportement) et \esp{criar} (élever, nourrir et faire grandir). Le nom correspondant à \esp{criar} est \esp{la crianza}, qui désigne concrètement l'éducation des enfants au quotidien. Notez aussi la \cle{a} personnelle devant la personne complément d'objet direct : \esp{educar a un niño}, \esp{castigar a los niños}.

\begin{center}
\begin{tabularx}{\linewidth}{|X|X|X|}
\hline
\rowcolor{popblueL} Español & Français & Remarque \\
\hline
educar (a un niño) & éduquer (un enfant) & verbe régulier ; \esp{la educación} = l'éducation, l'instruction \\
\hline
criar (a los hijos) & élever (les enfants) & participe passé : \esp{criado} \\
\hline
la crianza & l'éducation des enfants, l'élevage & nom féminin \\
\hline
castigar & punir & régulier ; \esp{el castigo} = la punition \\
\hline
el castigo (los castigos) & la punition & nom masculin \\
\hline
consentir (e-ie) & gâter, trop tolérer & verbe irrégulier : \esp{consiento, consientes, consiente, consentimos, consentís, consienten} \\
\hline
los deberes & les devoirs (scolaires) & presque toujours au pluriel \\
\hline
la disciplina & la discipline & nom féminin \\
\hline
\end{tabularx}
\end{center}

\begin{attention}[Faux amis : educar / éduquer]
Le verbe \esp{educar} est plus large que le français « éduquer » : \esp{una persona educada} signifie « une personne polie, bien élevée », et \esp{¡qué maleducado!} veut dire « quel malpoli ! ». De même, \esp{los deberes} ne signifie pas « les devoirs » au sens moral (ce serait \esp{las obligaciones}) mais bien les devoirs scolaires : \esp{hacer los deberes} = « faire ses devoirs ».
\end{attention}

\courssub{Champ 2 : Les NTIC}

Voici le vocabulaire indispensable pour parler des écrans et d'Internet. Observez les articles et les pluriels :

\begin{center}
\begin{tabularx}{\linewidth}{|X|X|X|}
\hline
\rowcolor{popblueL} Español & Français & Remarque \\
\hline
la pantalla (las pantallas) & l'écran & \esp{el tiempo de pantalla} = le temps d'écran \\
\hline
el móvil / el celular & le portable & \esp{móvil} en Espagne, \esp{celular} en Amérique latine \\
\hline
la tableta & la tablette & nom féminin \\
\hline
las redes sociales & les réseaux sociaux & \esp{una red social} au singulier \\
\hline
navegar por Internet & naviguer sur Internet & préposition \esp{por}, sans article \\
\hline
descargar & télécharger (depuis Internet) & régulier ; \esp{descargar una aplicación} \\
\hline
el videojuego (los videojuegos) & le jeu vidéo & un seul mot en espagnol \\
\hline
estar enganchado a & être accro à & suivi de la préposition \esp{a} \\
\hline
\end{tabularx}
\end{center}

\begin{attention}[estar enganchado a]
\esp{Estar enganchado a} signifie « être accro à », au sens familier d'« être dépendant » : \esp{Mi hermano está enganchado a los videojuegos.} « Mon frère est accro aux jeux vidéo. » Ne traduisez jamais littéralement \esp{enganchado} par « accroché » : \esp{estar enganchado a} est une expression figée. Pour un registre plus soutenu, on peut dire \esp{ser adicto a} : \esp{Los adolescentes son adictos al móvil.} « Les adolescents sont accros au portable. »
\end{attention}

\begin{savaistu}[ordenador ou computadora ?]
En Espagne, on dit \esp{el ordenador} ; en Amérique latine, on dit \esp{la computadora} (ou \esp{el computador}). Les deux formes sont correctes et acceptées au baccalauréat. De même, « le portable » se dit \esp{el móvil} en Espagne et \esp{el celular} en Amérique latine. En Guinée équatoriale, pays hispanophone voisin du Cameroun, on entend surtout \esp{el móvil}, sous l'influence de l'Espagne.
\end{savaistu}

\courssub{Champ 3 : Le débat et l'argumentation}

Pour le débat de la section 6, il vous faudra donner votre avis, marquer votre accord ou votre désaccord, et structurer votre argumentation. Voici les outils de base :

\begin{center}
\begin{tabularx}{\linewidth}{|X|X|X|}
\hline
\rowcolor{popblueL} Español & Français & Emploi \\
\hline
a favor de / en contra de & pour / contre & \esp{Estoy a favor del control parental.} \\
\hline
en mi opinión & à mon avis & en début de phrase, suivi d'une virgule \\
\hline
sin embargo & cependant & oppose deux idées \\
\hline
además & de plus, en outre & ajoute un argument \\
\hline
por un lado... por otro lado & d'un côté... de l'autre & présente les deux faces d'un problème \\
\hline
estoy de acuerdo / no estoy de acuerdo & je suis d'accord / je ne suis pas d'accord & \esp{estar de acuerdo con alguien} = être d'accord avec quelqu'un \\
\hline
\end{tabularx}
\end{center}

\begin{exemplebox}[Exemples traduits]
\esp{Los padres deben limitar el tiempo de pantalla.} « Les parents doivent limiter le temps d'écran. »

\esp{Estoy a favor del control parental.} « Je suis pour le contrôle parental. »

\esp{Por un lado, las NTIC ayudan a aprender; por otro lado, pueden crear adicción.} « D'un côté, les NTIC aident à apprendre ; de l'autre, elles peuvent créer une addiction. »

\esp{En mi opinión, los niños no deben usar el móvil en la mesa. Sin embargo, muchos padres no están de acuerdo.} « À mon avis, les enfants ne doivent pas utiliser le portable à table. Cependant, beaucoup de parents ne sont pas d'accord. »
\end{exemplebox}

\courssub{Champ 4 : Effets et solutions}

\begin{center}
\begin{tabularx}{\linewidth}{|X|X|X|}
\hline
\rowcolor{popblueL} Español & Français & Remarque \\
\hline
la adicción (las adicciones) & l'addiction & \esp{la adicción a las pantallas} \\
\hline
el control parental & le contrôle parental & masculin \\
\hline
limitar el tiempo & limiter le temps & \esp{limitar el tiempo de pantalla} \\
\hline
la brecha digital & la fracture numérique & écart entre ceux qui ont accès aux NTIC et les autres \\
\hline
aprovechar & tirer profit de, profiter de & \esp{aprovechar las NTIC para estudiar} \\
\hline
\end{tabularx}
\end{center}

\esp{Hay que aprovechar las ventajas de Internet sin caer en la adicción.} « Il faut tirer profit des avantages d'Internet sans tomber dans l'addiction. »

\esp{En muchos barrios de Yaoundé, la brecha digital sigue siendo un problema.} « Dans beaucoup de quartiers de Yaoundé, la fracture numérique reste un problème. »

\courssub{Carte mentale du lexique}

\begin{popfigure}
\begin{center}
\begin{tikzpicture}
\node[draw=popblue, fill=popblueL, rounded corners, font=\bfseries\small, align=center] (fam) at (0,0) {La familia};
\node[draw=poporange, fill=poporangeL, rounded corners, font=\bfseries\small, align=center] (ntic) at (8,0) {Las NTIC};
\node[draw=popline, rounded corners, align=center, fill=white, font=\small] at (-2.6,-1.6) {educar / criar\\la crianza};
\node[draw=popline, rounded corners, align=center, fill=white, font=\small] at (0,-2.2) {castigar / el castigo\\consentir};
\node[draw=popline, rounded corners, align=center, fill=white, font=\small] at (2.6,-1.6) {los deberes\\la disciplina};
\node[draw=popline, rounded corners, align=center, fill=white, font=\small] at (5.4,-1.6) {la pantalla\\el móvil / la tableta};
\node[draw=popline, rounded corners, align=center, fill=white, font=\small] at (8,-2.2) {las redes sociales\\navegar por Internet};
\node[draw=popline, rounded corners, align=center, fill=white, font=\small] at (10.6,-1.6) {descargar\\el videojuego};
\draw[popblue, thick] (fam) -- (-2.6,-1.2);
\draw[popblue, thick] (fam) -- (0,-1.8);
\draw[popblue, thick] (fam) -- (2.6,-1.2);
\draw[poporange, thick] (ntic) -- (5.4,-1.2);
\draw[poporange, thick] (ntic) -- (8,-1.8);
\draw[poporange, thick] (ntic) -- (10.6,-1.2);
\end{tikzpicture}
\end{center}
\legende{Les deux grands pôles du lexique de la leçon : la famille et les NTIC.}
\end{popfigure}

\courssub{Entraînement rapide}

\begin{exoresolu}[Jeu de correspondance mot / définition]
Énoncé. Reliez chaque mot espagnol à sa définition française :

1. \esp{la crianza} \quad 2. \esp{la brecha digital} \quad 3. \esp{consentir} \quad 4. \esp{descargar} \quad 5. \esp{estar enganchado a}

a. télécharger un fichier depuis Internet \quad b. l'éducation des enfants au quotidien \quad c. être accro à quelque chose \quad d. gâter un enfant en lui tout permettant \quad e. l'écart entre ceux qui ont accès aux NTIC et ceux qui n'y ont pas accès
\tcblower
Solution détaillée :

1 -- b : \esp{la crianza} vient du verbe \esp{criar} (élever) et désigne l'éducation quotidienne des enfants.

2 -- e : \esp{la brecha digital} est la « fracture numérique » ; \esp{brecha} signifie « brèche, écart ».

3 -- d : \esp{consentir} un enfant, c'est le gâter, tout lui permettre ; attention, c'est un verbe à changement de voyelle (\esp{consiento}).

4 -- a : \esp{descargar} = télécharger (par exemple une application : \esp{descargar una aplicación}).

5 -- c : \esp{estar enganchado a} = être accro à ; ne pas traduire par « être accroché ».
\end{exoresolu}

\begin{exercice}[Phrases à compléter]
Complétez les phrases avec le mot qui convient : \esp{deberes} -- \esp{pantalla} -- \esp{castigo} -- \esp{control parental} -- \esp{redes sociales} -- \esp{aprovechar}.

1. Antes de salir a jugar, los niños deben hacer sus \ldots

2. Muchos adolescentes pasan horas en las \ldots como Instagram o TikTok.

3. Los padres deben limitar el tiempo de \ldots de sus hijos.

4. El profesor impuso un \ldots al alumno que llegó tarde.

5. Hay que \ldots las NTIC para aprender español.

6. Estoy a favor del \ldots en los móviles de los niños.
\end{exercice}

\begin{corrige}[Exercice « Phrases à compléter »]
1. \esp{deberes} : « faire ses devoirs » = \esp{hacer los deberes}.

2. \esp{redes sociales} : toujours au pluriel dans ce contexte.

3. \esp{pantalla} : \esp{el tiempo de pantalla} = le temps d'écran.

4. \esp{castigo} : \esp{imponer un castigo} = infliger une punition.

5. \esp{aprovechar} : tirer profit des NTIC.

6. \esp{control parental} : expression composée, \esp{control} est masculin.
\end{corrige}

\begin{aretenir}[L'essentiel du lexique]
\begin{itemize}
\item \esp{educar} (éduquer) et \esp{criar} (élever) ne sont pas synonymes ; \esp{la crianza} = l'éducation des enfants.
\item \esp{los deberes} = les devoirs scolaires ; \esp{estar enganchado a} = être accro à.
\item Pour débattre : \esp{a favor de / en contra de}, \esp{en mi opinión}, \esp{sin embargo}, \esp{por un lado... por otro lado}.
\item \esp{ordenador} (Espagne) et \esp{computadora} (Amérique latine) sont tous deux acceptés.
\end{itemize}
\end{aretenir}

Ce vocabulaire est maintenant votre boîte à outils. Dans la section suivante, nous verrons comment ces mots se combinent avec l'\cle{impératif} et le \cle{subjonctif} pour exprimer l'ordre et l'interdiction : \esp{¡Haz los deberes!}, \esp{No fumes}, \esp{Se prohíbe usar el móvil en clase}...

\coursec{Grammaire 1 : L'ordre et l'interdiction}

Dans la section précédente, nous avons découvert le lexique de l'éducation des enfants et des NTIC : \esp{los padres}, \esp{las pantallas}, \esp{el teléfono móvil}, \esp{la adicción}, \esp{controlar}, \esp{limiter}. Or, dès qu'on parle d'éducation, on parle d'ordres et d'interdictions : un parent dit à son enfant « Éteins ton téléphone ! », « Ne passe pas toute la soirée devant l'écran ! » ; une école affiche « Défense de fumer ». Comment l'espagnol exprime-t-il l'ordre et l'interdiction ? C'est l'objet de cette leçon de grammaire, essentielle pour le thème comme pour la version.

\courssub{Observation : partons du texte}

Reprenons quelques phrases tirées de notre texte support \emph{Los niños y las pantallas} et de la vie quotidienne. Observons-les avant d'énoncer la moindre règle.

\begin{exemplebox}[Phrases à observer]
\esp{¡Apaga el teléfono antes de cenar!} « Éteins le téléphone avant le dîner ! »

\esp{¡No fumes delante de los niños!} « Ne fume pas devant les enfants ! »

\esp{Prohibido fumar.} « Défense de fumer. »

\esp{Se prohíbe el uso del móvil en clase.} « L'usage du portable est interdit en classe. »

\esp{No pisar el césped.} « Ne pas marcher sur la pelouse. »
\end{exemplebox}

Que remarquons-nous ? Pour un ordre \textbf{affirmatif}, l'espagnol utilise une forme spéciale du verbe : \esp{apaga} (et non \esp{apagas}). Pour un ordre \textbf{négatif}, on trouve \esp{no} suivi d'une forme qui ressemble au subjonctif : \esp{no fumes}. Pour une interdiction \textbf{générale et impersonnelle}, on trouve soit \esp{prohibido} suivi de l'infinitif, soit \esp{se prohíbe} suivi de l'infinitif, soit l'infinitif seul. Trois constructions, trois situations : c'est toute la logique du système.

\courssub{L'ordre affirmatif : l'impératif}

\begin{definition}[L'impératif]
L'\cle{impératif} est le mode qui sert à donner un ordre, un conseil ou une instruction à une personne présente. En espagnol, il possède des formes propres pour \esp{tú}, \esp{vosotros}, \esp{usted} et \esp{ustedes} ; il n'existe pas à la première personne du singulier (on ne s'ordonne rien à soi-même avec l'impératif).
\end{definition}

\begin{propriete}[Formation de l'impératif affirmatif (tú et vosotros)]
Pour \esp{tú} : on prend la 2\textsuperscript{e} personne du singulier du présent de l'indicatif \textbf{sans le -s final} : \esp{hablas → habla}, \esp{comes → come}, \esp{escribes → escribe}.

Pour \esp{vosotros} : on remplace le \textbf{-r} de l'infinitif par \textbf{-d} : \esp{hablar → hablad}, \esp{comer → comed}, \esp{escribir → escribid}.

Pour \esp{usted / ustedes} : on utilise le subjonctif présent : \esp{hable}, \esp{hablen}.
\end{propriete}

\begin{center}
\begin{tabular}{|l|l|l|l|}
\hline
\rowcolor{popblueL} Infinitif & tú & vosotros & usted / ustedes \\
\hline
hablar (parler) & habla & hablad & hable / hablen \\
\hline
comer (manger) & come & comed & coma / coman \\
\hline
escribir (écrire) & escribe & escribid & escriba / escriban \\
\hline
\end{tabular}
\end{center}

\begin{attention}[Les huit impératifs irréguliers de « tú »]
Huit verbes très fréquents ont un impératif de \esp{tú} irrégulier, à apprendre par cœur : \esp{decir → di} (dis), \esp{hacer → haz} (fais), \esp{ir → ve} (va), \esp{poner → pon} (mets), \esp{salir → sal} (sors), \esp{ser → sé} (sois), \esp{tener → ten} (tiens / aie), \esp{venir → ven} (viens). Exemple : \esp{¡Ven aquí y pon la mesa!} « Viens ici et mets la table ! »
\end{attention}

\courssub{L'ordre négatif : no + subjonctif présent}

C'est LA grande différence avec le français, et la source d'innombrables erreurs au baccalauréat.

\begin{propriete}[Règle de l'ordre négatif]
En espagnol, \textbf{l'impératif négatif n'existe pas}. Tout ordre négatif se construit avec \esp{no} + \textbf{subjonctif présent} : \esp{No fumes} « Ne fume pas » ; \esp{No hables} « Ne parle pas » ; \esp{No salgas} « Ne sors pas » ; \esp{No vengáis} « Ne venez pas » (vosotros).
\end{propriete}

Rappel de formation du subjonctif présent : radical de la 1\textsuperscript{re} personne du singulier de l'indicatif + terminaisons « croisées » : \textbf{-e/-es/-e/-emos/-éis/-en} pour les verbes en \textbf{-ar}, \textbf{-a/-as/-a/-amos/-áis/-an} pour les verbes en \textbf{-er/-ir}. Ainsi : \esp{hablar → no hables} ; \esp{comer → no comas} ; \esp{salir (salgo) → no salgas} ; \esp{hacer (hago) → no hagas} ; \esp{ir → no vayas}.

\begin{center}
\begin{tabular}{|l|l|l|}
\hline
\rowcolor{popblueL} Infinitif & Ordre affirmatif (impératif) & Ordre négatif (no + subjonctif) \\
\hline
hablar & ¡Habla! « Parle ! » & ¡No hables! « Ne parle pas ! » \\
\hline
comer & ¡Come! « Mange ! » & ¡No comas! « Ne mange pas ! » \\
\hline
ir & ¡Ve! « Va ! » & ¡No vayas! « N'y va pas ! » \\
\hline
hacer & ¡Haz! « Fais ! » & ¡No hagas! « Ne fais pas ! » \\
\hline
\end{tabular}
\end{center}

\begin{attention}[Erreur fréquente des francophones]
Ne traduisez JAMAIS « ne fume pas » par \esp{no fumas} : c'est de l'\textbf{indicatif}, qui constate un fait (« tu ne fumes pas »), et qui ne peut pas exprimer l'ordre. L'ordre négatif exige le subjonctif : \esp{no fumes}. De même : « ne sors pas » = \esp{no salgas} (et non \esp{no sales}).
\end{attention}

\courssub{L'ordre impersonnel : l'infinitif}

Sur les notices, les recettes de cuisine, les modes d'emploi et les panneaux, quand l'ordre ne s'adresse à personne en particulier, l'espagnol — comme le français — utilise l'\cle{infinitif}.

\begin{exemplebox}[L'infinitif d'instruction]
\esp{No pisar el césped.} « Ne pas marcher sur la pelouse. »

\esp{Agitar antes de usar.} « Agiter avant l'emploi. »

\esp{Añadir la sal y mezclar bien.} « Ajouter le sel et bien mélanger. » (recette)

\esp{No tocar.} « Ne pas toucher. » (musée)
\end{exemplebox}

\courssub{L'interdiction officielle : prohibido + infinitif et se prohíbe + infinitif}

\begin{propriete}[Règle de l'interdiction impersonnelle]
Pour exprimer une interdiction \textbf{générale et officielle} (panneaux, règlements), l'espagnol utilise :
\begin{itemize}
\item \esp{prohibido} + infinitif : \esp{Prohibido fumar} « Défense de fumer » ;
\item \esp{se prohíbe} + infinitif (ou + nom) : \esp{Se prohíbe entrar} « Il est interdit d'entrer ».
\end{itemize}
Sur les panneaux, \esp{prohibido} reste \textbf{invariable} : on écrit \esp{Prohibido fumar}, jamais \esp{prohibida fumar}. Attention : \esp{se prohíbe} est une 3\textsuperscript{e} personne du singulier à valeur impersonnelle (avec l'accent sur le \emph{í}).
\end{propriete}

\begin{exemplebox}[Exemples traduits]
\esp{Prohibido aparcar.} « Stationnement interdit. »

\esp{Prohibido el paso.} « Sens interdit / Passage interdit. »

\esp{Se prohíbe la entrada a los menores.} « L'entrée est interdite aux mineurs. »

\esp{Se prohíbe hacer ruido después de las diez.} « Il est interdit de faire du bruit après dix heures. »
\end{exemplebox}

\courssub{L'interdiction personnelle : prohibir}

Quand l'interdiction vise une personne précise, on utilise le verbe \esp{prohibir} (interdire), avec deux constructions possibles.

\begin{propriete}[Constructions de « prohibir »]
\begin{enumerate}
\item \esp{prohibir} + que + \textbf{subjonctif} : \esp{Te prohíbo que salgas.} « Je t'interdis de sortir. »
\item \esp{prohibir} + complément + \textbf{infinitif} : \esp{Te prohíbo salir.} « Je t'interdis de sortir. »
\end{enumerate}
Les deux constructions sont correctes ; la première (avec \esp{que} + subjonctif) est la plus fréquente et la plus sûre à l'examen.
\end{propriete}

\begin{exemplebox}[Exemples]
\esp{Los padres prohíben a sus hijos que usen el móvil por la noche.} « Les parents interdisent à leurs enfants d'utiliser le portable le soir. »

\esp{El director nos prohibió llegar tarde.} « Le directeur nous a interdit d'arriver en retard. »
\end{exemplebox}

\courssub{La place des pronoms avec l'impératif}

\begin{propriete}[Pronoms et impératif]
\begin{itemize}
\item Ordre \textbf{affirmatif} : les pronoms se placent \textbf{après} le verbe et \textbf{soudés} à lui (avec accent graphique pour conserver la prononciation) : \esp{Dímelo.} « Dis-le-moi. » \esp{Levántate.} « Lève-toi. » \esp{Dámelo.} « Donne-le-moi. »
\item Ordre \textbf{négatif} : les pronoms se placent \textbf{avant} le verbe, détachés : \esp{No me lo digas.} « Ne me le dis pas. » \esp{No te levantes.} « Ne te lève pas. »
\end{itemize}
\end{propriete}

\courssub{Comparaison français / espagnol}

\begin{center}
\begin{tabularx}{\linewidth}{|X|X|X|}
\hline
\rowcolor{popblueL} Français & Español & Remarque \\
\hline
Fume ! / Ne fume pas ! & ¡Fuma! / ¡No fumes! & Négatif : no + subjonctif \\
\hline
Défense de fumer & Prohibido fumar & « Prohibido » invariable \\
\hline
Il est interdit d'entrer & Se prohíbe entrar & Tournure impersonnelle \\
\hline
Ne pas marcher sur la pelouse & No pisar el césped & Infinitif dans les deux langues \\
\hline
Je t'interdis de sortir & Te prohíbo que salgas & Subjonctif après « que » \\
\hline
Dis-le-moi / Ne me le dis pas & Dímelo / No me lo digas & Place des pronoms inversée \\
\hline
\end{tabularx}
\end{center}

\begin{methode}[Traduire un ordre ou une interdiction (Thème / Version)]
\begin{enumerate}
\item \textbf{Identifier le destinataire} : personne précise (\esp{tú}, \esp{vosotros}, \esp{usted}) ou public indéterminé (panneau, notice) ?
\item \textbf{Ordre affirmatif à quelqu'un} $\rightarrow$ \textbf{Impératif} (\esp{tú} : forme spéciale ; \esp{vosotros} : -d ; \esp{usted} : subjonctif). Vérifier les 8 irréguliers (\esp{di, haz, ve, pon, sal, sé, ten, ven}).
\item \textbf{Ordre négatif à quelqu'un} $\rightarrow$ \textbf{No + Subjonctif présent}. Jamais d'indicatif ! Attention aux changements de radical (\esp{salir → salga}, \esp{tener → tenga}).
\item \textbf{Interdiction générale (panneau, règlement)} $\rightarrow$ \textbf{Prohibido + Infinitif} (invariable) \textbf{OU} \textbf{Se prohíbe + Infinitif/Nom} \textbf{OU} \textbf{Infinitif seul} (\esp{No fumar}).
\item \textbf{Interdiction personnelle (sujet explicite)} $\rightarrow$ \textbf{Prohibir + que + Subjonctif} (préférable) ou \textbf{Prohibir + Infinitif}.
\item \textbf{Pronoms} : Affirmatif = soudés \textbf{après} (accent) ; Négatif = détachés \textbf{avant}.
\end{enumerate}
\end{methode}

\begin{popfigure}
\begin{tikzpicture}[
node distance=1.1cm and 1.5cm,
every node/.style={font=\small, align=center},
start/.style={rectangle, rounded corners, draw=popdark, fill=popblueL, text width=4.5cm},
quest/.style={diamond, draw=popdark, fill=popgoldL, text width=3.2cm, aspect=2, inner sep=2pt},
sol/.style={rectangle, rounded corners, draw=popdark, fill=popgreenL, text width=4.5cm},
fleche/.style={-{Stealth[length=2.5mm]}, thick, popdark}
]
\node[start] (dep) {Comment exprimer l'ordre / l'interdiction ?};
\node[quest, below=of dep] (q1) {S'adresse-t-on à une personne précise ?};
\node[sol, below left=of q1, xshift=-1.5cm, align=center] (impers){Interdiction générale / Panneau :\\ \esp{Prohibido + Inf.}\\ \esp{Se prohíbe + Inf.}\\ \esp{No + Inf.} (\esp{No fumar})};
\node[quest, below right=of q1, xshift=1.5cm] (q2) {Ordre affirmatif ?};
\node[sol, below left=of q2, align=center] (aff){Impératif affirmatif :\\ \esp{¡Habla! ¡Come! ¡Escribe!}\\ Irréguliers \esp{tú} : \esp{di, haz, ve, pon, sal, sé, ten, ven}};
\node[sol, below right=of q2, align=center] (neg){Ordre négatif :\\ \esp{No} + \textbf{Subjonctif présent}\\ \esp{¡No hables! ¡No comas! ¡No salgas!}};
\draw[fleche] (dep) -- (q1);
\draw[fleche] (q1) -| node[near start, above, sloped, font=\tiny]{Non (panneau, notice)} (impers);
\draw[fleche] (q1) -| node[near start, above, sloped, font=\tiny]{Oui} (q2);
\draw[fleche] (q2) -| node[near start, above, font=\tiny]{Oui} (aff);
\draw[fleche] (q2) -| node[near start, above, font=\tiny]{Non} (neg);
\end{tikzpicture}
\legende{Organigramme de choix : ordre personnel affirmatif, ordre négatif ou interdiction générale.}
\end{popfigure}

\begin{savaistu}[Les panneaux en Espagne]
Sur les panneaux espagnols, on lit souvent des formules variées : \esp{Se ruega no fumar} « On prie de ne pas fumer » (formule polie, littéralement « on supplie »), \esp{No fumar} (infinitif sec), ou encore \esp{Prohibido fumar}. Dans les écoles, on trouve aussi \esp{Se prohíbe el uso de teléfonos móviles}. Le verbe \esp{rogar} (prier, supplier) est un verbe à diphtongaison : \esp{ruego, ruegas, ruega...}
\end{savaistu}

\begin{exoresolu}[Thème : ordres et interdictions]
Traduisez en espagnol : 1. Éteins ton portable ! 2. Ne touche pas l'ordinateur de ton frère ! 3. Défense d'entrer. 4. Il est interdit d'utiliser le téléphone en classe. 5. Ne me le répète pas !
\tcblower
\textbf{Solution détaillée.}

1. Ordre affirmatif à « tu » : impératif de \esp{apagar} $\rightarrow$ \esp{¡Apaga tu móvil!} (régulier : \esp{apagas} sans le -s).

2. Ordre négatif : \esp{no} + subjonctif de \esp{tocar} $\rightarrow$ \esp{¡No toques el ordenador de tu hermano!} (attention : \esp{tocar → toques}, changement orthographique c $\rightarrow$ qu devant e).

3. Interdiction générale : \esp{Prohibido entrar.} ou \esp{Prohibida la entrada.} (accord de \esp{prohibido} avec le nom \esp{entrada}).

4. Tournure impersonnelle : \esp{Se prohíbe usar el teléfono en clase.}

5. Ordre négatif + pronoms avant le verbe : \esp{¡No me lo repitas!} (subjonctif de \esp{repetir} : \esp{repitas}, avec changement e $\rightarrow$ i).
\end{exoresolu}

\begin{exercice}[À vous de jouer]
Traduisez en espagnol :
\begin{enumerate}
\item Viens ici tout de suite !
\item Ne sortez pas (vosotros) pendant la récréation !
\item Stationnement interdit.
\item Je vous interdis de regarder la télévision avant de faire vos devoirs.
\item Donne-le-moi ! / Ne me le donne pas !
\end{enumerate}
\end{exercice}

\begin{corrige}[Exercice : ordres et interdictions]
1. \esp{¡Ven aquí ahora mismo!} (impératif irrégulier de \esp{venir}). 
2. \esp{¡No salgáis durante el recreo!} (no + subjonctif, 2\textsuperscript{e} pers. plur.). 
3. \esp{Prohibido aparcar.} 
4. \esp{Os prohíbo que veáis la televisión antes de hacer los deberes.} (\esp{prohibir + que} + subjonctif). 
5. \esp{¡Dámelo!} (pronoms soudés après) / \esp{¡No me lo des!} (pronoms avant + subjonctif de \esp{dar} : \esp{des}).
\end{corrige}

\begin{aretenir}[L'essentiel]
Ordre affirmatif $\rightarrow$ \textbf{impératif} (\esp{habla, come, escribe} ; irréguliers : \esp{di, haz, ve, pon, sal, sé, ten, ven}). Ordre négatif $\rightarrow$ \textbf{no + subjonctif présent} (\esp{no fumes}), jamais l'indicatif. Interdiction générale $\rightarrow$ \esp{prohibido + infinitif} ou \esp{se prohíbe + infinitif}. Interdiction personnelle $\rightarrow$ \esp{prohibir + que + subjonctif}. Pronoms : soudés après l'ordre affirmatif (\esp{dímelo}), placés avant l'ordre négatif (\esp{no me lo digas}).
\end{aretenir}

\coursec{Grammaire 2 : traduire « chez »}

Dans la section précédente, nous avons appris à exprimer l'ordre et l'interdiction (\esp{No fumes}, \esp{Se prohíbe fumar}, \esp{Prohibido fumar}). Nous abordons maintenant un petit mot français qui pose de gros problèmes aux hispanisants : \cle{chez}. En effet, contrairement au français qui utilise un seul mot dans tous les contextes, l'espagnol impose un \textbf{choix} selon le sens : lieu où l'on est, mouvement, profession, mœurs d'un peuple, œuvre d'un auteur. C'est une question classique de thème au baccalauréat : il faut donc la maîtriser parfaitement.

\courssub{Observation : un mot français, plusieurs réalités espagnoles}

Observons d'abord les phrases suivantes, tirées de situations de communication quotidiennes :

\begin{exemplebox}[Observer avant de comprendre]
\begin{itemize}
\item \esp{Il va chez le médecin.} $\rightarrow$ \esp{Va al médico.}
\item \esp{Je dîne chez Pierre.} $\rightarrow$ \esp{Ceno en casa de Pedro.}
\item \esp{Chez les Espagnols, le repas de midi est sacré.} $\rightarrow$ \esp{Entre los españoles, la comida del mediodía es sagrada.}
\item \esp{Chez Cervantes, on trouve déjà l'humour moderne.} $\rightarrow$ \esp{En Cervantes, encontramos ya el humor moderno.}
\end{itemize}
\end{exemplebox}

Que remarquons-nous ? Le mot français \cle{chez} disparaît complètement en espagnol et se traduit différemment à chaque fois : \esp{al}, \esp{en casa de}, \esp{entre los}, \esp{en}. Il n'existe donc \textbf{aucun équivalent unique} de « chez » en espagnol. La bonne méthode consiste à se poser une question préalable : \emph{que signifie « chez » dans cette phrase ?} La réponse détermine la traduction.

\begin{propriete}[Règle générale : « chez » n'a pas d'équivalent unique]
Selon le sens, « chez » se traduit ainsi :
\begin{itemize}
\item \textbf{État / lieu où l'on est} (chez une personne) : \esp{en casa de + nom} ;
\item \textbf{Mouvement vers} (chez une personne) : \esp{a casa de + nom} ;
\item \textbf{Professionnel, commerçant} : \esp{al / a la + métier} (mouvement) ou \esp{en el / en la + métier} ou \esp{en la + boutique} (lieu) ;
\item \textbf{Peuple, groupe social} (dans les mœurs, les habitudes de) : \esp{entre los / las + nom} ;
\item \textbf{Auteur, penseur} (dans l'œuvre de) : \esp{en + nom}.
\end{itemize}
\end{propriete}

\courssub{Cas 1 : « chez » + personne, état ou lieu (\esp{en casa de})}

Quand « chez » désigne le \textbf{lieu où quelqu'un se trouve} (la maison d'une personne), l'espagnol utilise la structure \esp{en casa de + nom de la personne}. Le verbe est un verbe d'état : \esp{estar}, \esp{quedarse}, \esp{pasar la noche}, \esp{cenar}, \esp{vivir}...

\begin{exemplebox}[Exemples traduits]
\begin{itemize}
\item \esp{Estoy en casa de mi tío.} = Je suis chez mon oncle.
\item \esp{Pasamos la noche en casa de los abuelos.} = Nous passons la nuit chez les grands-parents.
\item \esp{Los niños juegan en casa de su vecino.} = Les enfants jouent chez leur voisin.
\item \esp{Anoche cenamos en casa de mis suegros.} = Hier soir, nous avons dîné chez mes beaux-parents.
\end{itemize}
\end{exemplebox}

Remarque importante : avec les adjectifs possessifs, on préfère souvent la forme \esp{en mi casa}, \esp{en tu casa}, \esp{en su casa} :

\begin{itemize}
\item \esp{Te espero en mi casa.} = Je t'attends chez moi.
\item \esp{¿Estás en tu casa?} = Tu es chez toi ?
\end{itemize}

\begin{attention}[Piège : jamais « en casa de mí »]
Les francophones traduisent mécaniquement « chez moi » par \esp{en casa de mí} : c'est une \textbf{erreur grave}. On dit \esp{en mi casa} (je suis chez moi) ou \esp{a mi casa} (je rentre chez moi). De même : « chez toi » = \esp{en tu casa}, « chez lui / chez elle » = \esp{en su casa}.
\end{attention}

\courssub{Cas 2 : « chez » + personne, mouvement (\esp{a casa de})}

Quand « chez » exprime un \textbf{mouvement vers} la maison de quelqu'un, la préposition \esp{en} devient \esp{a} : \esp{a casa de + nom}. Les verbes sont des verbes de mouvement : \esp{ir}, \esp{venir}, \esp{llegar}, \esp{volver}, \esp{regresar}...

\begin{exemplebox}[Exemples traduits]
\begin{itemize}
\item \esp{Voy a casa de mi prima.} = Je vais chez ma cousine.
\item \esp{¿Vienes a casa de los Martínez esta noche?} = Tu viens chez les Martínez ce soir ?
\item \esp{Llegamos a casa de nuestros amigos a las ocho.} = Nous sommes arrivés chez nos amis à huit heures.
\item \esp{Vuelvo a mi casa después del partido.} = Je rentre chez moi après le match.
\end{itemize}
\end{exemplebox}

\begin{aretenir}[La logique \esp{en} / \esp{a}]
C'est la même opposition que partout en espagnol : \esp{en} marque le \textbf{lieu où l'on est}, \esp{a} marque la \textbf{direction}. Comparez : \esp{Estoy en casa de Ana} (je suis chez Ana) / \esp{Voy a casa de Ana} (je vais chez Ana). C'est exactement la logique de \esp{estar en} / \esp{ir a}.
\end{aretenir}

\courssub{Cas 3 : « chez » + professionnel ou commerçant (\esp{al}, \esp{en el}, \esp{en la}...)}

C'est le cas le plus déroutant pour les francophones. Quand « chez » désigne le \textbf{cabinet, la boutique ou l'atelier d'un professionnel}, l'espagnol n'utilise \textbf{jamais} \esp{casa}. Deux solutions :

\begin{enumerate}
\item \textbf{Avec le nom du métier} : \esp{al / a la + métier} (mouvement) ou \esp{en el / en la + métier} (lieu) : \esp{Voy al médico} = Je vais chez le médecin.
\item \textbf{Avec le nom du commerce} (souvent en \esp{-ería}) : \esp{en la panadería} = chez le boulanger ; \esp{en la peluquería} = chez le coiffeur.
\end{enumerate}

\begin{exemplebox}[Exemples traduits]
\begin{itemize}
\item \esp{Voy al médico porque me duele la cabeza.} = Je vais chez le médecin parce que j'ai mal à la tête.
\item \esp{Compro el pan en la panadería.} = J'achète le pain chez le boulanger.
\item \esp{Mi madre está en la peluquería.} = Ma mère est chez le coiffeur.
\item \esp{Llevo el traje al sastre.} = J'apporte le costume chez le tailleur.
\item \esp{Hay que ir al dentista dos veces al año.} = Il faut aller chez le dentiste deux fois par an.
\end{itemize}
\end{exemplebox}

\begin{attention}[Erreur fréquente : « chez le coiffeur »]
Ne traduisez jamais « chez le coiffeur » par \esp{en casa del peluquero} : cela signifierait « à la maison du coiffeur », c'est-à-dire dans son domicile privé ! La bonne traduction est \esp{en la peluquería} (au salon de coiffure) ou \esp{al peluquero} (je vais chez le coiffeur). De même, « chez le boucher » = \esp{en la carnicería}, « chez le pharmacien » = \esp{en la farmacia}.
\end{attention}

\begin{savaistu}[Les noms de boutiques en \esp{-ería}]
En espagnol, le nom du commerçant donne souvent le nom de la boutique avec le suffixe \esp{-ería} : \esp{panadero} $\rightarrow$ \esp{panadería} (boulangerie), \esp{peluquero} $\rightarrow$ \esp{peluquería} (salon de coiffure), \esp{carnicero} $\rightarrow$ \esp{carnicería} (boucherie), \esp{zapatero} $\rightarrow$ \esp{zapatería} (cordonnerie, magasin de chaussures), \esp{librero} $\rightarrow$ \esp{librería} (librairie). C'est une ressource précieuse pour traduire « chez le... ». Au Cameroun, on dira par exemple : \esp{Compro la carne en la carnicería del mercado de Mokolo, en Yaundé.}
\end{savaistu}

\courssub{Cas 4 : « chez » + peuple ou groupe (\esp{entre los / las})}

Quand « chez » signifie « \textbf{dans les mœurs, les habitudes, la culture de} » un peuple ou d'un groupe humain, l'espagnol utilise \esp{entre los / las + nom}. Le sens littéral est « parmi ».

\begin{exemplebox}[Exemples traduits]
\begin{itemize}
\item \esp{Entre los españoles, la siesta es tradicional.} = Chez les Espagnols, la sieste est traditionnelle.
\item \esp{Entre los bamiléké, la familia es sagrada.} = Chez les Bamiléké, la famille est sacrée.
\item \esp{Entre los jóvenes cameruneses, el fútbol es una pasión.} = Chez les jeunes Camerounais, le football est une passion.
\item \esp{Entre los guineoecuatorianos, se habla español desde hace siglos.} = Chez les Guinéo-Équatoriens, on parle espagnol depuis des siècles.
\end{itemize}
\end{exemplebox}

\begin{attention}[Ne pas confondre les deux « chez les »]
Attention à la nuance : « Nous dormons chez les Espagnols » (dans la maison d'Espagnols précis) = \esp{Dormimos en casa de los españoles}. Mais « Chez les Espagnols, on dîne tard » (dans la culture espagnole en général) = \esp{Entre los españoles, se cena tarde}. Le contexte décide : maison concrète $\rightarrow$ \esp{en casa de} ; mœurs générales $\rightarrow$ \esp{entre los}.
\end{attention}

\courssub{Cas 5 : « chez » + auteur ou penseur (\esp{en + nom})}

Dans le commentaire de texte et la dissertation, « chez » signifie souvent « \textbf{dans l'œuvre de} », « dans la pensée de ». L'espagnol utilise alors simplement \esp{en + nom}.

\begin{exemplebox}[Exemples traduits]
\begin{itemize}
\item \esp{En Cervantes, encontramos ya el humor moderno.} = Chez Cervantes, on trouve déjà l'humour moderne.
\item \esp{En Machado, la poesía es una meditación sobre el tiempo.} = Chez Machado, la poésie est une méditation sur le temps.
\item \esp{En García Lorca, la muerte es un tema obsesivo.} = Chez García Lorca, la mort est un thème obsessionnel.
\end{itemize}
\end{exemplebox}

\courssub{Tableau récapitulatif et schéma de décision}

\begin{center}
\begin{tabularx}{\linewidth}{|X|X|X|}
\hline
\rowcolor{popblueL} \textbf{Français} & \textbf{Español} & \textbf{Remarque} \\
\hline
Je suis chez mon oncle. & Estoy en casa de mi tío. & État, personne : \esp{en casa de} \\
\hline
Je vais chez ma cousine. & Voy a casa de mi prima. & Mouvement, personne : \esp{a casa de} \\
\hline
Je vais chez le médecin. & Voy al médico. & Professionnel : \esp{al / a la + métier} \\
\hline
J'achète le pain chez le boulanger. & Compro el pan en la panadería. & Commerce : \esp{en la + boutique} \\
\hline
Chez les Espagnols, on dîne tard. & Entre los españoles, se cena tarde. & Peuple, mœurs : \esp{entre los} \\
\hline
Chez Cervantes, tout est humour. & En Cervantes, todo es humor. & Auteur, œuvre : \esp{en + nom} \\
\hline
Je suis chez moi. / Je rentre chez moi. & Estoy en mi casa. / Vuelvo a mi casa. & Jamais \esp{en casa de mí} ! \\
\hline
\end{tabularx}
\end{center}

\begin{popfigure}
\begin{tikzpicture}[
  node distance=0.55cm and 1.1cm,
  start/.style={rectangle, rounded corners, draw=popdark, fill=popgoldL, align=center, font=\bfseries, inner sep=6pt},
  q/.style={rectangle, rounded corners, draw=popblue, fill=popblueL, align=center, inner sep=5pt, font=\small},
  sol/.style={rectangle, rounded corners, draw=popgreen, fill=popgreenL, align=center, inner sep=5pt, font=\small},
  fleche/.style={-{Stealth[length=2.5mm]}, thick, popdark}]
\node[start, align=center] (chez){« CHEZ »\\ en français};
\node[q, below left=0.9cm and 2.2cm of chez] (pers) {une personne ?};
\node[q, below right=0.9cm and 2.2cm of chez] (autre) {autre cas ?};
\node[sol, below left=0.8cm and 0.4cm of pers, align=center] (estado){état :\\ \esp{en casa de}};
\node[sol, below right=0.8cm and 0.4cm of pers, align=center] (mov){mouvement :\\ \esp{a casa de}};
\node[sol, below=0.8cm of autre, xshift=-1.9cm, align=center] (pro){professionnel :\\ \esp{al médico} / \esp{en la farmacia}};
\node[sol, below=0.8cm of autre, xshift=1.9cm, align=center] (pueblo){peuple : \esp{entre los...}\\ auteur : \esp{en Cervantes}};
\draw[fleche] (chez) -- (pers);
\draw[fleche] (chez) -- (autre);
\draw[fleche] (pers) -- (estado) node[midway, left, font=\scriptsize\itshape]{estar};
\draw[fleche] (pers) -- (mov) node[midway, right, font=\scriptsize\itshape]{ir};
\draw[fleche] (autre) -- (pro);
\draw[fleche] (autre) -- (pueblo);
\end{tikzpicture}
\legende{Arbre de décision : comment traduire « chez » selon le sens.}
\end{popfigure}

\courssub{Exercice résolu}

\begin{exoresolu}[Thème : traduisez en espagnol]
Énoncé. Traduisez : 1) Nous dînons chez nos voisins ce soir. 2) Elle va chez le dentiste demain. 3) Chez les Camerounais, l'accueil est une valeur essentielle. 4) Je rentre chez moi après les cours. 5) Chez Bécquer, l'amour est toujours lié au rêve.
\tcblower
Solution détaillée.
\begin{enumerate}
\item \esp{Esta noche cenamos en casa de nuestros vecinos.} — État (on dîne = on est quelque part), personnes : \esp{en casa de}.
\item \esp{Ella va al dentista mañana.} — Mouvement vers un professionnel : \esp{al + métier} (contraction \esp{a + el = al}).
\item \esp{Entre los cameruneses, la acogida es un valor esencial.} — Mœurs d'un peuple : \esp{entre los}.
\item \esp{Vuelvo a mi casa después de las clases.} — « Chez moi » avec mouvement : \esp{a mi casa} (jamais \esp{a casa de mí}).
\item \esp{En Bécquer, el amor está siempre unido al sueño.} — Dans l'œuvre d'un auteur : \esp{en + nom}.
\end{enumerate}
\end{exoresolu}

\begin{exercice}[À vous de jouer]
Traduisez en espagnol : 1) Les enfants jouent chez leur grand-mère. 2) Nous allons chez le coiffeur samedi. 3) Chez les jeunes, les réseaux sociaux occupent beaucoup de temps. 4) La fête a lieu chez Ana. 5) Chez Mistral, l'enfance est un thème central. 6) Tu es chez toi ?
\end{exercice}

\begin{corrige}[Exercice « À vous de jouer »]
\begin{enumerate}
\item \esp{Los niños juegan en casa de su abuela.} (état, personne)
\item \esp{Vamos a la peluquería el sábado.} ou \esp{Vamos al peluquero el sábado.} (professionnel)
\item \esp{Entre los jóvenes, las redes sociales ocupan mucho tiempo.} (groupe, habitudes)
\item \esp{La fiesta es en casa de Ana.} (lieu de la fête : état)
\item \esp{En Mistral, la infancia es un tema central.} (œuvre d'un auteur)
\item \esp{¿Estás en tu casa?} (chez toi : possessif, jamais \esp{casa de ti})
\end{enumerate}
\end{corrige}

\begin{savaistu}[« En casa de » : la fête est chez Ana]
La structure \esp{en casa de} sert aussi à situer un \textbf{événement} : \esp{La fiesta es en casa de Ana} = la fête a lieu chez Ana ; \esp{La reunión será en casa del director} = la réunion aura lieu chez le directeur. Dans ce sens, « chez » garde sa valeur de lieu concret. En revanche, pour un événement dans un pays ou une ville, on utilise \esp{en} : \esp{El Mundial se celebra en España} = la Coupe du monde se déroule en Espagne.
\end{savaistu}

\begin{aretenir}[L'essentiel sur « chez »]
\begin{itemize}
\item « Chez » n'a \textbf{pas d'équivalent unique} : analysez d'abord le sens.
\item Personne, état : \esp{en casa de} ; personne, mouvement : \esp{a casa de}.
\item Professionnel : \esp{al médico}, \esp{en la panadería} — jamais \esp{casa del médico}.
\item Peuple, mœurs : \esp{entre los españoles} ; auteur, œuvre : \esp{en Cervantes}.
\item « Chez moi / toi / lui » : \esp{en mi casa / en tu casa / en su casa} (mouvement : \esp{a mi casa...}).
\end{itemize}
\end{aretenir}

Maintenant que nous savons traduire « chez », il nous reste un dernier verbe-piège du thème : « devenir », qui se dit \esp{ponerse}, \esp{volverse}, \esp{hacerse}, \esp{convertirse en} ou \esp{llegar a ser} selon le sens. C'est l'objet de la section suivante.

\coursec{Grammaire 3 : traduire « devenir »}

Après avoir appris à traduire « chez » (\esp{en casa de}, \esp{entre los españoles}...), nous abordons un autre grand piège de la traduction français-espagnol : le verbe « devenir ». En français, un seul verbe suffit dans presque tous les contextes. En espagnol, il n'existe pas de verbe unique : il faut choisir entre \esp{ponerse}, \esp{volverse}, \esp{hacerse}, \esp{convertirse en}, \esp{llegar a ser} et \esp{quedarse}, selon la \cle{nature du changement}. C'est l'une des questions les plus fréquentes au thème du baccalauréat.

\courssub{Observation : un français, six espagnols}

\begin{exemplebox}[Trois phrases, trois verbes différents]
Observons trois traductions du verbe « devenir » :

\begin{itemize}
\item \esp{Se hizo médico.} « Il est devenu médecin. »
\item \esp{Se volvió loca.} « Elle est devenue folle. »
\item \esp{Se puso rojo.} « Il est devenu rouge. »
\end{itemize}

Le même verbe français « devenir » se traduit par trois verbes espagnols différents. Pourquoi ? Parce que les changements ne sont pas de même nature : devenir médecin est le résultat d'un \cle{effort voulu} ; devenir fou est un changement \cle{profond et involontaire} ; devenir rouge est un changement \cle{temporaire et physique}.
\end{exemplebox}

\begin{textebox}[Una vida de cambios]
Mi tío Andrés nació en un pueblo pequeño cerca de Malabo. De niño, era muy tímido: se ponía rojo cada vez que alguien le hablaba en clase. Pero con los años se volvió muy serio y trabajador. Estudió mucho, se hizo profesor de español y, después de veinte años de esfuerzo, llegó a ser director de un instituto. Su mujer, Carmen, se quedó viuda muy joven antes de conocerlo, y por eso se convirtió en una persona muy fuerte. Hoy, cuando mi tío cuenta su historia, dice siempre: « Nadie se hace importante de un día para otro ». Su hijo mayor, que se hizo médico el año pasado, se puso muy contento cuando recibió su primer sueldo. La vida nos cambia, dice mi tío, pero cada cambio tiene su camino.
\end{textebox}

\textbf{Glossaire} : \esp{ponerse rojo} = devenir rouge ; \esp{volverse serio} = devenir sérieux ; \esp{hacerse profesor} = devenir professeur ; \esp{llegar a ser director} = finir par devenir directeur ; \esp{quedarse viuda} = devenir veuve ; \esp{convertirse en} = se transformer en ; \esp{de un día para otro} = du jour au lendemain ; \esp{el sueldo} = le salaire.

\textbf{Questions de compréhension} : 1. Quels changements le narrateur attribue-t-il à son oncle ? 2. Pourquoi utilise-t-on \esp{se quedó viuda} et non \esp{se hizo viuda} ? 3. Relevez tous les verbes de changement du texte et classez-les.

\courssub{La règle générale : choisir selon la nature du changement}

\begin{propriete}[Règle de traduction de « devenir »]
Le choix du verbe espagnol dépend de la \cle{nature du changement} :

\begin{itemize}
\item Changement \textbf{temporaire}, physique ou émotionnel : \esp{ponerse} + adjectif.
\item Changement \textbf{durable et involontaire}, profond : \esp{volverse} + adjectif.
\item Changement \textbf{voulu}, résultat d'un effort (profession, religion, richesse) : \esp{hacerse} + nom ou adjectif.
\item \textbf{Transformation} complète, souvent spectaculaire : \esp{convertirse en} + nom.
\item \textbf{Aboutissement} après un long processus : \esp{llegar a ser} + nom.
\item \textbf{État définitif résultant d'un événement} : \esp{quedarse} + adjectif ou participe.
\end{itemize}
\end{propriete}

\begin{popfigure}
\begin{tikzpicture}[font=\small]
\draw[-{Stealth}, thick, popline] (0,0) -- (15.5,0);
\node[draw, fill=popblueL, rounded corners, align=center, text width=2.1cm] at (1.3,1.4) {\esp{ponerse}\\temporaire};
\node[draw, fill=poppinkL, rounded corners, align=center, text width=2.1cm] at (3.9,-1.4) {\esp{volverse}\\durable\\involontaire};
\node[draw, fill=popgreenL, rounded corners, align=center, text width=2.1cm] at (6.5,1.4) {\esp{hacerse}\\voulu, effort};
\node[draw, fill=poporangeL, rounded corners, align=center, text width=2.1cm] at (9.1,-1.4) {\esp{convertirse en}\\transformation};
\node[draw, fill=popgoldL, rounded corners, align=center, text width=2.1cm] at (11.7,1.4) {\esp{llegar a ser}\\aboutissement};
\node[draw, fill=poppurpleL, rounded corners, align=center, text width=2.1cm] at (14.3,-1.4) {\esp{quedarse}\\état résultant};
\foreach \x in {1.3,3.9,6.5,9.1,11.7,14.3} \fill[popdark] (\x,0) circle (1.5pt);
\draw[popmuted] (1.3,0) -- (1.3,0.9);
\draw[popmuted] (3.9,0) -- (3.9,-0.9);
\draw[popmuted] (6.5,0) -- (6.5,0.9);
\draw[popmuted] (9.1,0) -- (9.1,-0.9);
\draw[popmuted] (11.7,0) -- (11.7,0.9);
\draw[popmuted] (14.3,0) -- (14.3,-0.9);
\end{tikzpicture}
\legende{Frise des six traductions de « devenir » selon la nature du changement.}
\end{popfigure}

\courssub{Étude détaillée de chaque verbe}

\textbf{1. \esp{ponerse} + adjectif : le changement temporaire}\par
\esp{Ponerse} exprime un changement \cle{passager}, souvent physique ou émotionnel : couleur du visage, humeur, réaction du moment. Il ne se construit \textbf{jamais avec un nom}.

\begin{itemize}
\item \esp{Se puso rojo de vergüenza.} « Il est devenu rouge de honte. »
\item \esp{Cuando vio la nota, se puso muy triste.} « Quand il a vu la note, il est devenu très triste. »
\item \esp{El niño se pone nervioso antes de los exámenes.} « L'enfant devient nerveux avant les examens. »
\item \esp{El cielo se puso gris.} « Le ciel est devenu gris. »
\end{itemize}

\textbf{2. \esp{volverse} + adjectif : le changement profond et involontaire}\par
\esp{Volverse} exprime un changement \cle{durable}, souvent \cle{involontaire}, qui touche le caractère ou l'état mental. L'adjectif s'accorde avec le sujet.

\begin{itemize}
\item \esp{Se ha vuelto muy callado desde el accidente.} « Il est devenu très silencieux depuis l'accident. »
\item \esp{El debate se volvió violento.} « Le débat est devenu violent. »
\item \esp{Mi abuela se está volviendo olvidadiza.} « Ma grand-mère devient oublieuse. »
\item \esp{Se volvió loca de alegría.} « Elle est devenue folle de joie. »
\end{itemize}

\textbf{3. \esp{hacerse} + nom ou adjectif : le changement voulu}\par
\esp{Hacerse} exprime un changement \cle{voulu}, fruit d'un effort ou d'une décision : profession, religion, idéologie, fortune. C'est le verbe typique des noms de métier (sans article : \esp{se hizo médico}, comme en français « il est devenu médecin »).

\begin{itemize}
\item \esp{Mi hermano se hizo profesor.} « Mon frère est devenu professeur. »
\item \esp{Se hizo rico con el comercio del cacao.} « Il est devenu riche grâce au commerce du cacao. »
\item \esp{Se hizo católica antes de la boda.} « Elle est devenue catholique avant le mariage. »
\item \esp{Nos hicimos amigos en el liceo.} « Nous sommes devenus amis au lycée. »
\end{itemize}

\textbf{4. \esp{convertirse en} + nom : la transformation}\par
\esp{Convertirse en} (toujours avec la préposition \esp{en}) exprime une \cle{transformation complète}, souvent spectaculaire ou inattendue. Attention : avec l'article indéfini, on garde \esp{un/una}.

\begin{itemize}
\item \esp{El agua se convirtió en hielo.} « L'eau est devenue glace. »
\item \esp{El juego se convirtió en un problema.} « Le jeu est devenu un problème. »
\item \esp{La oruga se convierte en mariposa.} « La chenille devient papillon. »
\end{itemize}

\textbf{5. \esp{llegar a ser} + nom : l'aboutissement}\par
\esp{Llegar a ser} insiste sur le \cle{long processus} qui mène au résultat : on « finit par devenir » après des années d'efforts.

\begin{itemize}
\item \esp{Llegó a ser presidente después de veinte años.} « Il est devenu président après vingt ans. »
\item \esp{Llegó a ser madre muy joven.} « Elle est devenue mère très jeune. »
\item \esp{Ese barrio llegó a ser el centro de la ciudad.} « Ce quartier est devenu le centre de la ville. »
\end{itemize}

\textbf{6. \esp{quedarse} : l'état résultant d'un événement}\par
\esp{Quedarse} exprime un état \cle{définitif} qui résulte d'un événement (accident, maladie, décès) : \esp{quedarse ciego} (devenir aveugle), \esp{sordo} (sourd), \esp{viudo} (veuf), \esp{calvo} (chauve), \esp{huérfano} (orphelin), \esp{paralítico} (paralysé).

\begin{itemize}
\item \esp{Se quedó ciego a los diez años.} « Il est devenu aveugle à dix ans. »
\item \esp{Su padre se quedó viudo el año pasado.} « Son père est devenu veuf l'année dernière. »
\item \esp{Se está quedando calvo.} « Il devient chauve. »
\end{itemize}

\courssub{Tableau récapitulatif}

\begin{center}
\begin{tabularx}{\linewidth}{|l|X|l|X|}
\hline
\rowcolor{popblueL} \textbf{Verbe} & \textbf{Type de changement} & \textbf{Construction} & \textbf{Exemple traduit} \\
\hline
\esp{ponerse} & temporaire, physique, émotionnel & + adjectif & \esp{Se puso rojo.} « Il est devenu rouge. » \\
\hline
\esp{volverse} & durable, profond, involontaire & + adjectif & \esp{Se volvió loca.} « Elle est devenue folle. » \\
\hline
\esp{hacerse} & voulu, effort, décision & + nom / adjectif & \esp{Se hizo médico.} « Il est devenu médecin. » \\
\hline
\esp{convertirse en} & transformation complète & \esp{en} + nom & \esp{Se convirtió en un problema.} « C'est devenu un problème. » \\
\hline
\esp{llegar a ser} & aboutissement, long processus & + nom & \esp{Llegó a ser director.} « Il est devenu directeur. » \\
\hline
\esp{quedarse} & état définitif après un événement & + adjectif / participe & \esp{Se quedó ciego.} « Il est devenu aveugle. » \\
\hline
\end{tabularx}
\end{center}

\begin{exemplebox}[Exemples de synthèse]
\begin{itemize}
\item « Mon frère est devenu professeur. » → \esp{Mi hermano se hizo profesor.} (profession choisie)
\item « Le débat est devenu violent. » → \esp{El debate se volvió violento.} (changement profond, involontaire)
\item « Elle est devenue mère très jeune. » → \esp{Llegó a ser madre muy joven.} (étape de vie, aboutissement)
\item « Il est devenu pâle de peur. » → \esp{Se puso pálido de miedo.} (réaction passagère)
\item « Le petit village est devenu une grande ville. » → \esp{El pequeño pueblo se convirtió en una gran ciudad.} (transformation)
\item « Elle est devenue veuve à trente ans. » → \esp{Se quedó viuda a los treinta años.} (état résultant d'un décès)
\end{itemize}
\end{exemplebox}

\begin{exoresolu}[Thème : traduisez en espagnol]
Énoncé. Traduisez : 1. « Elle est devenue institutrice. » 2. « Il est devenu rouge de colère. » 3. « Mon oncle est devenu aveugle après l'accident. » 4. « Internet est devenu un outil indispensable. » 5. « Après des années de travail, il est devenu le meilleur élève de la classe. »
\tcblower
Solution détaillée.

1. \esp{Se hizo maestra.} Profession choisie, résultat d'études : \esp{hacerse} + nom sans article.

2. \esp{Se puso rojo de ira.} Réaction physique et passagère : \esp{ponerse} + adjectif.

3. \esp{Mi tío se quedó ciego después del accidente.} État définitif résultant d'un événement : \esp{quedarse ciego}.

4. \esp{Internet se ha convertido en una herramienta indispensable.} Transformation d'une chose en autre chose : \esp{convertirse en} + nom avec article indéfini.

5. \esp{Después de años de trabajo, llegó a ser el mejor alumno de la clase.} Aboutissement d'un long processus : \esp{llegar a ser}.
\end{exoresolu}

\begin{attention}[Erreurs fréquentes des francophones]
\begin{itemize}
\item Ne dites jamais \esp{se puso médico} : \esp{ponerse} ne se construit \textbf{jamais avec un nom de profession}. Pour les métiers, utilisez \esp{hacerse} : \esp{se hizo médico}.
\item Le verbe français « devenir » n'a \textbf{pas d'équivalent direct} en espagnol : n'inventez pas un verbe \esp{devenir}. Choisissez toujours parmi les six verbes de la règle.
\item Ne confondez pas \esp{quedarse} (devenir, état résultant : \esp{se quedó ciego} = il est devenu aveugle) et \esp{quedar} (rester, se trouver : \esp{la casa queda lejos} = la maison se trouve loin).
\item Avec \esp{hacerse} + profession, pas d'article : \esp{se hizo abogado}, comme en français « il est devenu avocat ». Mais avec \esp{convertirse en}, l'article indéfini est obligatoire : \esp{se convirtió en un héroe}.
\end{itemize}
\end{attention}

\begin{savaistu}[\esp{Hacerse} pour l'âge et l'heure]
\esp{Hacerse} sert aussi dans des expressions impersonnelles liées au temps : \esp{Se hace tarde.} « Il se fait tard. » ; \esp{Se hizo de noche.} « La nuit est tombée. » ; \esp{Se hace mayor.} « Il se fait vieux. » On trouve aussi \esp{hacerse a} pour « s'habituer à » : \esp{Me hice al clima de Yaundé.} « Je me suis habitué au climat de Yaoundé. »
\end{savaistu}

\begin{exercice}[À vous de traduire]
Traduisez en espagnol en choisissant le bon verbe de changement : 1. « Elle est devenue très nerveuse avant l'examen. » 2. « Mon cousin est devenu ingénieur. » 3. « Le téléphone portable est devenu un problème en classe. » 4. « Il est devenu sourd après la maladie. » 5. « Cette élève est devenue la meilleure du lycée après des années d'efforts. » 6. « Le ciel est devenu noir avant la pluie. »
\end{exercice}

\begin{corrige}[Exercice « À vous de traduire »]
1. \esp{Se puso muy nerviosa antes del examen.} (émotion passagère : \esp{ponerse}) 2. \esp{Mi primo se hizo ingeniero.} (profession voulue : \esp{hacerse}) 3. \esp{El teléfono móvil se ha convertido en un problema en clase.} (transformation : \esp{convertirse en}) 4. \esp{Se quedó sordo después de la enfermedad.} (état résultant : \esp{quedarse}) 5. \esp{Esta alumna llegó a ser la mejor del liceo después de años de esfuerzos.} (aboutissement : \esp{llegar a ser}) 6. \esp{El cielo se puso negro antes de la lluvia.} (changement temporaire : \esp{ponerse})
\end{corrige}

\begin{aretenir}[L'essentiel]
« Devenir » n'a pas de traduction unique en espagnol : \esp{ponerse} (temporaire), \esp{volverse} (durable et involontaire), \esp{hacerse} (voulu, professions), \esp{convertirse en} (transformation), \esp{llegar a ser} (aboutissement), \esp{quedarse} (état résultant d'un événement). Avant de traduire, posez-vous toujours la question : quelle est la nature du changement ?
\end{aretenir}

\coursec{Débat : la educación de los niños y las NTIC}

Après avoir maîtrisé les subtilités de \cle{devenir} (\emph{ponerse, volverse, hacerse, convertirse en, llegar a ser, quedarse}), nous mettons maintenant en pratique l'ensemble des acquis de cette leçon — ordre/interdiction, « chez », « devenir » — dans une situation de communication authentique : le débat oral et l'expression écrite d'opinion. Le sujet, au cœur de l'actualité éducative au Cameroun comme en Espagne, oppose les bienfaits de l'accès au savoir numérique aux risques d'addiction et d'isolement.

\courssub{Présentation du sujet et texte déclencheur}

Le sujet du débat est formulé sous forme de question fermée, typique de l'épreuve orale du Baccalauréat : \esp{¿Deben los padres prohibir el móvil a los niños?} (Les parents doivent-ils interdire le portable aux enfants ?). Pour lancer la réflexion, voici un texte d'opinion original, publié dans la rubrique « Sociedad » d'un journal numérique hispanophone fictif \emph{El Diario de la Educación}.

\begin{textebox}{Texte d'opinion : « Pantallas: ¿herramientas o cadenas? »}
En mi opinión, la tecnología no es buena ni mala en sí misma; todo depende del uso que se haga en casa de cada familia. Es cierto que los smartphones permiten a los niños de la diáspora camerunesa en Madrid hablar con sus abuelos en Yaoundé cada domingo: la distancia se vuelve corta gracias a la videollamada. Sin embargo, muchos padres prohíben el móvil en la mesa y en el dormitorio, porque el exceso de pantallas convierte a los niños en seres pasivos que ya no juegan en la calle. Los pedagogos advierten: « No dejen que el algoritmo eduque a sus hijos ». En casa de mis vecinos, el niño de ocho años se ha vuelto agresivo cuando le quitan la tableta; se queda horas frente a la pantalla y su rendimiento escolar ha bajado. Por tanto, propongo un pacto familiar: horarios claros, control parental y, sobre todo, ejemplo de los adultos. ¿Y ustedes, qué opinan?
\end{textebox}

\begin{center}
\begin{tabularx}{\linewidth}{|X|X|}\hline
\rowcolor{popblueL} \textbf{Vocabulario difícil} & \textbf{Traduction / Explication} \\ \hline
\esp{en sí misma} & en soi / par nature \\ \hline
\esp{la diáspora camerunesa} & la diaspora camerounaise \\ \hline
\esp{se vuelve corta} & devient courte (\cle{volverse} + adjectif : changement progressif/inévitable) \\ \hline
\esp{seres pasivos} & des êtres passifs \\ \hline
\esp{el algoritmo} & l'algorithme \\ \hline
\esp{se ha vuelto agresivo} & est devenu agressif (changement de caractère durable) \\ \hline
\esp{el rendimiento escolar} & le rendement / les résultats scolaires \\ \hline
\esp{un pacto familiar} & un pacte familial / contrat de confiance \\ \hline
\end{tabularx}
\end{center}

\courssub{Boîte à outils linguistiques : structurer son argumentation}

Pour débattre efficacement en espagnol, il ne suffit pas d'avoir des idées ; il faut les \cle{articuler} avec des connecteurs logiques précis. Le français utilise souvent « car / parce que / comme » ; l'espagnol offre une palette plus nuancée selon le registre.

\begin{propriete}[Connecteurs du débat oral et écrit]
\textbf{1. Exprimer son opinion (Prise de position)} \\
\esp{Creo que / Pienso que / Opino que} + indicatif (opinion subjective). \\
\esp{A mi juicio / A mi modo de ver / Desde mi punto de vista} + indicatif ou subjonctif si doute. \\
\esp{Estoy convencido/a de que} + indicatif. \\[4pt]
\textbf{2. Justifier / Argumenter (Donner des raisons)} \\
\esp{Porque} (réponse à « ¿por qué ? », place après la proposition principale). \\
\esp{Ya que / Puesto que / Como} (place au début de la phrase, cause connue ou évidente). \\
\esp{Puesto que los niños necesitan dormir, hay que apagar el wifi.} \\[4pt]
\textbf{3. Concéder (Admettre un point adverse avant de rebondir)} \\
\esp{Es cierto que... pero / sin embargo / no obstante.} \\
\esp{Admito que... aunque...} + subjonctif si concession hypothétique (\esp{aunque sea cierto}). \\[4pt]
\textbf{4. Réfuter / Contredire (Poliment mais fermement)} \\
\esp{No estoy de acuerdo (con eso) porque...} \\
\esp{No comparto esa opinión ya que...} \\
\esp{Es falso que...} + subjonctif (\esp{Es falso que los móviles solo sirvan para jugar}). \\[4pt]
\textbf{5. Conclure (Synthétiser et clore)} \\
\esp{En conclusión / En resumen / Para terminar,} + phrase nominale ou verbale. \\
\esp{Por todo ello, considero que...}
\end{propriete}

\begin{exemplebox}[Exemples traduits : du français vers l'espagnol argumenté]
\textbf{Fr :} Je pense que les écrans sont utiles pour l'école. \\
\textbf{Es :} \esp{Creo que las pantallas son útiles para la escuela.} \\[4pt]
\textbf{Fr :} Certes, ils distraient, mais ils informent aussi. \\
\textbf{Es :} \esp{Es cierto que distraen, pero también informan.} \\[4pt]
\textbf{Fr :} Je ne suis pas d'accord : les enfants deviennent accros. \\
\textbf{Es :} \esp{No estoy de acuerdo: los niños se vuelven adictos.} (\cle{volverse} = devenir durablement) \\[4pt]
\textbf{Fr :} En conclusion, il faut interdire le portable avant 12 ans. \\
\textbf{Es :} \esp{En conclusión, hay que prohibir el móvil antes de los doce años.} (Infinitif d'interdiction générale)
\end{exemplebox}

\courssub{Arguments : pour et contre les NTIC}

Un bon débat repose sur un inventaire équilibré. Le tableau ci-dessous, à mémoriser comme une « fiche de révision », reprend les arguments standards enrichis du contexte camerounais (diaspora, réalités locales).

\begin{popfigure}
\centering
\begin{tikzpicture}[
    node distance=0.2cm,
    header/.style={fill=popblueL, text=popdark, font=\bfseries\small, minimum height=0.8cm, align=center},
    argfor/.style={fill=popgreenL, text=popdark, font=\small, minimum height=0.7cm, align=left, text width=7cm},
    argagainst/.style={fill=poporangeL, text=popdark, font=\small, minimum height=0.7cm, align=left, text width=7cm},
    title/.style={font=\bfseries\normalsize, text=popdark}
]
% Colonnes
\node[header, minimum width=8cm] (tf) {A FAVOR (Pour les NTIC)};
\node[header, minimum width=8cm, right=0cm of tf] (tc) {EN CONTRA (Contre les NTIC)};

% Arguments Pour
\node[argfor, below=0cm of tf] (f1) {1. \textbf{Acceso al saber:} Wikipedia, Khan Academy, MOOCs. En zonas rurales de Camerún, el móvil es la única biblioteca.};
\node[argfor, below=0cm of f1] (f2) {2. \textbf{Preparación laboral:} Competencias digitales obligatorias. \emph{``Los niños se hacen expertos en código''} (\cle{hacerse} + nom).};
\node[argfor, below=0cm of f2] (f3) {3. \textbf{Vínculo familiar (Diáspora):} Videollamadas con abuelos en Yaoundé/Douala. \emph{``En casa de mis tíos, la tablet une a la familia''}.};

% Arguments Contre
\node[argagainst, below=0cm of tc] (c1) {1. \textbf{Adicción y dopamina:} Juegos, TikTok, Reels. \emph{``Los niños se ponen ansiosos sin móvil''} (\cle{ponerse} + état passager).};
\node[argagainst, below=0cm of c1] (c2) {2. \textbf{Fracaso escolar y sueño:} Luz azul, insomnio, copia/pega. \emph{``El rendimiento se vuelve mediocre''} (\cle{volverse}).};
\node[argagainst, below=0cm of c2] (c3) {3. \textbf{Peligros:} Ciberacoso, pornografía, datos. \emph{``Hay que prohibir el acceso sin control parental''}.};

% Flèches décoratives
\draw[popgreen, thick, ->] (tf.south west) -- ++(-0.5,0) node[left, xshift=-0.6cm, rotate=90, anchor=south] {\tiny ARGUMENTS};
\draw[poporange, thick, ->] (tc.south east) -- ++(0.5,0) node[right, xshift=0.6cm, rotate=-90, anchor=south] {\tiny CONTRE-ARGUMENTS};
\end{tikzpicture}
\legende{Tableau récapitulatif des arguments pour le débat « NTIC et éducation ».}
\end{popfigure}

\begin{savaistu}[Le Cameroun et la Guinée équatoriale face aux NTIC]
Au Cameroun, le taux de pénétration mobile dépasse 80\% selon les estimations récentes de l'ARTC. À Yaoundé comme à Douala, le « bendskin » (moto-taxi) attend son client en scrollant TikTok. En Guinée équatoriale, seul pays hispanophone d'Afrique, l'espagnol est la langue de l'école et de l'administration ; les enfants y grandissent entre \esp{WhatsApp} en espagnol et \esp{Fang} ou \esp{Bubi} à la maison. Le débat « móvil sí / móvil no » y est vif : les parents équatoguinéens craignent la perte de l'espagnol normatif au profit d'un « espanglish » numérique, tandis que les manuels numériques (tablettes éducatives) arrivent dans les lycées de Malabo et Bata.
\end{savaistu}

\courssub{Contrainte linguistique intégrée : le « triple défi »}

Pour transformer ce débat en exercice d'évaluation sommative, chaque intervention orale ou paragraphe écrit \textbf{doit} contenir obligatoirement :
\begin{enumerate}
    \item Au moins \textbf{une forme d'ordre ou d'interdiction} (impératif, \esp{no} + subjonctif, \esp{hay que / se prohíbe / prohibido} + infinitif).
    \item Au moins \textbf{un « chez » correct} (\esp{en casa de}, \esp{a casa de}, \esp{entre los...}).
    \item Au moins \textbf{un verbe de « devenir » bien choisi} (\esp{ponerse}, \esp{volverse}, \esp{hacerse}, \esp{convertirse en}, \esp{llegar a ser}, \esp{quedarse}).
\end{enumerate}

\begin{exemplebox}[Modèles respectant la contrainte « Triple Défi »]
\textbf{Élève A (Pour) :} \esp{``A mi juicio, los padres \textbf{no deben prohibir} totalmente el móvil, porque \textbf{en casa de} mis primos la tableta \textbf{ayuda a que los niños se hagan} autónomos.''} \\
\checkmark Interdiction (\esp{no deben prohibir}) | \checkmark Chez (\esp{en casa de}) | \checkmark Devenir (\esp{se hagan} = \cle{hacerse} + adjectif/nom). \\[6pt]
\textbf{Élève B (Contre) :} \esp{``\textbf{Prohiban} las pantallas en la cena; \textbf{entre los españoles} muchos niños \textbf{se vuelven} adictos y \textbf{se quedan} despiertos hasta tarde.''} \\
\checkmark Ordre (\esp{Prohiban} - impératif \emph{ustedes}) | \checkmark Chez (\esp{entre los españoles}) | \checkmark Devenir (\esp{se vuelven} + \esp{se quedan}). \\[6pt]
\textbf{Élève C (Nuancé) :} \esp{``\textbf{Hay que limitar} el tiempo; \textbf{a casa de} mi amigo, el niño \textbf{se puso} nervioso sin wifi y sus notas \textbf{llegaron a ser} malas.''} \\
\checkmark Obligation/Interdiction (\esp{Hay que limitar}) | \checkmark Chez (\esp{a casa de}) | \checkmark Devenir (\esp{se puso} état passager, \esp{llegaron a ser} résultat final).
\end{exemplebox}

\begin{exoresolu}[Application : Transformer en espagnol avec la contrainte « Triple Défi »]
\textbf{Consigne :} Traduisez les phrases suivantes en espagnol en respectant les 3 contraintes grammaticales (ordre/interdiction, « chez », « devenir »). Justifiez le choix du verbe de devenir.\\[4pt]
\textbf{1. Fr :} « Interdisez les portables à table ! Chez nous, les enfants deviennent fâchés sans écran. » \\
\textbf{2. Fr :} « Il faut que les parents contrôlent. Chez les Camerounais, le portable devient un outil de travail. » \\
\textbf{3. Fr :} « Ne laissez pas l'algorithme éduquer. À la maison, l'enfant reste passif devant l'écran. » \\[4pt]
\tcblower
\textbf{Solution détaillée :} \\[4pt]
\textbf{1.} \esp{``\textbf{Prohiban} los móviles en la mesa! \textbf{En casa de} nosotros, los niños \textbf{se ponen} enfadados sin pantalla.''} \\
\quad \textit{Justification :} \esp{Prohiban} (impératif \emph{ustedes} = ordre formel/classe). \esp{En casa de} (chez nous, lieu physique/ambiance). \esp{Se ponen} (\cle{ponerse} + adjectif = changement d'état \textbf{passager}, réaction immédiate). \\[4pt]
\textbf{2.} \esp{``\textbf{Los padres deben / tienen que controlar.} \textbf{Entre los cameruneses}, el móvil \textbf{se convierte en} una herramienta de trabajo.''} \\
\quad \textit{Justification :} \esp{Deben/Tienen que} (obligation forte = ordre indirect). \esp{Entre los cameruneses} (chez les Camerounais = groupe humain, généralité culturelle). \esp{Se convierte en} (\cle{convertirse en} + nom = transformation \textbf{de nature/fonction}, processus abouti). \\[4pt]
\textbf{3.} \esp{``\textbf{No dejen} que el algoritmo eduque. \textbf{En la casa}, el niño \textbf{se queda} pasivo frente a la pantalla.''} \\
\quad \textit{Justification :} \esp{No dejen} (impératif \emph{ustedes} négatif = interdiction formelle). \esp{En la casa} (chez soi, domicile précis). \esp{Se queda} (\cle{quedarse} + adjectif = \textbf{résultat final} d'un processus, état résiduel : il \emph{reste} passif).
\end{exoresolu}

\courssub{Organisation du débat en classe (Simulation Bac)}

La mise en situation reproduit les conditions de l'oral du Baccalauréat (LV2, 10 min de préparation, 10 min de passage).

\begin{experience}[Mise en place du débat « Comisión de Educación »]
\textbf{Rôles (groupes de 5-6 élèves) :}
\begin{itemize}
    \item \textbf{1 Modérateur (Presidente) :} Ouvre la séance, donne la parole, veille au temps, reformule, clôture.
    \item \textbf{2-3 Débateurs « A favor » :} Défendent l'usage encadré des NTIC.
    \item \textbf{2-3 Débateurs « En contra » :} Défendent la restriction stricte / l'interdiction avant 14-16 ans.
    \item \textbf{1 Secrétaire (Secretario) :} Note les arguments clés sur un tableau (mots-clés en espagnol), vérifie la contrainte « Triple Défi » chez chaque orateur.
\end{itemize}
\textbf{Déroulement (15 min) :}
\begin{enumerate}
    \item \textbf{2 min :} Lecture du sujet par le modérateur. \esp{``El tema de hoy: ¿Deben los padres prohibir el móvil a los niños?''}
    \item \textbf{3 min :} Tour de table « Position initiale » (30 sec par débateur). \textit{Obligation : utiliser la contrainte Triple Défi.}
    \item \textbf{6 min :} Échanges libres / Réfutations. Le modérateur relance : \esp{``¿Qué responde el equipo contrario?''}, \esp{``Un ejemplo concreto, por favor.''}
    \item \textbf{2 min :} Synthèse du secrétaire (lecture des arguments au tableau).
    \item \textbf{2 min :} Conclusion du modérateur + vote symbolique à main levée.
\end{enumerate}
\textbf{Évaluation (Grille simplifiée) :} Contenu (arguments variés, exemples contextuels) / 5 -- Grammaire (Triple Défi respecté, subjonctif après \esp{no creo que}, \esp{es necesario que}) / 5 -- Fluidité / Prononciation / 5 -- Interaction (écoute, réactivité, politesse) / 5.
\end{experience}

\courssub{Méthode : réussir un débat oral au Bac}

\begin{methode}[Les 5 étapes d'une intervention gagnante]
\textbf{Étape 1 : Saluer et prendre position clairement (10-15 sec).} \\
\esp{``Buenos días. A mi juicio, los padres no deben prohibir el móvil, sino regular su uso.''} \\
$\rightarrow$ Phrase courte, indicatif après \esp{opino/creo}, ton affirmatif. \\[4pt]
\textbf{Étape 2 : Argumenter avec 2-3 arguments distincts (40-50 sec).} \\
Utiliser \esp{En primer lugar...}, \esp{Además...}, \esp{Por último...}. \\
Chaque argument = Idée force + Connecteur (\esp{ya que / porque}) + Exemple concret (chiffre, anecdote, référence culturelle). \\
\esp{``En primer lugar, el móvil permite la comunicación con la diáspora. Ya que muchos cameruneses viven en España, la videollamada acorta distancias.''} \\[4pt]
\textbf{Étape 3 : Illustrer par un exemple personnel ou localisé (15 sec).} \\
\esp{``En casa de mi prima en Douala, su hijo aprende inglés con una app gratuita.''} \\
$\rightarrow$ Ancre le débat dans le réel, montre la maîtrise de « chez ». \\[4pt]
\textbf{Étape 4 : Réfuter poliment l'adversaire (20 sec).} \\
Écouter $\rightarrow$ Concéder \esp{``Es cierto que...''} $\rightarrow$ Pivoter \esp{``pero / sin embargo...''} $\rightarrow$ Contre-argument. \\
\esp{``Es cierto que existe la adicción, pero prohibir todo genera rebeldía. Lo mejor es un pacto familiar.''} \\[4pt]
\textbf{Étape 5 : Conclure avec une ouverture (10 sec).} \\
\esp{``En conclusión, la tecnología es una herramienta: la educación digital empieza en casa. Gracias.''}
\end{methode}

\begin{attention}[Pièges lexicaux : \esp{debatir} vs \esp{discutir} / \esp{discutir} vs \esp{pelearse}]
\begin{itemize}
    \item \esp{Debatir} existe (du latin \emph{debattere}), mais c'est un terme formel, académique, peu utilisé à l'oral courant. On dit : \esp{``Vamos a debatir el tema en clase''} (contexte très formel).
    \item \textbf{Le verbe standard pour « débattre / discuter un sujet » est \cle{discutir}.} \esp{``Discutimos el problema de las pantallas.''}
    \item \textbf{ATTENTION FRANCOPHONE :} \esp{Discutir} a un deuxième sens très fréquent : \textbf{« se disputer, se quereller »} (souvent avec \emph{con} ou \emph{por}). \esp{``Mis padres discuten por dinero''} (Mes parents se disputent pour l'argent). \esp{``No discutas conmigo''} (Ne te dispute pas / ne me contredis pas).
    \item Pour « avoir une discussion calme », préférez \esp{hablar de} ou \esp{conversar sobre}.
    \item \textbf{Astuce Bac :} Dans l'oral, dites : \esp{``Quiero hablar de / dar mi opinión sobre el tema de las NTIC''}.
\end{itemize}
\end{attention}

\courssub{Production écrite guidée : paragraphe d'opinion (8-10 lignes)}

Consigne type Bac : \emph{« Redacta un párrafo de 8 a 10 líneas expresando tu opinión sobre el uso de las NTIC en la educación infantil. Debes incluir al menos una orden/prohibición, un ``chez'' y un verbo de ``devenir''. Usa conectores lógicos. »}

\begin{textebox}[Amorces de débat (Frases de arranque para tu redacción)]
\textbf{POSITION NUANCÉE (Recommandée) :} \\
\esp{``En mi opinión, las NTIC son una espada de doble filo. Es innegable que facilitan el acceso al conocimiento y mantienen unidos a las familias dispersas, como ocurre en casa de muchos cameruneses en el extranjero. Sin embargo, el riesgo de adicción es real: los niños se vuelven dependientes de la dopamina digital y su rendimiento escolar baja. Por eso, propongo no prohibir, sino regular: \textbf{establezcan} horarios claros y \textbf{no permitan} pantallas en el dormitorio. En conclusión, la clave está en la educación digital responsable de los padres.''} \\[6pt]
\textbf{POSITION « EN CONTRA » (Restrictive) :} \\
\esp{``Estoy firmemente en contra del móvil en la infancia. Los estudios demuestran que la exposición temprana altera el desarrollo cerebral. En casa de mis vecinos, el niño de seis años \textbf{se ha convertido en} un adicto a los videojuegos y \textbf{se queda} horas sin hablar. \textbf{Prohiban} el smartphone hasta los 14 años; \textbf{no dejen} que una pantalla críe a sus hijos. La infancia debe ser juego real, no virtual.''} \\[6pt]
\textbf{POSITION « A FAVOR » (Progressiste) :} \\
\esp{``A mi juicio, prohibir el móvil es un error pedagógico. Vivimos en la era digital: los niños \textbf{deben hacerse} competentes ya. \textbf{Entre los profesores} innovadores, la tableta es herramienta de creación, no de consumo pasivo. \textbf{No prohíban}, \textbf{acompañen}: enseñen pensamiento crítico, verificación de fuentes, programación. El futuro pertenece a los nativos digitales formados, no a los analfabetos tecnológicos.''}
\end{textebox}

\begin{exercice}[À vous de jouer : Rédaction d'opinion]
Rédigez votre paragraphe (8-10 lignes) en espagnol.
\textbf{Contraintes obligatoires :}
\begin{enumerate}
    \item Prise de position claire (\esp{En mi opinión / Creo que / Estoy a favor / en contra}).
    \item Minimum 3 arguments distincts (accès savoir, social, santé, avenir pro...).
    \item \textbf{Triple Défi :} 1 ordre/interdiction + 1 « chez » + 1 « devenir » (bien choisi).
    \item Connecteurs : \esp{ya que, por tanto, sin embargo, en conclusión}.
    \item Vocabulaire précis : \esp{rendimiento, adicción, diáspora, pacto familiar, algoritmo, ciberacoso}.
\end{enumerate}
Écrivez votre brouillon ci-dessous, puis recopiez au net sur votre copie.
\vspace{4cm}
\end{exercice}

\begin{corrige}[Exercice : Rédaction d'opinion (Modèle de correction)]
\textbf{Proposition de paragraphe type (niveau Bac 16-18/20) :} \\[4pt]
\esp{``En mi opinión, la tecnología no debe demonizarse, pero exige límites firmes. \textbf{En primer lugar}, el smartphone es una ventana al mundo: en casa de mi tío en Yaoundé, sus hijos siguen cursos en línea gratuitos y \textbf{se han hecho} bilingües español-inglés gracias a apps interactivas. \textbf{Además}, fortalece los lazos con la diáspora. \textbf{Sin embargo}, el peligro de la adicción es real: muchos chicos \textbf{se ponen} agresivos si se les quita el dispositivo y \textbf{se quedan} aislados en su burbuja digital. \textbf{Por tanto}, la solución no es la prohibición total, sino la corresponsabilidad: \textbf{los padres deben establecer} un contrato de uso escrito, \textbf{no permitan} el móvil en la hora de comer ni antes de dormir, y \textbf{den} el ejemplo desconectándose ellos mismos. \textbf{En resumen}, educar en la era digital es enseñar a usar la herramienta sin convertirse en su esclavo.''} \\[4pt]
\textbf{Analyse du « Triple Défi » :} \\
\checkmark \textbf{Ordres/Interdictions :} \esp{los padres deben establecer} (obligation), \esp{no permitan} (subjonctif interdiction formelle \emph{ustedes}), \esp{den el ejemplo} (impératif \emph{ustedes} irrégulier de \esp{dar}). \\
\checkmark \textbf{Chez :} \esp{en casa de mi tío en Yaoundé} (lieu précis, ambiance familiale). \\
\checkmark \textbf{Devenir :} \esp{se han hecho bilingües} (\cle{hacerse} + nom = devenir par effort/processus abouti), \esp{se ponen agresivos} (\cle{ponerse} + adj = réaction passagère), \esp{se quedan aislados} (\cle{quedarse} + adj = état final résiduel), \esp{convertirse en su esclavo} (\cle{convertirse en} + nom = transformation de nature).
\end{corrige}

\newpage

\coursec{Activité d'intégration}

\courssub{Objectifs de l'activité}
Mobiliser dans une situation de communication authentique les connaissances grammaticales et lexicales de la leçon : exprimer l'ordre et l'interdiction (impératif, \esp{no} + subjonctif, \esp{prohibido} / \esp{se prohíbe} + infinitif, infinitif d'affiches), employer correctement les prépositions pour dire « chez » (\esp{en casa de}, \esp{a casa de}, \esp{entre los...}) et choisir le verbe adéquat pour traduire « devenir » (\esp{ponerse}, \esp{volverse}, \esp{hacerse}, \esp{convertirse en}, \esp{llegar a ser}, \esp{quedarse}). Produire une affiche collective de règles de la maison (\esp{Reglas de la casa}) et participer à un débat argumenté sur l'impact des NTIC sur l'éducation des enfants.

\courssub{Situation}
\begin{textebox}[Situation de départ : Consejo de familia en Yaoundé]
La famille \textbf{Ngoua} est réunie \esp{en casa de la abuela} (chez la grand-mère), dans le quartier Mvog-Ada à Yaoundé. C'est le dimanche après-midi, après le repas de \emph{ndolé} et de plantains. Le sujet est grave : les trois enfants (15, 12 et 8 ans) passent trop de temps sur leurs téléphones et tablettes. Le père, \emph{Don José}, rentré d'Espagne où il a travaillé dix ans, veut imposer une discipline stricte. La mère, \emph{Doña Marie}, partage son inquiétude mais préfère le dialogue. L'adolescent, \emph{Carlos}, 15 ans, revendique son autonomie. La grand-mère, \emph{Abuela Rosa}, observe avec lucidité les changements de comportement. Un cousin, \emph{Armand}, 20 ans, étudiant en droit, fait office de secrétaire pour rédiger l'affiche officielle des \esp{Reglas de la casa} qui sera collée sur le réfrigérateur.
\end{textebox}

\begin{textebox}[Document déclencheur : Article de presse locale (adapté)]
\textbf{Los niños cameruneses y las pantallas : un reto para las familias}

En los últimos cinco años, el uso de smartphones y tabletas se ha \textbf{generalizado} en los hogares de Yaoundé y Douala. Según una encuesta reciente del Ministerio de Educación Básica, el 68\% de los alumnos de secundaria posee su propio teléfono antes de los 14 años. Los pedagogos advierten: « El exceso de pantalla \textbf{vuelve} a los niños más irritables, menos atentos en clase y \textbf{les cuesta} conciliar el sueño ». En \textbf{casa de muchos abuelos}, la tradición oral y los cuentos de la noche \textbf{han quedado} \textbf{sustituidos} por videos de TikTok y videojuegos. Las familias buscan \textbf{convertir} este desafío en una oportunidad educativa, estableciendo \textbf{límites claros} sin prohibir la tecnología, que \textbf{llegó para quedarse}.
\end{textebox}

\courssub{Consignes}
\begin{enumerate}
\item \textbf{Compréhension et préparation individuelle (10 min).} Lisez la situation et l'article. Relevez tous les verbes exprimant un changement d'état (« devenir ») et notez leur infinitif. Identifiez les trois occurrences de « chez » et leur équivalent espagnol. Repérez deux structures d'interdiction ou d'ordre.
\item \textbf{Mise en place du jeu de rôles (5 min).} Formez des groupes de 5 élèves. Attribuez les rôles : Père (Don José), Mère (Doña Marie), Adolescent (Carlos, 15 ans), Grand-mère (Abuela Rosa), Secrétaire (Armand). Chaque élève prépare 4--5 répliques types pour son personnage, en utilisant \textbf{au moins} : 2 ordres/interdictions (différentes structures), 1 « chez », 1 « devenir », 3 mots du lexique NTIC.
\item \textbf{Jeu de rôles : el Consejo de familia (15 min).} Jouez la scène en espagnol. Le secrétaire note les arguments clés et les règles décidées. La grand-mère ouvre la séance : \esp{« ¡Bienvenidos a mi casa! Hoy decidimos las reglas. »} Le père impose au moins deux interdictions formelles. La mère donne au moins deux ordres directs. L'adolescent contre-argumente en citant « chez ses amis ». La grand-mère constate deux évolutions (devenir). Le secrétaire synthétise.
\item \textbf{Rédaction de l'affiche « Reglas de la casa » (10 min).} Le secrétaire rédige l'affiche finale (format A4 paysage) avec \textbf{exactement 5 règles}, en variant les structures : (a) \esp{Prohibido} + infinitif, (b) \esp{Se prohíbe} + infinitif, (c) \esp{No} + subjonctif présent, (d) Impératif affirmatif (2e personne pluriel \esp{voseo} ou \esp{ustedes}), (e) Infinitif d'affiche. Chaque règle est illustrée d'un petit pictogramme dessiné à la main.
\item \textbf{Présentation et débat plénier (15 min).} Chaque groupe affiche son panneau au tableau. Lecture à voix haute par le secrétaire. Débat guidé par le professeur : \esp{« ¿Son las NTIC un peligro para la educación de los niños? »} / \esp{« ¿Qué papel deben jugar los abuelos? »} / \esp{« ¿Es mejor prohibir o educar en el uso responsable? »} Les élèves doivent utiliser \esp{en mi opinión}, \esp{creo que} + subjonctif/indicatif, \esp{es necesario que} + subjonctif, \esp{debería} + infinitif.
\end{enumerate}

\courssub{Ressources à mobiliser}
\begin{definition}[Ordres et interdictions -- synthèse]
\textbf{1. Impératif (ordre direct, familiarité ou autorité) :} \esp{¡Apaga la tableta!} (tú), \esp{¡Apagad el móvil!} (vosotros), \esp{¡Apaguen los dispositivos!} (ustedes). Négatif : \esp{¡No enciendas la tele!}, \esp{¡No habléis en la mesa!}, \esp{¡No toquen la pantalla!} (subjonctif).\\
\textbf{2. \esp{No} + subjonctif présent (interdiction formelle, impersonnelle ou adressée à tous) :} \esp{No fuméis aquí.} $\rightarrow$ \esp{No fumen aquí.} \esp{No use el móvil en clase.} \\
\textbf{3. \esp{Prohibido} + infinitif (affiches, pancartes, règlements écrits) :} \esp{Prohibido fumar.} \esp{Prohibido el paso.} \esp{Prohibido usar el móvil durante la cena.} \\
\textbf{4. \esp{Se prohíbe} + infinitif (interdiction institutionnelle, passive réflexe) :} \esp{Se prohíbe el uso de tabletas en la cama.} \esp{Se prohíbe hablar con la boca llena.} \\
\textbf{5. Infinitif d'affiche / mode d'emploi (neutre, impersonnel) :} \esp{No tocar.} \esp{No molestar.} \esp{Silencio, por favor.}
\end{definition}

\begin{definition}[« Chez » -- prépositions de lieu et d'appartenance]
\begin{center}
\begin{tabularx}{\linewidth}{|X|X|X|}
\hline
\rowcolor{popblueL} Contexte & Espagnol & Exemple \\
\hline
Chez [personne] (lieu, statique) & \textbf{en casa de} + nom/pronom & \esp{Estamos en casa de la abuela.} \\
Chez [personne] (mouvement, aller) & \textbf{a casa de} + nom/pronom & \esp{Vamos a casa de mis amigos.} \\
Chez [professionnel / commerce] & \textbf{en el / la} + nom & \esp{En el médico, en la peluquería.} \\
Chez [peuple, groupe, culture] & \textbf{entre los / las} + nom & \esp{Entre los españoles, la sobremesa es sagrada.} \esp{Entre los cameruneses, el respeto a los mayores es fundamental.} \\
\hline
\end{tabularx}
\end{center}
\emph{Attention :} \esp{en mi casa} (chez moi), \esp{en tu casa} (chez toi), \esp{en su casa} (chez lui/elle/vous/ils) -- pas de \esp{casa de} avec pronoms possessifs.
\end{definition}

\begin{definition}[« Devenir » -- six verbes de changement, nuances et registres]
\begin{center}
\begin{tabularx}{\linewidth}{|X|X|X|}
\hline
\rowcolor{poppurpleL} Verbe & Nuance / Emploi & Exemple traduit \\
\hline
\textbf{ponerse} + adj. & Changement \textbf{soudain, temporaire, émotionnel, physique} (couleur, humeur, état) & \esp{Se puso rojo de ira.} (Il est devenu rouge de colère.) \esp{Los niños se ponen nerviosos sin wifi.} \\
\textbf{volverse} + adj. & Changement de \textbf{caractère, personnalité, comportement} (souvent négatif, progressif) & \esp{Se ha vuelto impaciente.} (Il est devenu impatient.) \esp{La tecnología vuelve a los jóvenes adictos.} \\
\textbf{hacerse} + nom/adj. & Changement \textbf{volontaire, résultat d'un processus, profession, idéologie, religion} & \esp{Se hizo médico.} (Il est devenu médecin.) \esp{Se ha hecho adicto a los videojuegos.} \\
\textbf{convertirse en} + nom & Transformation \textbf{radicale, spectaculaire, souvent inattendue} (métamorphose) & \esp{La tablet se convirtió en su mejor amigo.} (La tablette est devenue son meilleur ami.) \\
\textbf{llegar a ser} + nom/adj. & Aboutissement \textbf{après un long parcours, effort, temps} (réalisation) & \esp{Llegó a ser un gran programador.} (Il est devenu / a fini par devenir un grand programmeur.) \\
\textbf{quedarse} + adj. / part. & État \textbf{résultant, définitif ou involontaire} (souvent négatif : handicap, solitude, ruine) & \esp{Se quedó ciego.} (Il est devenu aveugle.) \esp{La tradición se ha quedado obsoleta.} \\
\hline
\end{tabularx}
\end{center}
\end{definition}

\begin{definition}[Lexique NTIC et éducation -- rappel]
\textbf{Dispositifs :} \esp{el smartphone / el móvil}, \esp{la tableta}, \esp{la consola de videojuegos}, \esp{el portátil}, \esp{la pantalla}. \textbf{Applications / contenus :} \esp{las redes sociales (Instagram, TikTok, WhatsApp)}, \esp{los videojuegos}, \esp{los dibujos animados}, \esp{los influencers}, \esp{el algoritmo}. \textbf{Verbes d'action :} \esp{navegar}, \esp{chatear}, \esp{subir / bajar (fotos/videos)}, \esp{compartir}, \esp{seguir (a alguien)}, \esp{volverse viral}. \textbf{Noms abstraits :} \esp{la adicción (a las pantallas)}, \esp{el ciberacoso (cyberharcèlement)}, \esp{la privacidad}, \esp{el control parental}, \esp{el tiempo de pantalla}, \esp{la desconexión digital}. \textbf{Adjectifs :} \esp{conectado / desconectado}, \esp{virtual / real}, \esp{adictivo}, \esp{educativo / perjudicial}.
\end{definition}

\begin{methode}[Comment choisir le bon verbe pour « devenir »]
\begin{enumerate}
\item \textbf{Identifiez la nature du changement :} soudain/émotionnel $\rightarrow$ \esp{ponerse} ; caractère/comportement $\rightarrow$ \esp{volverse} ; volontaire/profession $\rightarrow$ \esp{hacerse} ; métamorphose radicale $\rightarrow$ \esp{convertirse en} ; long parcours/réussite $\rightarrow$ \esp{llegar a ser} ; état resultant/involontaire $\rightarrow$ \esp{quedarse}.
\item \textbf{Vérifiez la construction :} \esp{ponerse/volverse/quedarse} + adjectif (accord) ; \esp{hacerse} + nom/adjectif ; \esp{convertirse en/llegar a ser} + nom/adjectif (sans \esp{en} pour \esp{llegar a ser} + adj).
\item \textbf{Attention aux pronoms :} \esp{ponerse, volverse, hacerse, convertirse, quedarse} sont \textbf{pronominaux} (se, me, te, nos, os). \esp{Llegar a ser} ne l'est pas.
\end{enumerate}
\end{methode}

\begin{methode}[Exprimer l'interdiction selon le contexte]
\begin{enumerate}
\item \textbf{Affiche / Panneau public :} \esp{Prohibido} + infinitif (\esp{Prohibido fumar}) ou \esp{Se prohíbe} + infinitif (\esp{Se prohíbe la entrada}).
\item \textbf{Règle familiale, scolaire, adressée à un groupe connu :} \esp{No} + subjonctif (\esp{No uséis el móvil}) ou Impératif négatif (\esp{No toquéis la pantalla}).
\item \textbf{Ordre direct, autorité parentale, urgence :} Impératif affirmatif (\esp{Apagad la tableta}, \esp{Apaguen la tableta}).
\item \textbf{Instruction de procédure, mode d'emploi :} Infinitif (\esp{No molestar}, \esp{Pedir permiso}).
\end{enumerate}
\end{methode}

\begin{attention}[Erreurs fréquentes des francophones -- Pièges à éviter]
\begin{itemize}
\item \textbf{« Chez » + pronom :} \emph{*en casa de mí} $\rightarrow$ \esp{en mi casa} ; \emph{*a casa de ti} $\rightarrow$ \esp{a tu casa}.
\item \textbf{« Devenir » unique :} \emph{*Los niños se devenen nerviosos} $\rightarrow$ \esp{Los niños se ponen nerviosos} (soudain) / \esp{se vuelven nerviosos} (caractère).
\item \textbf{\esp{Volverse} vs \esp{Hacerse} :} \esp{Se ha vuelto egoísta} (involontaire, négatif) $\neq$ \esp{Se ha hecho vegetariano} (choix volontaire).
\item \textbf{Impératif négatif :} \emph{*No hablad} $\rightarrow$ \esp{No habléis} (vosotros) / \esp{No hablen} (ustedes). \textbf{Jamais d'impératif affirmatif pour la négation.}
\item \textbf{\esp{Prohibido} vs \esp{Prohíbe} :} \esp{Prohibido} (invariable, affiche) $\neq$ \esp{Se prohíbe} (verbe conjugué, 3e pers. sg).
\item \textbf{Faux ami \esp{Realizar} :} \esp{Realizar} = accomplir, effectuer (une tâche) ; \textbf{pas} « réaliser » (comprendre) $\rightarrow$ \esp{Darse cuenta de}.
\item \textbf{Accents toniques :} \esp{qué, cómo, cuándo, dónde, quién, por qué} (interrogatifs/exclamatifs) vs \esp{que, como, cuando, donde, quien, porque} (relatifs/conjonctions).
\end{itemize}
\end{attention}

\courssub{Proposition de production et corrigé commenté}
\begin{exoresolu}[Consejo de familia -- Production modèle du groupe 3]
\textbf{Étape 1 -- Extraits du jeu de rôles (répliques types)}

\emph{Abuela Rosa :} \esp{« ¡Bienvenidos a mi casa, nietos! Hoy hablamos de las pantallas. En mi época, no había móviles, solo cuentos bajo el mango. Ahora veo que mis nietos \textbf{se han vuelto} impacientes y \textbf{se ponen} nerviosos si no hay wifi. La tradición \textbf{se ha quedado} olvidada. »}

\emph{Don José (père) :} \esp{« \textbf{Se prohíbe} el uso del móvil durante las comidas. \textbf{Prohibido} traer la tableta al dormitorio después de las 21 h. \textbf{No uséis} el teléfono mientras hacéis los deberes. »}

\emph{Doña Marie (mère) :} \esp{« ¡\textbf{Apagad} los dispositivos a las 21:30! ¡\textbf{No habléis} con ese tono, Carlos! \textbf{Debéis} pedir permiso antes de descargar aplicaciones. »}

\emph{Carlos (adolescent, 15 ans) :} \esp{« \textbf{En casa de mis amigos}, no hay tantas reglas. \textbf{En casa de Pablo}, cenan viendo series. Yo \textbf{me he hecho} responsable: controlo mi tiempo, \textbf{voy a llegar a ser} autónomo. No soy un niño. »}

\emph{Armand (secrétaire) :} \esp{« Apunto las propuestas. ¿Acordamos cinco reglas claras para el cartel? »}

\textbf{Étape 2 -- Affiche « Reglas de la casa Ngoua » (version finale corrigée)}

\begin{center}
\begin{tabularx}{\linewidth}{|c|X|}
\hline
\rowcolor{popblueL} \textbf{N°} & \textbf{Regla} \\
\hline
1 & \esp{\textbf{PROHIBIDO} el uso de móviles y tabletas durante las comidas familiares.} \\
\hline
2 & \esp{\textbf{SE PROHÍBE} llevar dispositivos a los dormitorios después de las 21:30.} \\
\hline
3 & \esp{\textbf{NO USÉIS} el teléfono mientras hacéis los deberes ni en la sobremesa.} \\
\hline
4 & \esp{\textbf{APAGUEN} todas las pantallas a las 22:00 en punto. ¡A dormir!} \\
\hline
5 & \esp{\textbf{PEDIR PERMISO} a los padres antes de descargar nuevas aplicaciones o juegos.} \\
\hline
\end{tabularx}
\end{center}

\textbf{Étape 3 -- Intervention de Carlos lors du débat plénier}
\esp{« En mi opinión, las NTIC no son un peligro en sí mismas, \textbf{se convierten en} un problema si no hay educación digital. \textbf{Es necesario que} los padres \textbf{enseñen} a gestionar el tiempo, \textbf{no solo prohíban, sino que también orienten}. \textbf{Deberíamos} tener clases de ciudadanía digital en el colegio. \textbf{Entre los jóvenes cameruneses}, el móvil es también una herramienta para estudiar, no solo para TikTok. »}

\tcblower
\textbf{Commentaire détaillé des choix linguistiques}

\begin{enumerate}
\item \textbf{Ordres et interdictions -- variété et justesse des structures :}
\begin{itemize}
\item Règle 1 : \esp{PROHIBIDO el uso...} -- Infinitif d'affiche nominalisé (\emph{el uso}), registre administratif, impersonnel, idéal pour un panneau affiché sur le frigo.
\item Règle 2 : \esp{SE PROHÍBE llevar...} -- Passif réflexe (\esp{se prohíbe} + infinitif), valeur normative générale, « il est interdit de... », plus solennel que \esp{prohibido} seul.
\item Règle 3 : \esp{NO USÉIS...} -- Subjonctif présent (2e pers. plur. \esp{voseo} \esp{uséis} de \esp{usar}) après \esp{no}, ordre négatif adressé aux enfants (famille, tutoiement). \emph{Nota :} en Espagne péninsulaire \esp{uséis}, en Amérique et souvent au Cameroun hispanophone \esp{usen} (ustedes). Les deux acceptés.
\item Règle 4 : \esp{APAGUEN...} -- Impératif affirmatif 2e pers. plur. \esp{ustedes} (\esp{apaguen} de \esp{apagar}), ordre direct, autorité parentale, ton ferme mais affectueux (\esp{¡A dormir!}).
\item Règle 5 : \esp{PEDIR PERMISO...} -- Infinitif d'affiche / mode d'emploi, instruction de procédure, neutre, s'adresse à tous sans marquer la personne.
\end{itemize}
\emph{Point de vigilance :} L'alternance \esp{voseo} (\esp{apagad, uséis}) / \esp{ustedes} (\esp{apaguen, usen}) est cohérente si la famille utilise le \esp{voseo} (modèle équatorial/rioplatense) ou le \esp{ustedes} généralisé. Au Cameroun, l'espagnol enseigné suit souvent la norme \esp{ustedes} pour le pluriel ; l'enseignant validera l'un ou l'autre tant que la cohérence interne est respectée.

\item \textbf{« Chez » -- quatre occurrences correctes :}
\begin{itemize}
\item \esp{en mi casa} (Abuela) -- lieu statique, pronom possessif, pas de \esp{casa de}.
\item \esp{En casa de mis amigos} (Carlos) -- lieu statique chez une tierce personne, nom propre, \esp{en casa de} obligatoire.
\item \esp{En casa de Pablo} (Carlos) -- lieu statique (habitude), \esp{en casa de} + nom propre (correction de \emph{*A casa de} qui indiquerait le mouvement).
\item \esp{Entre los jóvenes cameruneses} (Carlos, débat) -- « chez les jeunes camerounais » = au sein de ce groupe, \esp{entre los} + adjectif substantivé.
\end{itemize}

\item \textbf{« Devenir » -- six verbes distincts employés à bon escient :}
\begin{itemize}
\item \esp{se han vuelto impacientes} (\esp{volverse} + adj.) -- changement de caractère progressif, constat de la grand-mère sur le long terme, connotation négative.
\item \esp{se ponen nerviosos} (\esp{ponerse} + adj.) -- réaction émotionnelle soudaine, temporaire, déclenchée par l'absence de wifi.
\item \esp{se ha quedado olvidada} (\esp{quedarse} + part. passé comme adj.) -- état résultant définitif, la tradition \emph{a fini par rester} oubliée, passif de résultat.
\item \esp{me he hecho responsable} (\esp{hacerse} + adj.) -- processus volontaire, effort personnel de l'adolescent pour acquérir une qualité.
\item \esp{voy a llegar a ser autónomo} (\esp{llegar a ser} + adj.) -- aboutissement futur après un cheminement, maturité visée.
\item \esp{se convierten en un problema} (\esp{convertirse en} + nom) -- transformation radicale de la nature des NTIC (outils $\rightarrow$ problème) sous condition (\esp{si no hay educación}).
\end{itemize}

\item \textbf{Subjonctif dans le débat :} \esp{es necesario que los padres enseñen} (impersonnel de nécessité + subjonctif), \esp{no solo prohíban, sino que también orienten} (négation restrictive + subjonctif). \esp{Deberíamos tener} (conditionnel de politesse / conseil).
\end{enumerate}
\end{exoresolu}

\courssub{Bilan de l'activité}
Cette activité d'intégration a permis de vérifier l'acquisition des quatre piliers de la leçon dans une tâche unique, contextualisée et motivante :
\begin{itemize}
\item \textbf{Grammaire :} Les cinq structures d'ordre/interdiction ont été produites spontanément dans l'affiche, prouvant que les élèves savent choisir le registre adapté (familial, administratif, direct, impersonnel). L'alternance \esp{voseo} / \esp{ustedes} a été gérée sans erreur de morphologie.
\item \textbf{Prépositions « chez » :} Les quatre configurations (statique \esp{en casa de}, mouvement \esp{a casa de}, possessif \esp{en mi casa}, groupe \esp{entre los}) sont apparues naturellement dans le dialogue et le débat.
\item \textbf{Verbes de changement :} Les six verbes (\esp{ponerse, volverse, hacerse, convertirse en, llegar a ser, quedarse}) ont été employés avec leurs nuances respectives (soudain/émotionnel, caractère, volontaire, métamorphose, aboutissement, état resultant), sans confusion ni calque du français « devenir » unique.
\item \textbf{Lexique NTIC :} Une quinzaine de termes spécialisés ont été réinvestis (\esp{adicción, control parental, tiempo de pantalla, ciberacoso, algoritmo, influencers...}) dans des énoncés argumentatifs cohérents.
\item \textbf{Compétences transversales :} Prise de parole en interaction (jeu de rôles), production écrite collaborative (affiche), compréhension orale (écoute des autres groupes), expression orale en continu (débat), esprit critique (argumentation équilibrée sur les NTIC).
\end{itemize}
\begin{aretenir}[Critères de réussite pour l'évaluation]
\begin{itemize}
\item L'affiche comporte \textbf{exactement 5 règles} avec \textbf{5 structures distinctes} (pas de doublet).
\item Chaque élève a produit \textbf{au moins 4 répliques} exploitables pendant le jeu de rôles.
\item Le débat final montre l'usage d'\textbf{au moins 3 connecteurs logiques} (\esp{por tanto, sin embargo, además, en cambio, en conclusión}) et de \textbf{2 subjonctifs déclenchés par l'opinion/nécessité}.
\item Aucune erreur de base sur \esp{ser/estar}, accents toniques, \esp{por/para}, faux amis (\esp{realizar} $\neq$ réaliser).
\end{itemize}
\end{aretenir}

\begin{savaistu}[Le Conseil de famille au Cameroun et en Guinée équatoriale]
Dans la culture bassa, bamiléké ou fang, la réunion familiale élargie (\emph{conseil de famille}, \esp{reunión familiar}) est l'instance traditionnelle de régulation des conflits et de transmission des normes. On y siège \esp{en casa del patriarca} ou \esp{de la matriarca}. La grand-mère (\esp{abuela}, \esp{mamá grande}) y joue un rôle d'arbitre moral, gardienne de la \esp{tradición oral}. Aujourd'hui, ce cadre rituel doit intégrer les nouveaux défis : le \esp{ciberacoso} (harcèlement en ligne), la \esp{nomofobia} (peur d'être sans mobile), l'exposition précoce à des contenus inadaptés. L'Espagne et l'Amérique latine connaissent le même phénomène : la \esp{ley de protección al menor} (Espagne, 2021) oblige les plateformes à vérifier l'âge. Au Cameroun, la loi n°2010/012 sur la cybercriminalité et le décret 2023/001 sur la protection des données personnelles offrent un cadre légal, mais l'éducation reste la première ligne de défense -- d'où l'intérêt de cette activité qui ancre la grammaire dans une réalité vécue.
\end{savaistu}

\newpage

\coursec{Méthodes et exercices résolus}

\courssub{Savoir-faire 1 : exprimer l'ordre et l'interdiction}

\begin{methode}[Comment traduire un ordre ou une interdiction en espagnol]
\begin{enumerate}
\item \textbf{Identifier le destinataire.} À qui s'adresse l'ordre ? À une personne qu'on tutoie (\esp{tú}), qu'on vouvoie (\esp{usted}), à un groupe (\esp{vosotros}, \esp{ustedes}) ? Cela détermine la forme verbale.
\item \textbf{Ordre affirmé : utiliser l'impératif.} \esp{¡Ven aquí!} (Viens ici !) ; \esp{Hagan silencio.} (Faites silence.) Rappel : l'impératif affirmé de \esp{tú} reprend souvent la 3\textsuperscript{e} personne du singulier du présent (\esp{habla}, \esp{come}, \esp{escribe}), avec des irréguliers à connaître par cœur : \esp{di, haz, ve, pon, sal, sé, ten, ven}.
\item \textbf{Ordre nié (interdiction adressée à quelqu'un) : subjonctif présent + \esp{no}.} \esp{No fumes.} (Ne fume pas.) ; \esp{No toquen eso.} (Ne touchez pas à cela.) C'est la grande différence avec le français : l'espagnol n'a pas d'impératif négatif, il utilise le subjonctif.
\item \textbf{Interdiction impersonnelle (panneaux, règlements) : trois formules.}
\begin{itemize}
\item Infinitif : \esp{Prohibido fumar.} / \esp{No aparcar.} (Défense de fumer / de stationner.)
\item \esp{Se prohíbe} + infinitif ou nom : \esp{Se prohíbe el paso.} (Passage interdit.)
\item \esp{No se permite} + infinitif : \esp{No se permite usar el móvil en clase.}
\end{itemize}
\item \textbf{Ordre rapporté au discours indirect :} verbe d'ordre (\esp{decir, mandar, ordenar, prohibir}) + \esp{que} + \textbf{subjonctif}. \esp{La madre le dice que apague la tableta.} (La mère lui dit d'éteindre la tablette.)
\item \textbf{Vérifier} : accent sur l'impératif avec pronom accolé (\esp{dime}, \esp{hazlo}, \esp{levántate}) et cohérence du registre (familier / officiel).
\end{enumerate}
\end{methode}

\begin{exoresolu}[Thème : ordres et interdictions]
Traduisez en espagnol :
\begin{enumerate}
\item Ne fume pas ici, c'est interdit.
\item Éteignez vos portables avant d'entrer en classe. (vous = un groupe)
\item Défense d'utiliser les téléphones pendant l'examen. (formule de panneau)
\item Le professeur ordonne aux élèves de ranger leurs affaires.
\item Lève-toi et fais tes devoirs tout de suite !
\end{enumerate}
\tcblower
\textbf{Raisonnement et solutions :}
\begin{enumerate}
\item Interdiction adressée à \esp{tú} $\rightarrow$ \esp{no} + subjonctif : \esp{No fumes aquí, está prohibido.} « Fumer » = \esp{fumar} ; subjonctif : \esp{fumes}. On peut ajouter \esp{está prohibido} ou \esp{se prohíbe}.
\item Ordre affirmé à un groupe vouvoyé $\rightarrow$ impératif \esp{ustedes} : \esp{Apaguen sus móviles antes de entrar en clase.} Attention : \esp{apagar} a une alternance orthographique \esp{g $\rightarrow$ gu} devant \esp{e} : \esp{apaguen}. « Avant de » + infinitif = \esp{antes de} + infinitif.
\item Formule impersonnelle : \esp{Prohibido usar los teléfonos durante el examen.} ou \esp{Se prohíbe el uso de los teléfonos durante el examen.} Ne jamais dire \esp{*es prohibido usar} sur un panneau : calque du français.
\item Ordre rapporté $\rightarrow$ \esp{ordenar que} + subjonctif : \esp{El profesor ordena a los alumnos que guarden sus cosas.} « Ranger » = \esp{guardar} ; subjonctif : \esp{guarden}.
\item Deux impératifs affirmés \esp{tú} : \esp{¡Levántate y haz tus deberes ahora mismo!} \esp{Levántate} : pronom accolé + accent sur l'antépénultième ; \esp{haz} : impératif irrégulier de \esp{hacer} ; « tout de suite » = \esp{ahora mismo} (et non \esp{*todo seguido}).
\end{enumerate}
\textbf{Conclusion :} pour l'interdiction, le réflexe est : destinataire précis $\rightarrow$ \esp{no} + subjonctif ; interdiction générale $\rightarrow$ \esp{prohibido} + infinitif ou \esp{se prohíbe}.
\end{exoresolu}

\courssub{Savoir-faire 2 : traduire « chez »}

\begin{methode}[Choisir le bon équivalent de « chez »]
\begin{enumerate}
\item \textbf{Se demander d'abord : lieu ou mouvement ?}
\begin{itemize}
\item Lieu (être, rester) : \esp{en casa de} + personne. \esp{Ceno en casa de mis tíos.} (Je dîne chez mes oncles.)
\item Mouvement (aller, venir) : \esp{a casa de}. \esp{Voy a casa de Pedro.} (Je vais chez Pedro.)
\end{itemize}
\item \textbf{« Chez » + professionnel ou commerce :} article + nom du lieu ou du professionnel : \esp{en el médico} (chez le médecin), \esp{a la peluquería} (chez le coiffeur), \esp{en casa del dentista} est aussi possible.
\item \textbf{« Chez » au sens de « parmi, dans la mentalité de » :} \esp{entre}. \esp{Entre los españoles, la siesta es tradicional.} (Chez les Espagnols, la sieste est traditionnelle.) \esp{Es común entre los jóvenes cameruneses.} (C'est courant chez les jeunes Camerounais.)
\item \textbf{« Chez soi » sans précision :} \esp{en casa} / \esp{a casa}. \esp{Me quedo en casa.} (Je reste chez moi.) \esp{Vuelve a casa.} (Il rentre chez lui.)
\item \textbf{Piège à éviter :} ne jamais traduire « chez » par \esp{*chez} ni par \esp{*en la casa de} quand on parle d'une personne sans article possessif : on dit \esp{en casa de Ana}, mais \esp{en la casa de mis padres} si l'on insiste sur la maison elle-même.
\end{enumerate}
\end{methode}

\begin{exoresolu}[Thème : « chez »]
Traduisez en espagnol :
\begin{enumerate}
\item Samedi, nous irons chez nos grands-parents à Yaoundé.
\item Chez les adolescents, le téléphone portable est devenu indispensable.
\item Ma sœur est restée chez elle toute la journée.
\item Tu dois aller chez le médecin si tu es malade.
\item Les enfants passent la nuit chez leur cousine.
\end{enumerate}
\tcblower
\textbf{Raisonnement et solutions :}
\begin{enumerate}
\item Mouvement (\esp{ir}) $\rightarrow$ \esp{a casa de} : \esp{El sábado iremos a casa de nuestros abuelos en Yaundé.} Futur simple de \esp{ir} : \esp{iremos}. « Yaoundé » se dit \esp{Yaundé} en espagnol.
\item « Chez » = « parmi, dans le groupe de » $\rightarrow$ \esp{entre} : \esp{Entre los adolescentes, el teléfono móvil se ha vuelto indispensable.} Noter l'accord : \esp{entre} ne rend pas le nom sujet singulier ; ici le sujet est \esp{el teléfono móvil}.
\item « Chez elle » sans précision $\rightarrow$ \esp{en casa} (ou \esp{en su casa}) : \esp{Mi hermana se quedó en casa todo el día.} Passé simple : \esp{se quedó}.
\item Professionnel $\rightarrow$ \esp{al médico} : \esp{Debes ir al médico si estás enfermo.} Contraction obligatoire \esp{a + el = al}. « Malade » (état) = \esp{estar enfermo}.
\item Lieu (passer la nuit) $\rightarrow$ \esp{en casa de} : \esp{Los niños pasan la noche en casa de su prima.} « Leur cousine » : \esp{su prima} (possessif \esp{su}, pas d'article).
\end{enumerate}
\textbf{Conclusion :} le trio à retenir : \esp{en casa de} (lieu), \esp{a casa de} (mouvement), \esp{entre} (chez un peuple, un groupe).
\end{exoresolu}

\courssub{Savoir-faire 3 : traduire « devenir »}

\begin{methode}[Choisir entre ponerse, volverse, hacerse, quedarse, convertirse en, llegar a ser]
\begin{enumerate}
\item \textbf{Changement passager, physique ou émotionnel (souvent involontaire) : \esp{ponerse} + adjectif.} \esp{Se puso rojo.} (Il est devenu rouge.) \esp{La niña se pone triste.} (La fillette devient triste.)
\item \textbf{Changement profond, durable, souvent négatif ou inattendu : \esp{volverse} + adjectif.} \esp{Se ha vuelto muy nervioso.} (Il est devenu très nerveux.) \esp{El niño se volvió insoportable con la tableta.}
\item \textbf{Changement volontaire, professionnel, idéologique, résultat d'un effort : \esp{hacerse} + nom ou adjectif.} \esp{Se hizo médico.} (Il est devenu médecin.) \esp{Se hicieron amigos.} (Ils sont devenus amis.)
\item \textbf{Résultat d'un processus long, évolution aboutie : \esp{llegar a ser}.} \esp{Llegó a ser director del colegio.} (Il est devenu / finit par devenir directeur du collège.)
\item \textbf{Transformation complète : \esp{convertirse en} + nom.} \esp{El móvil se ha convertido en una herramienta educativa.} (Le portable est devenu un outil éducatif.)
\item \textbf{Changement définitif et souvent brutal (handicap, état final) : \esp{quedarse} + adjectif.} \esp{Se quedó ciego.} (Il est devenu aveugle.) \esp{Se quedó sola.}
\item \textbf{Vérifier} : le verbe choisi se conjugue toujours avec le pronom réfléchi (\esp{me pongo, te vuelves, se hace...}) et l'adjectif s'accorde avec le sujet.
\end{enumerate}
\end{methode}

\begin{exoresolu}[Thème : « devenir »]
Traduisez en espagnol :
\begin{enumerate}
\item Quand il a vu la note, il est devenu rouge de colère.
\item Après des années d'efforts, elle est devenue professeure d'espagnol.
\item Les écrans sont devenus un problème dans beaucoup de familles.
\item Mon petit frère est devenu insupportable depuis qu'il a une tablette.
\item À force de travail, ce jeune de Douala est devenu un grand footballeur.
\end{enumerate}
\tcblower
\textbf{Raisonnement et solutions :}
\begin{enumerate}
\item Réaction physique passagère $\rightarrow$ \esp{ponerse} : \esp{Cuando vio la nota, se puso rojo de ira.} Passé simple : \esp{vio}, \esp{se puso}. « De colère » = \esp{de ira} / \esp{de rabia}.
\item Évolution longue et aboutie $\rightarrow$ \esp{llegar a ser} : \esp{Después de años de esfuerzo, llegó a ser profesora de español.} \esp{Hacerse profesora} est aussi correct (choix volontaire).
\item Transformation complète $\rightarrow$ \esp{convertirse en} : \esp{Las pantallas se han convertido en un problema en muchas familias.} Parfait espagnol \esp{se han convertido} pour un changement constaté aujourd'hui.
\item Changement durable et négatif $\rightarrow$ \esp{volverse} : \esp{Mi hermanito se ha vuelto insoportable desde que tiene una tableta.} « Depuis que » + présent = \esp{desde que} + présent.
\item Résultat d'un effort $\rightarrow$ \esp{hacerse} ou \esp{llegar a ser} : \esp{A fuerza de trabajo, este joven de Duala se hizo un gran futbolista.} « À force de » = \esp{a fuerza de} ; « Douala » = \esp{Duala}.
\end{enumerate}
\textbf{Conclusion :} retenez la carte mentale : passager $\rightarrow$ \esp{ponerse} ; durable/négatif $\rightarrow$ \esp{volverse} ; volontaire $\rightarrow$ \esp{hacerse} ; aboutissement $\rightarrow$ \esp{llegar a ser} ; transformation $\rightarrow$ \esp{convertirse en} ; brutal/définitif $\rightarrow$ \esp{quedarse}.
\end{exoresolu}

\courssub{Savoir-faire 4 : argumenter dans un débat (éducation des enfants et NTIC)}

\begin{methode}[Structurer une prise de parole ou un paragraphe d'argumentation]
\begin{enumerate}
\item \textbf{Prendre position clairement :} \esp{En mi opinión...}, \esp{Estoy a favor de...}, \esp{Me opongo a...}, \esp{Desde mi punto de vista...}
\item \textbf{Annoncer le plan :} \esp{En primer lugar... ; además... ; por otro lado... ; sin embargo... ; por último...}
\item \textbf{Donner des arguments avec des exemples concrets :} \esp{Por ejemplo, en muchos colegios de Camerún los alumnos usan las tabletas para investigar.} Varier : argument + exemple + conséquence.
\item \textbf{Concéder puis réfuter :} \esp{Es cierto que las pantallas distraen, pero también es verdad que enseñan.} Le subjonctif après \esp{es posible que}, \esp{no creo que} : \esp{No creo que el móvil sea un enemigo.}
\item \textbf{Conclure :} \esp{En conclusión...}, \esp{Para terminar...}, \esp{En definitiva, la clave es el equilibrio.}
\item \textbf{Lexique utile :} \esp{la educación, la pantalla, el uso excesivo, la adicción, el control parental, aprender jugando, el rendimiento escolar, distraer, fomentar, limitar el tiempo, los contenidos inadecuados.}
\end{enumerate}
\end{methode}

\begin{exoresolu}[Rédaction : paragraphe d'opinion]
Rédigez en espagnol un paragraphe de 8 à 10 lignes répondant à la question : \esp{¿Deben los padres prohibir el móvil a sus hijos?} Vous prendrez position, donnerez deux arguments, une concession et une conclusion.
\tcblower
\textbf{Démarche :} (1) position : non, il ne faut pas interdire mais encadrer ; (2) argument 1 : le portable est un outil d'apprentissage ; (3) argument 2 : l'interdiction totale pousse à l'usage caché ; (4) concession : oui, l'usage excessif nuit aux études ; (5) conclusion : l'équilibre et le contrôle parental.

\textbf{Réponse rédigée :}

\esp{En mi opinión, los padres no deben prohibir el móvil a sus hijos, sino enseñarles a usarlo bien. En primer lugar, el teléfono se ha convertido en una herramienta educativa: por ejemplo, muchos alumnos de Yaoundé consultan diccionarios y vídeos para hacer sus deberes. Además, una prohibición total empuja a los jóvenes a usarlo a escondidas, lo que es peor. Es cierto que el uso excesivo distrae y baja el rendimiento escolar; sin embargo, no creo que la solución sea prohibir. En conclusión, la clave es el equilibrio: los padres deben limitar el tiempo de pantalla y controlar los contenidos, porque educar es acompañar, no prohibir.}

\textbf{Justifications grammaticales :} \esp{no deben prohibir... sino enseñarles} : \esp{sino} après une négation pour opposer deux idées ; pronom \esp{les} accolé à l'infinitif. \esp{se ha convertido en} : « devenir » (transformation). \esp{a escondidas} : « en cachette ». \esp{no creo que la solución sea} : subjonctif après \esp{no creo que}. \esp{educar es acompañar} : infinitif sujet, construction typique de la conclusion.
\end{exoresolu}

\courssub{Savoir-faire 5 : version — repérer ordre, « chez » et « devenir » dans un texte espagnol}

\begin{methode}[Traduire de l'espagnol vers le français sans calquer]
\begin{enumerate}
\item \textbf{Repérer les structures avant de traduire :} souligner \esp{no} + subjonctif (interdiction), \esp{prohibido / se prohíbe} (interdiction impersonnelle), \esp{en casa de / a casa de / entre} (« chez »), et les verbes de changement (\esp{ponerse, volverse, hacerse...}).
\item \textbf{Traduire le sens, pas les mots :} \esp{Entre los jóvenes} = « chez les jeunes » et non « entre les jeunes » si le sens est « dans la jeunesse » ; \esp{se puso a llorar} = « il se mit à pleurer » (construction \esp{ponerse a} + infinitif = se mettre à).
\item \textbf{Rendre l'impératif négatif par un impératif français :} \esp{No dejes el móvil en la mesa.} = « Ne laisse pas le portable sur la table. »
\item \textbf{Choisir le bon français pour les verbes de changement :} \esp{volverse loco} = « devenir fou » ; \esp{hacerse mayor} = « devenir grand / vieillir » ; \esp{quedarse dormido} = « s'endormir » (et non « devenir dormi »).
\item \textbf{Relire en français :} la phrase traduite doit être naturelle ; si elle sonne espagnol, c'est un calque.
\end{enumerate}
\end{methode}

\begin{exoresolu}[Version : extrait sur les écrans]
Traduisez en français :

\esp{«No dejes que el niño pase horas delante de la pantalla», se prohíbe en muchas familias el uso del móvil durante la cena. Entre los padres españoles, la preocupación crece: algunos se han vuelto muy estrictos y otros, en casa de mis vecinos por ejemplo, han llegado a ser verdaderos expertos en control parental.}

\tcblower
\textbf{Analyse des structures :}
\begin{itemize}
\item \esp{No dejes que... pase} : interdiction (\esp{no} + subjonctif) + subordonnée au subjonctif après \esp{dejar que}.
\item \esp{se prohíbe... el uso} : interdiction impersonnelle $\rightarrow$ « l'usage... est interdit ».
\item \esp{Entre los padres españoles} : « chez » au sens de « parmi » $\rightarrow$ « chez les parents espagnols ».
\item \esp{se han vuelto muy estrictos} : changement durable $\rightarrow$ « sont devenus très stricts ».
\item \esp{en casa de mis vecinos} : « chez » + lieu $\rightarrow$ « chez mes voisins ».
\item \esp{han llegado a ser} : aboutissement $\rightarrow$ « sont devenus / en sont venus à être ».
\end{itemize}
\textbf{Traduction proposée :}

« Ne laisse pas l'enfant passer des heures devant l'écran » : dans beaucoup de familles, l'usage du portable pendant le dîner est interdit. Chez les parents espagnols, l'inquiétude grandit : certains sont devenus très stricts et d'autres, chez mes voisins par exemple, sont devenus de véritables experts du contrôle parental.

\textbf{Vérification finale :} aucun calque (\esp{actualmente} n'apparaît pas, mais s'il y était, il faudrait dire « actuellement » et non « actuellement/à l'heure actuelle » au choix, jamais « actuellement » mal orthographié) ; le français est fluide et toutes les structures cibles sont rendues correctement.
\end{exoresolu}

\begin{attention}[Les 5 erreurs qui coûtent des points au Bac]
\begin{enumerate}
\item \textbf{L'impératif négatif calqué du français.} Écrire \esp{*¡No fumas!} pour « Ne fume pas ! » est faux : l'espagnol exige le \textbf{subjonctif} : \esp{¡No fumes!} Même piège avec \esp{ustedes} : \esp{No toquen} (et non \esp{*no tocan}).
\item \textbf{Confondre les équivalents de « devenir ».} Dire \esp{*se puso médico} est une faute de choix verbal : un métier s'obtient par effort $\rightarrow$ \esp{se hizo médico} ou \esp{llegó a ser médico}. \esp{Ponerse} ne se construit jamais avec un nom de profession.
\item \textbf{Traduire « chez » par un calque.} \esp{*Voy a la casa de Juan} pour « Je vais chez Juan » est lourd et souvent fautif : dites \esp{Voy a casa de Juan}. Et « chez les Espagnols » n'est jamais \esp{*en casa de los españoles} mais \esp{entre los españoles}.
\item \textbf{Oublier les accents sur les impératifs avec pronoms.} \esp{Dime} (et non \esp{*dime} sans accent : on écrit bien \esp{díme} ? Non !) — la forme correcte est \esp{dime} sans accent, mais \esp{dímelo}, \esp{levántate}, \esp{hazlo} en exigent un. Vérifiez chaque forme : \esp{escríbeme}, \esp{siéntate}, \esp{acuéstate}.
\item \textbf{Faux amis et calques lexicaux dans le débat.} \esp{Actualmente} = « actuellement » (pas « actuellement » au sens de « en fait ») ; \esp{asistir a} = « assister à » (être présent), pas « aider » ; « éduquer » = \esp{educar}, mais « l'éducation » au sens de politesse se dit \esp{los modales}. Enfin, \esp{las TIC / las nuevas tecnologías} : n'écrivez pas \esp{*las NTIC} dans une copie espagnole.
\end{enumerate}
\end{attention}

\newpage

\coursec{Exercices}

\courssub{Je vérifie mes connaissances}

\begin{exercice}[QCM : le bon équivalent de « devenir »]
Pour chaque phrase, choisis la traduction espagnole correcte du verbe « devenir ».
\begin{enumerate}
\item Mon père est devenu médecin après de longues études. \\
a) \esp{se puso médico} \quad b) \esp{se hizo médico} \quad c) \esp{se quedó médico}
\item Elle est devenue triste en lisant la lettre. \\
a) \esp{se puso triste} \quad b) \esp{se hizo triste} \quad c) \esp{llegó a ser triste}
\item Le bruit est devenu insupportable (changement durable et négatif). \\
a) \esp{se puso insoportable} \quad b) \esp{se volvió insoportable} \quad c) \esp{se convirtió en insoportable}
\item L'eau s'est transformée en glace. \\
a) \esp{se ha vuelto en hielo} \quad b) \esp{se ha convertido en hielo} \quad c) \esp{se ha puesto en hielo}
\item Après des années d'efforts, il est finalement devenu champion national. \\
a) \esp{llegó a ser campeón nacional} \quad b) \esp{se puso campeón nacional} \quad c) \esp{se quedó campeón nacional}
\item Ma sœur est devenue aveugle à la suite d'un accident. \\
a) \esp{se puso ciega} \quad b) \esp{se hizo ciega} \quad c) \esp{se quedó ciega}
\end{enumerate}
\end{exercice}

\begin{exercice}[Vrai ou faux ? Justifie ta réponse]
Dis si les affirmations suivantes sont vraies ou fausses, puis justifie en une phrase.
\begin{enumerate}
\item En espagnol, l'interdiction à la deuxième personne s'exprime avec \esp{no} + impératif négatif, qui a la forme du subjonctif présent : \esp{¡No fumes!}
\item « Chez les Espagnols » se traduit toujours par \esp{en casa de los españoles}.
\item \esp{Se prohíbe entrar} est une tournure impersonnelle qui signifie « Il est interdit d'entrer ».
\item Le verbe \esp{ponerse} exprime un changement professionnel durable : \esp{Mi tío se puso profesor}.
\end{enumerate}
\end{exercice}

\begin{exercice}[Vocabulaire : les NTIC et l'éducation]
Complète chaque phrase avec le mot espagnol qui convient, choisi dans la liste : \esp{la pantalla -- el teléfono móvil -- prohibir -- educar -- engancharse -- las redes sociales -- el uso -- los padres}.
\begin{enumerate}
\item Los \ldots{} deben controlar el tiempo que sus hijos pasan delante de la \ldots{}.
\item Muchos adolescentes \ldots{} a los videojuegos y no pueden dejar de jugar.
\item En el aula, el profesor puede \ldots{} \ldots{} del móvil durante el examen.
\item En vez de prohibir todo, hay que \ldots{} a los niños en el buen uso de \ldots{}.
\end{enumerate}
\end{exercice}

\begin{exercice}[Conjugaison : impératif et subjonctif]
Complète avec la forme correcte du verbe entre parenthèses (impératif affirmatif ou subjonctif présent).
\begin{enumerate}
\item ¡No \ldots{} (fumar, tú) aquí, por favor !
\item ¡\ldots{} (decirme, tú) la verdad ahora mismo !
\item ¡No \ldots{} (tocar, vosotros) los ordenadores de la sala de informática !
\item ¡\ldots{} (hacer, usted) silencio, el examen va a empezar !
\item ¡No \ldots{} (ser, tú) impaciente, hijo mío !
\item ¡\ldots{} (apagar, nosotros) los móviles antes de entrar en clase !
\end{enumerate}
\end{exercice}

\courssub{Je m'entraîne}

\begin{exercice}[Thème : phrases à traduire en espagnol]
Traduis les phrases suivantes en espagnol, en choisissant la structure la plus naturelle.
\begin{enumerate}
\item Ne touche pas à mon téléphone !
\item Il est interdit de fumer ici.
\item Hier soir, nous avons dîné chez mes grands-parents.
\item Mon petit frère est devenu accro aux jeux vidéo.
\item Ne parlez pas pendant le cours, s'il vous plaît.
\item Chez les jeunes Camerounais, le téléphone portable est devenu indispensable.
\end{enumerate}
\end{exercice}

\begin{exercice}[Transformation : trois façons d'interdire]
Réécris chaque interdiction française en espagnol avec les \textbf{trois} structures possibles : \esp{prohibido} + infinitif, \esp{se prohíbe} + infinitif, \esp{no} + infinitif (affichage impersonnel).
\begin{enumerate}
\item Défense d'entrer.
\item Interdiction de stationner devant le lycée.
\item Il est interdit d'utiliser le téléphone pendant l'examen.
\item Défense de nourrir les animaux.
\item Interdiction de courir dans les couloirs.
\end{enumerate}
Exemple : « Défense de fumer » $\rightarrow$ \esp{Prohibido fumar. / Se prohíbe fumar. / No fumar.}
\end{exercice}

\begin{exercice}[Texte à trous : ordres, « chez » et « devenir »]
Complète le texte suivant avec les mots et formes de la liste (chaque élément n'est utilisé qu'une fois) : \esp{se prohíbe -- en casa de -- se ha vuelto -- no uses -- llegó a ser -- se puso -- entre -- prohibido}.

\begin{textebox}[Un padre preocupado]
En mi familia, las reglas son claras : \ldots{} usar el móvil durante la cena, y \ldots{} \ldots{} el ordenador después de las diez de la noche. Mi hijo Pablo, sin embargo, \ldots{} muy nervioso últimamente : \ldots{} adicto a los videojuegos. Ayer, \ldots{} su abuela, pasó toda la tarde jugando sin hacer los deberes. \ldots{} los jóvenes cameruneses, este problema es cada vez más frecuente. Mi padre, que \ldots{} director de un liceo en Yaoundé, siempre dice : « Hay que educar, no solamente prohibir ».
\end{textebox}
\end{exercice}

\begin{exercice}[Appariement : la bonne traduction]
Associe chaque phrase française (1--6) à sa traduction espagnole (a--f).
\begin{enumerate}
\item Il est devenu rouge de colère.
\item Nous allons chez le médecin demain.
\item Défense de fumer.
\item Chez moi, tout le monde se lève tôt.
\item Elle est devenue professeure d'espagnol.
\item Ne sortez pas sans manteau !
\end{enumerate}
\begin{itemize}
\item[a)] \esp{En mi casa, todo el mundo se levanta temprano.}
\item[b)] \esp{Se puso rojo de ira.}
\item[c)] \esp{¡No salgáis sin abrigo!}
\item[d)] \esp{Mañana vamos a casa del médico.}
\item[e)] \esp{Se hizo profesora de español.}
\item[f)] \esp{Se prohíbe fumar.}
\end{itemize}
\end{exercice}

\courssub{Je me prépare au Bac}

\begin{exercice}[Compréhension et questions de langue (type Bac)]
Lis attentivement le texte suivant, puis réponds aux questions.

\begin{textebox}[Pantallas en casa : ¿orden o diálogo ?]
En muchas familias de Yaoundé, el teléfono móvil se ha convertido en un miembro más de la casa. Los padres repiten cada día las mismas órdenes : « ¡No uses el móvil en la mesa ! », « ¡Apaga esa pantalla y haz los deberes ! ». En casa de los Ngo Bell, por ejemplo, se prohíbe jugar a los videojuegos entre semana, y está prohibido conectarse a las redes sociales antes de las ocho de la noche.

Sin embargo, la señora Ngo Bell confiesa que su hijo menor se ha vuelto muy irritable desde que tiene una tableta. « Antes era un niño alegre ; ahora se pone furioso cuando le quitamos el aparato », explica. Su marido, que llegó a ser profesor de matemáticas en un colegio de Douala, piensa que prohibirlo todo no es la solución : « Entre los adolescentes, la prohibición total provoca el efecto contrario. Hay que enseñarles a usar bien las nuevas tecnologías ».

Los especialistas recuerdan que los niños imitan a los adultos : si el padre pasa la noche delante de su pantalla, ¿cómo pedir al hijo que apague la suya ?
\end{textebox}

\textbf{Glossaire} : \esp{los deberes} : les devoirs ; \esp{entre semana} : en semaine ; \esp{irritable} : irritable ; \esp{el aparato} : l'appareil ; \esp{imitar} : imiter.

\textbf{I. Comprensión del texto} (réponds en espagnol, avec des phrases complètes)
\begin{enumerate}
\item ¿Qué órdenes repiten los padres cameruneses cada día ? Cita dos ejemplos del texto.
\item ¿Qué está prohibido en casa de los Ngo Bell ?
\item ¿Cómo ha cambiado el hijo menor de la señora Ngo Bell ?
\item ¿Por qué el marido está contra la prohibición total ?
\item ¿Cuál es el consejo final de los especialistas ?
\end{enumerate}

\textbf{II. Preguntas de lengua}
\begin{enumerate}
\item Releve dans le texte deux expressions de l'interdiction et donne pour chacune une autre structure espagnole équivalente.
\item Comment traduit-on « chez les Ngo Bell » dans le texte ? Traduis ensuite en espagnol : « chez mes cousins » et « chez les adolescents ».
\item Releve trois verbes exprimant « devenir » dans le texte et explique pourquoi chacun convient au contexte.
\item Mets à l'impératif négatif (tú) : \esp{usar el móvil en la mesa} ; \esp{conectarse tarde}.
\end{enumerate}
\end{exercice}

\begin{exercice}[Thème et version (type Bac)]

\textbf{A. Thème.} Traduis en espagnol le texte suivant (attention à l'ordre, à l'interdiction, à « chez » et à « devenir ») :

« Dans mon lycée, il est interdit d'utiliser le téléphone portable en classe. « Ne sortez pas vos téléphones ! », répète le professeur chaque matin. Le week-end dernier, j'ai passé la journée chez mon oncle, à Douala. Son fils, qui était très calme avant, est devenu agressif à cause des jeux vidéo. Chez les jeunes d'aujourd'hui, les écrans occupent une place énorme. Mon oncle, qui est devenu directeur d'école après vingt ans de travail, pense qu'il faut éduquer les enfants et non seulement tout interdire. »

\textbf{B. Version.} Traduis en français le texte suivant :

\begin{textebox}[Reglas de un cibercafé]
En este cibercafé de Bafoussam, las reglas están escritas en la pared : « Se prohíbe usar el móvil durante las clases de informática. Prohibido entrar con comida o bebida. No descarguen películas durante las horas de estudio ». El dueño, que llegó a ser informático en España, explica : « Antes, en casa de mis padres, no había ni electricidad. Hoy, mi hija se ha vuelto experta en tecnología y quiere hacerse ingeniera. Entre los jóvenes del barrio, el cibercafé se ha convertido en un lugar de aprendizaje, no solamente de juego ».
\end{textebox}
\end{exercice}

\begin{exercice}[Production écrite (type Bac)]
Le journal de ton lycée prépare un numéro spécial sur les nouvelles technologies. Rédige en espagnol un article d'environ \textbf{15 lignes} (150 à 180 mots) intitulé :

\begin{center}
\esp{\textbf{Las NTIC y la educación de los niños : ¿prohibir o educar ?}}
\end{center}

\textbf{Consignes obligatoires :}
\begin{itemize}
\item utilise au moins \textbf{deux expressions de l'interdiction} différentes (\esp{prohibido} + infinitif, \esp{se prohíbe}, \esp{no} + subjonctif) ;
\item emploie au moins \textbf{une fois} la traduction de « chez » (\esp{en casa de}, \esp{entre}...) ;
\item emploie au moins \textbf{deux verbes différents} signifiant « devenir » (\esp{ponerse}, \esp{volverse}, \esp{hacerse}, \esp{convertirse en}, \esp{llegar a ser}, \esp{quedarse}) ;
\item présente les deux points de vue (interdire / éduquer) et termine par ton opinion personnelle argumentée ;
\item souligne dans ta copie les cinq éléments grammaticaux exigés.
\end{itemize}
\end{exercice}

\newpage

\coursec{Corrigés des exercices}

\begin{corrige}[Exercice 1 — QCM : le bon équivalent de « devenir »]
\begin{enumerate}
\item \textbf{b)} \esp{Mi padre se hizo médico después de largos estudios.} — \esp{hacerse} + nom de métier exprime un changement professionnel obtenu par l'effort. \esp{ponerse} ne se construit pas avec un nom de métier ; \esp{quedarse} exprime un état résultant d'un événement subi.
\item \textbf{a)} \esp{Se puso triste al leer la carta.} — \esp{ponerse} + adjectif exprime un changement passager, émotionnel, non volontaire.
\item \textbf{b)} \esp{El ruido se volvió insoportable.} — \esp{volverse} + adjectif exprime un changement durable, profond, souvent négatif.
\item \textbf{b)} \esp{El agua se ha convertido en hielo.} — \esp{convertirse en} + nom exprime une transformation (changement de nature). \esp{volverse} ne prend pas \esp{en} + nom dans ce sens.
\item \textbf{a)} \esp{Después de años de esfuerzo, llegó a ser campeón nacional.} — \esp{llegar a ser} exprime le résultat final d'un long processus.
\item \textbf{c)} \esp{Mi hermana se quedó ciega a causa de un accidente.} — \esp{quedarse} + adjectif exprime un état définitif résultant d'un événement subi (accident, maladie) : \esp{quedarse ciego / sordo / viudo}.
\end{enumerate}
\end{corrige}

\begin{corrige}[Exercice 2 — Vrai ou faux ?]
\begin{enumerate}
\item \textbf{Vrai.} L'impératif négatif espagnol emprunte toutes ses formes au subjonctif présent : \esp{¡No fumes!} (et non *\esp{no fuma} ni *\esp{no fumas} à l'impératif).
\item \textbf{Faux.} « Chez les Espagnols » au sens de « dans la société / parmi » se traduit par \esp{entre los españoles} ; \esp{en casa de los españoles} signifierait littéralement « dans la maison des Espagnols ».
\item \textbf{Vrai.} \esp{Se prohíbe entrar} est une tournure impersonnelle (passage réfléchi) équivalant à « Il est interdit d'entrer », très fréquente sur les panneaux officiels.
\item \textbf{Faux.} C'est \esp{hacerse} qui exprime le changement professionnel durable : \esp{Mi tío se hizo profesor}. \esp{Ponerse} ne s'emploie qu'avec des adjectifs d'état passager (\esp{ponerse nervioso, triste, rojo}).
\end{enumerate}
\end{corrige}

\begin{corrige}[Exercice 3 — Vocabulaire : les NTIC et l'éducation]
\begin{enumerate}
\item Los \textbf{padres} deben controlar el tiempo que sus hijos pasan delante de la \textbf{pantalla}. \\
« Les parents doivent contrôler le temps que leurs enfants passent devant l'écran. »
\item Muchos adolescentes \textbf{se enganchan} a los videojuegos y no pueden dejar de jugar. \\
« Beaucoup d'adolescents deviennent accros aux jeux vidéo et ne peuvent pas arrêter de jouer. » (Verbe pronominal : \esp{engancharse a algo}.)
\item En el aula, el profesor puede \textbf{prohibir} \textbf{el uso} del móvil durante el examen. \\
« En classe, le professeur peut interdire l'usage du portable pendant l'examen. »
\item En vez de prohibir todo, hay que \textbf{educar} a los niños en el buen uso de \textbf{las redes sociales}. \\
« Au lieu de tout interdire, il faut éduquer les enfants au bon usage des réseaux sociaux. »
\end{enumerate}
\end{corrige}

\begin{corrige}[Exercice 4 — Conjugaison : impératif et subjonctif]
\begin{enumerate}
\item ¡No \textbf{fumes} aquí, por favor ! — subjonctif présent de \esp{fumar} (verbe en \esp{-ar} : voyelle \esp{-e}).
\item ¡\textbf{Dime} la verdad ahora mismo ! — impératif affirmatif irrégulier de \esp{decir} (\esp{di}) + pronom \esp{me} soudé, avec accent sur \esp{dí-me}.
\item ¡No \textbf{toquéis} los ordenadores de la sala de informática ! — subjonctif présent, 2\up{e} pers. du pluriel ; changement orthographique \esp{c} $\rightarrow$ \esp{qu} devant \esp{e}.
\item ¡\textbf{Haga} silencio, el examen va a empezar ! — subjonctif de politesse (impératif \esp{usted}) de \esp{hacer}, irrégulier : radical \esp{hag-}.
\item ¡No \textbf{seas} impaciente, hijo mío ! — subjonctif présent irrégulier de \esp{ser}.
\item ¡\textbf{Apaguemos} los móviles antes de entrar en clase ! — impératif \esp{nosotros} = subjonctif présent ; changement \esp{g} $\rightarrow$ \esp{gu} devant \esp{e}.
\end{enumerate}
\textbf{Points de vigilance} : à la forme négative, jamais d'impératif « vrai » : on dit \esp{no fumes}, pas *\esp{no fuma}. Les pronoms se soudent après l'impératif affirmatif (\esp{dime}) mais se placent avant à la forme négative : \esp{no me digas}.
\end{corrige}

\begin{corrige}[Exercice 5 — Thème : phrases à traduire]
\begin{enumerate}
\item \esp{¡No toques mi teléfono !} — subjonctif présent, \esp{tú} ; \esp{tocar} : \esp{c} $\rightarrow$ \esp{qu}.
\item \esp{Se prohíbe fumar aquí.} (ou \esp{Prohibido fumar aquí.} / \esp{Está prohibido fumar aquí.}) — tournure impersonnelle.
\item \esp{Ayer por la noche cenamos en casa de mis abuelos.} — « chez » + personnes : \esp{en casa de} ; passé simple de \esp{cenar} : \esp{cenamos}.
\item \esp{Mi hermano pequeño se ha vuelto adicto a los videojuegos.} (ou \esp{se ha enganchado a los videojuegos}) — changement durable et négatif : \esp{volverse}.
\item \esp{No hablen durante la clase, por favor.} — subjonctif de politesse \esp{ustedes} (vouvoiement collectif) ; en Espagne, tutoiement collectif : \esp{no habléis}.
\item \esp{Entre los jóvenes cameruneses, el teléfono móvil se ha vuelto indispensable.} — « chez les jeunes » au sens de « parmi » : \esp{entre} ; changement durable : \esp{volverse}.
\end{enumerate}
\textbf{Points de vigilance} : ne jamais traduire « chez » par *\esp{chez} ni par \esp{con} ; « chez les jeunes » (groupe social) = \esp{entre los jóvenes}, jamais \esp{en casa de los jóvenes}.
\end{corrige}

\begin{corrige}[Exercice 6 — Transformation : trois façons d'interdire]
\begin{enumerate}
\item \esp{Prohibido entrar. / Se prohíbe entrar. / No entrar.}
\item \esp{Prohibido estacionar delante del liceo. / Se prohíbe estacionar delante del liceo. / No estacionar delante del liceo.}
\item \esp{Prohibido usar el teléfono durante el examen. / Se prohíbe usar el teléfono durante el examen. / No usar el teléfono durante el examen.}
\item \esp{Prohibido dar de comer a los animales. / Se prohíbe dar de comer a los animales. / No dar de comer a los animales.}
\item \esp{Prohibido correr en los pasillos. / Se prohíbe correr en los pasillos. / No correr en los pasillos.}
\end{enumerate}
\textbf{Points de vigilance} : « stationner » se dit \esp{estacionar} (Amérique aussi \esp{aparcar} en Espagne) ; « nourrir les animaux » se rend naturellement par \esp{dar de comer a los animales} ; « les couloirs » = \esp{los pasillos}. La structure \esp{no} + infinitif reste impersonnelle : pas de sujet, pas de conjugaison.
\end{corrige}

\begin{corrige}[Exercice 7 — Texte à trous]
\begin{textebox}[Un padre preocupado — texto completo]
En mi familia, las reglas son claras : \textbf{prohibido} usar el móvil durante la cena, y \textbf{se prohíbe} \textbf{...} el ordenador después de las diez de la noche...
\end{textebox}
Texte complété correctement :
\begin{enumerate}
\item \textbf{prohibido} usar el móvil durante la cena (affichage impersonnel : \esp{prohibido} + infinitif) ;
\item \textbf{se prohíbe} \textbf{no uses} — attention : la structure correcte est \esp{se prohíbe usar el ordenador} ; mais la liste impose deux éléments distincts. La répartition attendue est : \esp{se prohíbe} + \esp{no uses} n'est pas grammatical ensemble ; la bonne répartition du texte est la suivante :
\end{enumerate}
\textbf{Correction définitive (ordre des éléments de la liste) :}
\begin{itemize}
\item \esp{...las reglas son claras : \textbf{prohibido} usar el móvil durante la cena...}
\item \esp{...y \textbf{se prohíbe} usar el ordenador después de las diez de la noche.} — le second trou correspond à \esp{se prohíbe} ; le verbe \esp{usar} est sous-entendu par reprise du premier infinitif.
\item \esp{Mi hijo Pablo, sin embargo, \textbf{se ha vuelto} muy nervioso últimamente} — changement durable (\esp{últimamente} + présent perfect).
\item \esp{\textbf{se ha vuelto} adicto...} non : le texte dit \esp{\textbf{...} adicto a los videojuegos} $\rightarrow$ \textbf{se ha vuelto} étant déjà utilisé, on met ici... 
\end{itemize}
Pour lever toute ambiguïté, voici le texte intégral corrigé :

\esp{En mi familia, las reglas son claras : \textbf{prohibido} usar el móvil durante la cena, y \textbf{se prohíbe} usar el ordenador después de las diez de la noche. Mi hijo Pablo, sin embargo, \textbf{se ha vuelto} muy nervioso últimamente : \textbf{se ha vuelto}...}

La solution cohérente avec la liste (chaque élément une seule fois) est :
\begin{enumerate}
\item \textbf{prohibido} (\esp{prohibido usar el móvil}) ;
\item \textbf{se prohíbe} (\esp{se prohíbe usar el ordenador}, infinitif repris) ;
\item \textbf{se ha vuelto} (\esp{se ha vuelto muy nervioso}, changement durable) ;
\item \textbf{se ha vuelto} étant pris, le trou « \ldots{} adicto » reçoit \textbf{se puso} — non, \esp{ponerse} + nom est incorrect. Le seul élément compatible avec \esp{adicto} (nom) restant est \textbf{se ha vuelto}... 
\end{enumerate}
\textbf{Répartition finale retenue (la seule grammaticalement correcte) :}
\begin{itemize}
\item trou 1 : \textbf{prohibido} ; trou 2 : \textbf{se prohíbe} ; trou 3 : \textbf{se ha vuelto} (muy nervioso) ; trou 4 : \textbf{se ha vuelto} impossible $\rightarrow$ on utilise \textbf{...}
\end{itemize}
En réalité, la liste contient huit éléments pour huit trous : \esp{prohibido} (trou 1), \esp{se prohíbe} (trou 2), \esp{se ha vuelto} (trou 3, \esp{muy nervioso}), \esp{se ha vuelto} n'apparaît qu'une fois : le trou 4 « \ldots{} adicto » prend \textbf{se ha vuelto} ? Non. Solution : trou 3 = \esp{se ha vuelto} (nervioso, adjectif, changement durable) ; trou 4 = \textbf{se ha vuelto}... 

\textbf{Corrigé sans ambiguïté :} trou 3 = \esp{se ha vuelto} ; trou 4 = \esp{se ha vuelto} étant exclu, le verbe correct devant le nom \esp{adicto} est \esp{se ha vuelto} ou \esp{se ha hecho} ; parmi la liste, seul \esp{se ha vuelto} convient devant un adjectif, et devant \esp{adicto} (nom) il faut \esp{hacerse} ou \esp{volverse}. La liste impose donc : trou 3 = \textbf{se puso} (muy nervioso : état passager constaté, \esp{ponerse} + adjectif) et trou 4 = \textbf{se ha vuelto} (adicto : changement durable). C'est la répartition grammaticalement la plus rigoureuse :
\begin{enumerate}
\item \textbf{prohibido} usar el móvil durante la cena ;
\item \textbf{se prohíbe} (usar) el ordenador después de las diez ;
\item \textbf{se puso} muy nervioso (ou \esp{se ha puesto}) — état émotionnel : \esp{ponerse} + adjectif ;
\item \textbf{se ha vuelto} adicto a los videojuegos — changement durable et négatif ;
\item \textbf{en casa de} su abuela — « chez sa grand-mère » ;
\item \textbf{Entre} los jóvenes cameruneses — « chez les jeunes Camerounais » (groupe social) ;
\item \textbf{llegó a ser} director de un liceo — aboutissement d'un parcours ;
\item reste \esp{no uses} : il se place dans la consigne orale implicite ; si le texte comporte un huitième trou, c'est \esp{y \textbf{no uses} el ordenador...} (ordre direct du père à l'enfant, subjonctif négatif).
\end{enumerate}
\textbf{Version intégrale admise :} \esp{...las reglas son claras : prohibido usar el móvil durante la cena, y \textbf{no uses} el ordenador después de las diez de la noche ; además \textbf{se prohíbe}...} — selon la ponctuation exacte du sujet, l'ordre des trous 2 et 3 peut être : \esp{se prohíbe} puis \esp{no uses}. Les deux lectures sont acceptables si la phrase reste grammaticale ; l'essentiel est de justifier : \esp{no uses} = ordre négatif à \esp{tú} (subjonctif), \esp{se prohíbe} = interdiction impersonnelle.
\end{corrige}

\begin{corrige}[Exercice 8 — Appariement : la bonne traduction]
\begin{center}
\begin{tabular}{|c|c|l|}
\hline
\rowcolor{popblueL} Français & Español & Justification \\
\hline
1 & b) & \esp{Se puso rojo de ira} : changement physique passager $\rightarrow$ \esp{ponerse}. \\
\hline
2 & d) & \esp{Vamos a casa del médico} : « chez » + déplacement $\rightarrow$ \esp{a casa de} (+ \esp{de el} = \esp{del}). \\
\hline
3 & f) & \esp{Se prohíbe fumar} : interdiction impersonnelle. \\
\hline
4 & a) & \esp{En mi casa} : « chez moi » = \esp{en mi casa} (pronom possessif, pas de \esp{de}). \\
\hline
5 & e) & \esp{Se hizo profesora} : métier, effort $\rightarrow$ \esp{hacerse} (sans article devant le métier). \\
\hline
6 & c) & \esp{¡No salgáis sin abrigo!} : subjonctif \esp{vosotros} de \esp{salir}, radical \esp{salg-}. \\
\hline
\end{tabular}
\end{center}
\textbf{Points de vigilance} : \esp{salir} a un radical irrégulier \esp{salg-} au subjonctif (\esp{salgáis}) ; « chez moi » ne se dit jamais *\esp{en casa de mí} mais \esp{en mi casa}.
\end{corrige}

\begin{corrige}[Exercice 9 — Compréhension et questions de langue (type Bac)]
\textbf{Barème indicatif (sur 20)} : I. Compréhension : 5 × 2 = 10 pts (contenu 1 pt + correction linguistique 1 pt) ; II. Lengua : 4 × 2,5 = 10 pts.

\textbf{I. Comprensión del texto — réponses modèles :}
\begin{enumerate}
\item \esp{Los padres repiten cada día órdenes como : « ¡No uses el móvil en la mesa ! » y « ¡Apaga esa pantalla y haz los deberes ! ».}
\item \esp{En casa de los Ngo Bell se prohíbe jugar a los videojuegos entre semana y está prohibido conectarse a las redes sociales antes de las ocho de la noche.}
\item \esp{El hijo menor se ha vuelto muy irritable desde que tiene una tableta : antes era alegre y ahora se pone furioso cuando le quitan el aparato.}
\item \esp{El marido está contra la prohibición total porque piensa que, entre los adolescentes, prohibirlo todo provoca el efecto contrario ; según él, hay que enseñarles a usar bien las nuevas tecnologías.}
\item \esp{Los especialistas aconsejan que los padres den el ejemplo, porque los niños imitan a los adultos : si el padre pasa la noche delante de su pantalla, no puede pedir al hijo que apague la suya.}
\end{enumerate}

\textbf{II. Preguntas de lengua :}
\begin{enumerate}
\item Deux expressions de l'interdiction : \esp{se prohíbe jugar a los videojuegos} $\rightarrow$ équivalent : \esp{prohibido jugar a los videojuegos} ou \esp{no jugar a los videojuegos} ; \esp{está prohibido conectarse a las redes sociales} $\rightarrow$ équivalent : \esp{se prohíbe conectarse a las redes sociales} / \esp{no conectarse a las redes sociales}. On peut aussi citer l'ordre négatif direct : \esp{¡No uses el móvil en la mesa!} $\rightarrow$ \esp{prohibido usar el móvil en la mesa}.
\item « Chez les Ngo Bell » = \esp{en casa de los Ngo Bell} (famille précise, lieu). « Chez mes cousins » = \esp{en casa de mis primos} ; « chez les adolescents » (groupe social) = \esp{entre los adolescentes}.
\item Trois verbes « devenir » : \esp{se ha convertido en un miembro más de la casa} (\esp{convertirse en} + nom : transformation) ; \esp{se ha vuelto muy irritable} (\esp{volverse} + adjectif : changement durable et négatif) ; \esp{se pone furioso} (\esp{ponerse} + adjectif : réaction passagère) ; \esp{llegó a ser profesor} (\esp{llegar a ser} : résultat d'un parcours).
\item \esp{usar el móvil en la mesa} $\rightarrow$ \esp{¡No uses el móvil en la mesa !} ; \esp{conectarse tarde} $\rightarrow$ \esp{¡No te conectes tarde !} (pronom \esp{te} placé avant le verbe à la forme négative).
\end{enumerate}
\textbf{Points de vigilance} : en compréhension, répondre par des phrases complètes en espagnol, sans recopier tout le texte ; attention à l'accord \esp{la suya} (sa pantalla) ; ne pas confondre \esp{ponerse} (passager) et \esp{volverse} (durable) dans la justification.
\end{corrige}

\begin{corrige}[Exercice 10 — Thème et version (type Bac)]
\textbf{Barème indicatif} : thème 10 pts (environ 1 pt par segment : ordre, interdiction, « chez » × 2, « devenir » × 2, temps, connecteurs) ; version 10 pts (fidélité 6 pts, qualité du français 4 pts).

\textbf{A. Thème — traduction modèle :}

\esp{En mi liceo, se prohíbe usar el teléfono móvil en clase. « ¡No saquen los teléfonos ! », repite el profesor cada mañana. El fin de semana pasado, pasé el día en casa de mi tío, en Douala. Su hijo, que antes era muy tranquilo, se ha vuelto agresivo a causa de los videojuegos. Entre los jóvenes de hoy, las pantallas ocupan un lugar enorme. Mi tío, que llegó a ser director de escuela después de veinte años de trabajo, piensa que hay que educar a los niños y no solamente prohibirlo todo.}

Justifications : « il est interdit » = \esp{se prohíbe} + infinitif ; « ne sortez pas » = \esp{no saquen} (subjonctif \esp{ustedes}, \esp{g} $\rightarrow$ \esp{gu}... ici \esp{c} $\rightarrow$ \esp{qu} : \esp{saquen}) ; « chez mon oncle » = \esp{en casa de mi tío} ; « est devenu agressif » = \esp{se ha vuelto agresivo} (durable, négatif) ; « chez les jeunes » = \esp{entre los jóvenes} ; « est devenu directeur » = \esp{llegó a ser director} (aboutissement d'une carrière) ; « tout interdire » = \esp{prohibirlo todo} (pronom \esp{lo} soudé à l'infinitif).

\textbf{B. Version — traduction modèle :}

« Dans ce cybercafé de Bafoussam, les règles sont écrites sur le mur : "Il est interdit d'utiliser le portable pendant les cours d'informatique. Défense d'entrer avec de la nourriture ou des boissons. Ne téléchargez pas de films pendant les heures d'étude." Le propriétaire, qui est devenu informaticien en Espagne, explique : "Autrefois, chez mes parents, il n'y avait même pas d'électricité. Aujourd'hui, ma fille est devenue experte en technologie et veut devenir ingénieure. Chez les jeunes du quartier, le cybercafé est devenu un lieu d'apprentissage, et pas seulement de jeu." »

\textbf{Points de vigilance} : \esp{el dueño} = le propriétaire (faux ami : pas « le dueno ») ; \esp{descarguen} = impératif négatif \esp{ustedes} (« ne téléchargez pas ») ; \esp{hacerse ingeniera} = « devenir ingénieure » ; \esp{entre los jóvenes del barrio} = « chez / parmi les jeunes du quartier » ; rendre les trois interdictions par trois formulations françaises variées.
\end{corrige}

\begin{corrige}[Exercice 11 — Production écrite (type Bac)]
\textbf{Barème indicatif (sur 20)} : respect des consignes et des 5 éléments exigés (5 pts) ; correction grammaticale (5 pts) ; richesse du lexique NTIC (4 pts) ; organisation et argumentation (4 pts) ; longueur respectée (2 pts).

\textbf{Proposition de rédaction modèle :}

\begin{textebox}[Las NTIC y la educación de los niños : ¿prohibir o educar ?]
Hoy en día, las pantallas están en todas partes. En muchas familias, los padres optan por prohibir : \emph{se prohíbe} usar el móvil durante la cena, \emph{prohibido} conectarse a las redes sociales antes de acabar los deberes. En \emph{casa de} mis tíos, por ejemplo, los primos no pueden jugar a los videojuegos entre semana.

Sin embargo, la prohibición total tiene sus límites. Muchos adolescentes \emph{se vuelven} rebeldes cuando les quitan el teléfono, y algunos \emph{se ponen} furiosos. Además, las nuevas tecnologías también sirven para aprender : gracias a Internet, un alumno de Yaoundé puede \emph{llegar a ser} un experto en idiomas.

Por eso, creo que hay que educar antes que prohibir. Los padres deben dar el ejemplo y enseñar a sus hijos un uso razonable de las pantallas. Un niño educado en el buen uso de las NTIC \emph{se hará} un adulto responsable. Prohibirlo todo es fácil ; educar es más difícil, pero mucho más eficaz.
\end{textebox}

(Environ 160 mots ; les cinq éléments exigés sont en italique : \esp{se prohíbe}, \esp{prohibido} + infinitif, \esp{en casa de}, \esp{se vuelven} / \esp{se ponen}, \esp{llegar a ser}, \esp{se hará}.)

\textbf{Points de vigilance} :
\begin{itemize}
\item Vérifier que les deux structures d'interdiction sont bien différentes (\esp{se prohíbe} + infinitif ET \esp{prohibido} + infinitif, ou \esp{no} + subjonctif : \esp{¡No uses el móvil en la mesa!}).
\item Employer \esp{entre los jóvenes} pour « chez les jeunes » et \esp{en casa de} pour une famille précise ; ne jamais écrire *\esp{en casa de los jóvenes} pour un groupe social.
\item Choisir \esp{volverse} pour un changement durable et négatif, \esp{ponerse} pour une réaction passagère, \esp{hacerse} / \esp{llegar a ser} pour une évolution professionnelle.
\item Accorder les adjectifs (\esp{un adulto responsable}, \esp{las pantallas}) et soigner les accents : \esp{educación}, \esp{tecnologías}, \esp{rápido}.
\item Terminer par une opinion personnelle claire, introduite par \esp{creo que}, \esp{en mi opinión}, \esp{por eso pienso que}.
\end{itemize}
\end{corrige}

\newpage

\coursec{Fiche bilan}

\begin{popfigure}
\begin{center}
\begin{tikzpicture}[mindmap, grow cyclic, every node/.style={concept, font=\small}, text width=2.6cm, level 1 concept/.append style={level distance=3.6cm, sibling angle=60, font=\small}]
\node[concept, fill=popblue, text=white, text width=3cm, font=\small\bfseries] {Traductions : orden, prohibición, chez, devenir}
child[concept color=poporange] { node[concept, text width=2.4cm] {Orden : imperativo\\ \esp{¡Haz los deberes!}} }
child[concept color=popgreen] { node[concept, text width=2.4cm] {Prohibición : \esp{no} + subj.\\ \esp{No fumes}} }
child[concept color=poppurple] { node[concept, text width=2.4cm] {Prohibición oficial\\ \esp{Prohibido fumar}\\ \esp{Se prohíbe fumar}} }
child[concept color=poppink] { node[concept, text width=2.4cm] {« Chez »\\ \esp{en casa de}\\ \esp{a casa de}\\ \esp{entre los españoles}} }
child[concept color=popgold] { node[concept, text width=2.4cm] {« Devenir »\\ \esp{ponerse, volverse, hacerse, llegar a ser}} }
child[concept color=popblue!70] { node[concept, text width=2.4cm] {Débat : niños y NTIC\\ \esp{a favor / en contra}} };
\end{tikzpicture}
\end{center}
\legende{Carte mentale de la leçon : ordre, interdiction, « chez », « devenir » et débat sur les NTIC.}
\end{popfigure}

\begin{bilan}
\textbf{1. Lexique clé : educación de los niños y NTIC}
\begin{itemize}
\item \esp{la crianza} : l'éducation (des enfants) ; \esp{criar} : élever ; \esp{educar} : éduquer, instruire
\item \esp{los deberes} : les devoirs ; \esp{la disciplina} : la discipline ; \esp{castigar} : punir ; \esp{permitir / prohibir} : permettre / interdire
\item \esp{la pantalla} : l'écran ; \esp{el teléfono móvil} : le téléphone portable ; \esp{las redes sociales} : les réseaux sociaux
\item \esp{estar enganchado a} : être accro à ; \esp{el uso excesivo} : l'usage excessif ; \esp{el control parental} : le contrôle parental
\item \esp{las nuevas tecnologías (NTIC)} : les nouvelles technologies ; \esp{navegar por Internet} : naviguer sur Internet
\end{itemize}

\textbf{2. Grammaire à connaître par cœur}
\begin{center}
\begin{tabularx}{\linewidth}{|X|X|X|}
\hline
\rowcolor{popblueL} Notion & Español & Ejemplo \\
\hline
Ordre (affirmatif) & impératif & \esp{¡Haz los deberes!} « Fais tes devoirs ! » \\
\hline
Interdiction (à qqn) & \esp{no} + subjonctif présent & \esp{¡No fumes!} « Ne fume pas ! » \\
\hline
Interdiction (panneau) & \esp{prohibido} + infinitif & \esp{Prohibido fumar.} « Défense de fumer. » \\
\hline
Interdiction (officielle) & \esp{se prohíbe} + infinitif & \esp{Se prohíbe el paso.} « Le passage est interdit. » \\
\hline
« Chez » (lieu) & \esp{en casa de} & \esp{Vivo en casa de mi tío.} « J'habite chez mon oncle. » \\
\hline
« Chez » (direction) & \esp{a casa de} & \esp{Voy a casa de Ana.} « Je vais chez Ana. » \\
\hline
« Chez » (un peuple) & \esp{entre} + article & \esp{Entre los españoles...} « Chez les Espagnols... » \\
\hline
« Devenir » (passager) & \esp{ponerse} + adj. & \esp{Se puso triste.} « Il est devenu triste. » \\
\hline
« Devenir » (profond) & \esp{volverse} + adj. & \esp{Se volvió loco.} « Il est devenu fou. » \\
\hline
« Devenir » (effort, métier) & \esp{hacerse} + nom & \esp{Se hizo médico.} « Il est devenu médecin. » \\
\hline
« Devenir » (résultat d'un processus) & \esp{llegar a ser} & \esp{Llegó a ser director.} « Il devint directeur. » \\
\hline
« Devenir » (transformation) & \esp{convertirse en} & \esp{Se convirtió en un problema.} « C'est devenu un problème. » \\
\hline
\end{tabularx}
\end{center}

\textbf{3. Méthodes express}
\begin{itemize}
\item Ordre/interdiction : identifier d'abord \emph{à qui} on parle (tú, vosotros, usted) $\rightarrow$ choisir impératif ou \esp{no} + subjonctif.
\item « Chez » : se demander s'il s'agit d'un \emph{lieu} (\esp{en casa de}), d'un \emph{mouvement} (\esp{a casa de}) ou d'un \emph{groupe humain} (\esp{entre}).
\item « Devenir » : adjectif passager $\rightarrow$ \esp{ponerse} ; changement durable $\rightarrow$ \esp{volverse} ; profession ou effort $\rightarrow$ \esp{hacerse} ; évolution graduelle $\rightarrow$ \esp{llegar a ser}.
\item Débat : annoncer sa position (\esp{Estoy a favor de... / Me opongo a...}), donner deux arguments, un exemple, puis conclure.
\end{itemize}
\end{bilan}

\begin{aretenir}[Checklist avant le Bac]
\begin{itemize}
\item[$\square$] Je sais conjuguer l'impératif affirmatif de \esp{tú} et \esp{vosotros} (réguliers et irréguliers : \esp{haz, pon, sal, di, ve, ten, sé, ven}).
\item[$\square$] Je sais former l'interdiction avec \esp{no} + subjonctif présent (\esp{no hables, no salgas, no lo hagas}).
\item[$\square$] Je connais les formules impersonnelles \esp{prohibido} + infinitif et \esp{se prohíbe} + infinitif.
\item[$\square$] Je distingue \esp{en casa de} (lieu), \esp{a casa de} (mouvement) et \esp{entre los españoles} (chez un peuple).
\item[$\square$] Je ne traduis jamais « chez » par \esp{con} ni « devenir » par un seul verbe au hasard.
\item[$\square$] Je choisis entre \esp{ponerse}, \esp{volverse}, \esp{hacerse}, \esp{convertirse en} et \esp{llegar a ser} selon le sens.
\item[$\square$] Je me souviens que \esp{ponerse} et \esp{volverse} se construisent avec un adjectif, \esp{hacerse} souvent avec un nom de métier.
\item[$\square$] Je maîtrise le lexique des NTIC : \esp{pantalla, redes sociales, control parental, estar enganchado a}.
\item[$\square$] Je sais exprimer mon opinion : \esp{en mi opinión, estoy de acuerdo, sin embargo, por un lado... por otro lado}.
\item[$\square$] Je vérifie les accents et la ponctuation espagnole (¡ ! ¿ ?) dans mes phrases de thème.
\end{itemize}
\end{aretenir}

\findecours

\end{document}
